% !TEX program = xelatex
% 经济与社科速成讲义：配合二试综合“经济”“社科”两份预测训练卷
\documentclass[cjk=noto, nomargin]{mojing-notes}
\usepackage{float}

\course[Crash Course]{经济与社科速成讲义}
\term{2026 年 10 月}
\instructor{Claude（配合本套预测训练卷编写）}
\institute{训练资料，非官方文件}

% ---------------------------------------------------------------------------
% 讲义专用的小盒子
% ---------------------------------------------------------------------------
\NewDocumentEnvironment{pitfall}{}
  {\begin{tcolorbox}[mojing/base, unbreakable, frame hidden, colback=mjVermilion!5,
     left=9pt, right=8pt, top=6pt, bottom=6pt,
     borderline west={1.1pt}{0pt}{mjVermilion}]
   {\mojingheadfont\color{mjVermilion}易错提醒}\par\vspace{.25\baselineskip}}
  {\end{tcolorbox}}
\NewDocumentEnvironment{oneline}{}
  {\begin{tcolorbox}[mojing/base, unbreakable, frame hidden, colback=mjPine!6,
     left=9pt, right=8pt, top=6pt, bottom=6pt,
     borderline west={1.1pt}{0pt}{mjPine}]
   {\mojingheadfont\color{mjPine}一句话记住}\quad}
  {\end{tcolorbox}}

% ---------------------------------------------------------------------------
% 作图约定：需求用黛蓝，供给用朱砂；剩余、税收、无谓损失用浅色填充
% ---------------------------------------------------------------------------
\tikzset{
  dem/.style={draw=mjAccent, line width=1pt},
  sup/.style={draw=mjSecond, line width=1pt},
  sup2/.style={draw=mjSecond, line width=1pt, dash pattern=on 4pt off 2pt},
  dem2/.style={draw=mjAccent, line width=1pt, dash pattern=on 4pt off 2pt},
  csfill/.style={fill=mjAccent!16},
  psfill/.style={fill=mjSecond!14},
  taxfill/.style={fill=mjOchre!28},
  dwlfill/.style={fill=mjPine!38},
  guide/.style={draw=mjGray, line width=.4pt, dash pattern=on 2pt off 1.6pt},
  lab/.style={font=\small},
  slab/.style={font=\footnotesize},
  eqpt/.style={circle, fill=mjInk, inner sep=0pt, minimum size=3.6pt},
  shiftarrow/.style={-{Stealth[length=2.4mm,width=1.8mm]}, draw=mjInk!70, line width=.7pt},
  node box/.style={draw=mjAccent!70, rounded corners=2pt, fill=white, align=center,
    inner sep=4pt, font=\small},
  cause/.style={-{Stealth[length=2.4mm,width=1.8mm]}, draw=mjInk!80, line width=.8pt},
}
% 坐标轴：\econaxes{横轴长}{纵轴长}{横轴名}{纵轴名}
\newcommand{\econaxes}[4]{%
  \draw[mj axis] (0,0) -- (#1,0) node[right, lab] {#3};
  \draw[mj axis] (0,0) -- (0,#2) node[above, lab] {#4};}
% 收益表中的一对收益：甲（行）黛蓝，乙（列）朱砂
\newcommand{\pay}[2]{$(\textcolor{mjAccent}{#1},\,\textcolor{mjSecond}{#2})$}
\newcommand{\upay}[2]{$(\textcolor{mjAccent}{\underline{#1}},\,\textcolor{mjSecond}{\underline{#2}})$}
\newcommand{\upayA}[2]{$(\textcolor{mjAccent}{\underline{#1}},\,\textcolor{mjSecond}{#2})$}
\newcommand{\upayB}[2]{$(\textcolor{mjAccent}{#1},\,\textcolor{mjSecond}{\underline{#2}})$}

\tcbset{mojing/tint/.append style={unbreakable}}  % 定义框较短，不跨页
\setlength{\parskip}{.35\baselineskip}
\captionsetup{font=small}

\begin{document}

\maketitle
\tableofcontents

% =============================================================================
% 使用说明
% =============================================================================
\clearpage
\section*{怎么用这份讲义}
\addcontentsline{toc}{lecture}{\protect\makebox[4.6em][l]{\protect\color{mjAccent}导读}怎么用这份讲义}

这份讲义配合本套预测训练卷中的二试综合“经济”“社科”两卷。它不是逐题讲解，而是把做这两卷要用到的\textbf{基本模型和思考方法}从头讲一遍，
再往外多走一步：卷子里没有直接考、但和这些题目同一类的内容，也一并讲到。

\paragraph{需要什么基础}
高中数学就够了：会解一次方程组、会算三角形面积、知道等比数列求和、会算简单概率。经济学和社会科学的名词全部从头解释。

\paragraph{每一讲怎么读}
\begin{itemize}
  \item 先读开头的\textbf{本讲要点}，知道这一讲要解决什么问题；
  \item 正文每讲一个概念都配一张图。\textbf{先看图，再看文字}，大部分结论在图上一眼就能看出来；
  \item \textcolor{mjVermilion}{\textbf{易错提醒}}是考试中最常见的失分点；\textcolor{mjPine}{\textbf{一句话记住}}是这一讲最该带走的结论；
  \item 讲末有 2--4 道“练一练”，答案在第 12 讲。
\end{itemize}
每讲正文大约需要 30--40 分钟。时间很紧时，可以只读要点、图和“一句话记住”。

\paragraph{各讲之间的关系}
下图是全部 12 讲的地图。箭头表示“后一讲要用到前一讲的东西”，虚线表示两讲之间有联系但不互相依赖。

\begin{figure}[H]
\centering
\begin{tikzpicture}[
  every node/.style={node box, minimum width=3.3cm, minimum height=.8cm},
  link/.style={-{Stealth[length=2.2mm,width=1.6mm]}, draw=mjInk!70, line width=.6pt},
  soft/.style={draw=mjGray, line width=.5pt, dash pattern=on 2.5pt off 1.8pt}]
  % 经济
  \node[draw=none, fill=none, font=\sffamily\bfseries\color{mjAccent}] at (0,0.9) {经济};
  \node (L1) at (0,0)    {1\ 供给与需求};
  \node (L2) at (0,-1.3) {2\ 弹性与剩余};
  \node (L3) at (0,-2.6) {3\ 税收};
  \node (L4) at (0,-3.9) {4\ 价格管制与外部性};
  \node (L5) at (4.2,0)    {5\ 博弈论入门};
  \node (L6) at (4.2,-1.3) {6\ 钱的时间价值};
  \node (L7) at (4.2,-3.9) {7\ 比较优势与贸易};
  % 社科
  \node[draw=none, fill=none, font=\sffamily\bfseries\color{mjSecond}] at (9.2,0.9) {社科};
  \node (L8)  at (9.2,0)    {8\ 投票与集体决策};
  \node (L9)  at (9.2,-1.3) {9\ 读懂数字};
  \node (L10) at (9.2,-2.6) {10\ 抽样调查};
  \node (L11) at (9.2,-3.9) {11\ 相关与因果};
  \node[minimum width=5cm] (L12) at (4.6,-5.4) {12\ 答题方法与练一练答案};
  \draw[link] (L1) -- (L2);
  \draw[link] (L2) -- (L3);
  \draw[link] (L3) -- (L4);
  \draw[link] (L6) -- node[draw=none, fill=none, minimum width=0pt, right, font=\footnotesize, inner sep=2pt] {贴现因子 $\delta$} (L5);
  \draw[link] (L9) -- (L10);
  \draw[link] (L10) -- (L11);
  \draw[soft] (L5) -- node[draw=none, fill=none, minimum width=0pt, above, font=\footnotesize, inner sep=2pt] {策略性投票} (L8);
  \draw[soft] (L4.north east) -- node[draw=none, fill=none, minimum width=0pt, right, font=\footnotesize, inner sep=2pt, pos=.35] {公地悲剧} (L5.south west);
  \draw[soft] (L9.west) to[bend right=18] node[draw=none, fill=none, minimum width=0pt, left, font=\footnotesize, inner sep=2pt, pos=.62] {辛普森悖论} (L11.west);
\end{tikzpicture}
\caption*{12 讲的关系}
\end{figure}

\noindent\begin{minipage}{\linewidth}
{\mojingheadfont\color{mjAccent}和训练卷的对应}\par
\begin{center}
\small
\begin{tabular}{lll}
  \toprule
  训练卷题目 & 主要用到 & 延伸阅读 \\
  \midrule
  经济第 1 题（税收由谁负担） & 第 1、2、3 讲 & 第 4 讲 \\
  经济第 2 题（囚徒困境与重复博弈） & 第 5 讲 & 第 6 讲的贴现 \\
  经济第 3 题（贷款的两种还款方式） & 第 6 讲 & — \\
  经济第 4 题（比较优势） & 第 7 讲 & — \\
  社科第 1 题（集体如何做决定） & 第 8 讲 & 第 5 讲 \\
  社科第 2 题（数据中的“悖论”） & 第 9 讲 & 第 11 讲 \\
  社科第 3 题（抽样调查） & 第 10 讲 & — \\
  社科第 4 题（相关与因果） & 第 11 讲 & 第 9 讲 \\
  \bottomrule
\end{tabular}
\end{center}
\end{minipage}

\paragraph{说明}
官方只公布了二试综合六选二的科目名称，没有公布经济、社科的考纲和题型。这份讲义的范围是按训练卷和高中《经济与社会》等课程的内容来定的，属于训练资料，不是官方材料。

% =============================================================================
\lecture[number=1, short=供给与需求]{供给与需求：价格是怎么定出来的}
% =============================================================================

\begin{keypoints}
\begin{itemize}
  \item 需求曲线向下倾斜：价格越高，想买的量越少。供给曲线向上倾斜：价格越高，想卖的量越多。
  \item 两条线的交点叫\textbf{均衡}。价格高于均衡会出现过剩，低于均衡会出现短缺，这两种力量都会把价格推回交点。
  \item 价格本身变了，是\textbf{沿着}曲线移动；价格以外的因素变了，是\textbf{整条}曲线移动。
  \item 分析题三步走：哪条线动？往哪边动？新的交点在哪里？
\end{itemize}
\end{keypoints}

\section{需求：价格和“想买多少”}

设想学校门口的一家奶茶店。经过一段时间的观察，店主发现价格和每天卖出的杯数大致是这样的：
\begin{center}
\small
\begin{tabular}{l*{6}{c}}
  \toprule
  价格 $P$（元） & 15 & 20 & 25 & 30 & 35 & 40 \\
  需求量 $Q$（杯） & 60 & 50 & 40 & 30 & 20 & 10 \\
  \bottomrule
\end{tabular}
\end{center}
这张表叫\textbf{需求表}。它正好满足 $Q=90-2P$。把它画出来，就是\textbf{需求曲线}（图~\ref{fig:es-dd-ss}(a)）。

价格越高，买的人越少，原因很朴素：
\begin{itemize}
  \item 每个顾客心里都有一个“最多愿意付多少”的价格。价格越高，超过这个心理价位的人越多，这些人就不买了；
  \item 价格高了，有人会改买别的饮料（这叫\textbf{替代}）。
\end{itemize}

需求曲线可以从两个方向读：
\begin{itemize}
  \item \textbf{横着读}：给定价格，读出想买多少。价格 30 元时，一天卖 30 杯。
  \item \textbf{竖着读}：给定数量，读出“最后那位顾客最多愿意付多少”。第 30 杯对应 30 元，意思是第 30 位顾客最多愿意付 30 元。
\end{itemize}
第二种读法在第 2 讲计算“消费者剩余”时要用到。

\begin{figure}[!htbp]
\centering
\begin{subfigure}{.48\textwidth}\centering
\begin{tikzpicture}[x=.075cm, y=.075cm]
  \econaxes{78}{52}{$Q$}{$P$}
  \draw[dem] (8,41) -- (72,9) node[right, slab, mjAccent] {需求 $D$};
  \foreach \q/\p in {60/15, 50/20, 40/25, 30/30, 20/35, 10/40}
    \node[eqpt, fill=mjAccent] at (\q,\p) {};
  \draw[guide] (0,30) node[left, slab] {30} -- (30,30) -- (30,0) node[below, slab] {30};
\end{tikzpicture}
\caption{需求曲线：向下倾斜}
\end{subfigure}\hfill
\begin{subfigure}{.48\textwidth}\centering
\begin{tikzpicture}[x=.075cm, y=.075cm]
  \econaxes{78}{52}{$Q$}{$P$}
  \draw[sup] (0,0) -- (47,47) node[right, slab, mjSecond] {供给 $S$};
  \foreach \q in {10,20,30,40}
    \node[eqpt, fill=mjSecond] at (\q,\q) {};
  \draw[guide] (0,30) node[left, slab] {30} -- (30,30) -- (30,0) node[below, slab] {30};
\end{tikzpicture}
\caption{供给曲线：向上倾斜}
\end{subfigure}
\caption{需求曲线和供给曲线。圆点是表中的数据。}
\label{fig:es-dd-ss}
\end{figure}

\begin{pitfall}
经济学习惯把\textbf{价格画在纵轴、数量画在横轴}，和数学里“自变量画在横轴”的习惯相反。所以需求 $Q=90-2P$ 在图上是直线 $P=45-\frac12Q$，
斜率是 $-\frac12$ 而不是 $-2$。读题、画图时先确认哪根轴是什么。
\end{pitfall}

\section{供给：价格和“想卖多少”}

再看卖方。价格越高，越多的店愿意卖奶茶，原有的店也愿意多做一些（加班、多雇人也划算了）。假设这个市场的\textbf{供给}是 $Q=P$：
价格 10 元时一天只供应 10 杯，价格 40 元时供应 40 杯（图~\ref{fig:es-dd-ss}(b)）。

供给曲线也可以竖着读：第 $Q$ 杯对应的价格，是卖方肯做这一杯的\textbf{最低要价}，也就是多做这一杯的成本。第 30 杯的成本是 30 元，价格低于 30 元，就没人愿意做第 30 杯。

\section{均衡：两条线的交点}

把两条线画在一起（图~\ref{fig:es-eq}）。想买的量等于想卖的量时，市场“出清”：
\[
  90-2P=P \quad\Longrightarrow\quad P^*=30,\qquad Q^*=30 .
\]
这一点叫\textbf{均衡}，$P^*$、$Q^*$ 叫均衡价格和均衡数量。

为什么价格会停在这里？看图中的两条虚线。
\begin{itemize}
  \item \textbf{价格偏高}（比如 36 元）：顾客只想买 18 杯，店家却想卖 36 杯，多出 18 杯卖不掉，叫\textbf{过剩}。卖不掉的店家会降价。
  \item \textbf{价格偏低}（比如 24 元）：顾客想买 42 杯，店家只肯做 24 杯，差 18 杯，叫\textbf{短缺}。排队买不到的顾客愿意多付钱，店家会涨价。
\end{itemize}
两种情况都会把价格推向 30 元。只有在交点上，没有人有动力改变价格。

\begin{figure}[!htbp]
\centering
\begin{tikzpicture}[x=.1cm, y=.09cm]
  \econaxes{88}{53}{数量 $Q$}{价格 $P$}
  \draw[dem] (0,45) -- (80,5) node[right, slab, mjAccent] {需求 $D$};
  \draw[sup] (0,0) -- (49,49) node[right, slab, mjSecond] {供给 $S$};
  % 过剩
  \draw[guide] (0,36) node[left, slab] {36} -- (36,36);
  \draw[decorate, decoration={brace, amplitude=4pt}, mjOchre, line width=.7pt] (18,37) -- (36,37)
    node[midway, above=4pt, slab, mjOchre] {过剩 18};
  \node[circle, fill=mjAccent, inner sep=1.4pt] at (18,36) {};
  \node[circle, fill=mjSecond, inner sep=1.4pt] at (36,36) {};
  % 短缺
  \draw[guide] (0,24) node[left, slab] {24} -- (42,24);
  \draw[decorate, decoration={brace, amplitude=4pt, mirror}, mjOchre, line width=.7pt] (24,23) -- (42,23)
    node[pos=.62, below=4pt, slab, mjOchre] {短缺 18};
  \node[circle, fill=mjSecond, inner sep=1.4pt] at (24,24) {};
  \node[circle, fill=mjAccent, inner sep=1.4pt] at (42,24) {};
  % 均衡
  \draw[guide] (0,30) node[left, slab] {$P^*=30$} -- (30,30) -- (30,0) node[below, slab] {$Q^*=30$};
  \node[eqpt] at (30,30) {};
  \node[slab, above=2pt] at (30,30) {$E$};
  % 价格调整箭头
  \draw[shiftarrow] (60,40) -- (60,33) node[midway, right, slab] {价格回落};
  \draw[shiftarrow] (60,18) -- (60,26) node[midway, right, slab] {价格上涨};
\end{tikzpicture}
\caption{均衡、过剩与短缺：价格偏离交点时，过剩和短缺会把它推回来}
\label{fig:es-eq}
\end{figure}

\section{曲线什么时候整体移动}

需求曲线描述的是“\textbf{其他条件不变}时，价格和想买的量的关系”。如果其他条件变了，整条曲线都会移动。
\begin{center}
\small
\begin{tabular}{p{6.6cm}p{6.6cm}}
  \toprule
  \textcolor{mjAccent}{\textbf{使需求曲线移动的因素}} & \textcolor{mjSecond}{\textbf{使供给曲线移动的因素}} \\
  \midrule
  收入变化 & 原料、工资等成本 \\
  相关商品的价格：替代品涨价，需求增加；互补品涨价，需求减少 & 生产技术 \\
  偏好、潮流、季节 & 卖家的数量 \\
  对未来价格的预期 & 对未来价格的预期 \\
  买家的数量 & 天气、自然灾害、政策 \\
  \bottomrule
\end{tabular}
\end{center}

需求“增加”指整条需求曲线向右移：在每个价格上，人们都想买得更多。
图~\ref{fig:es-shift}(a) 中，天气转热，奶茶的需求从 $Q=90-2P$ 增加到 $Q=105-2P$，新交点在 $P=35$，$Q=35$，价格和数量都上升。
图~\ref{fig:es-shift}(b) 中，牛奶涨价，每杯成本多了 15 元，供给从 $Q=P$ 减少为 $Q=P-15$（曲线向上、向左移），新交点在 $P=35$，$Q=20$，价格上升、数量下降。

\begin{figure}[!htbp]
\centering
\begin{subfigure}{.48\textwidth}\centering
\begin{tikzpicture}[x=.072cm, y=.072cm]
  \econaxes{80}{58}{$Q$}{$P$}
  \draw[dem] (0,45) -- (70,10) node[below, slab, mjAccent] {$D_1$};
  \draw[dem2] (0,52.5) -- (75,15) node[right, slab, mjAccent] {$D_2$};
  \draw[sup] (0,0) -- (50,50) node[above, slab, mjSecond] {$S$};
  \draw[shiftarrow] (52,19.5) -- (62,19.5);
  \node[eqpt] at (30,30) {}; \node[eqpt] at (35,35) {};
  \draw[guide] (0,30) node[left, slab] {30} -- (30,30) -- (30,0) node[below, slab] {30};
  \draw[guide] (0,35) node[left, slab] {35} -- (35,35) -- (35,0) node[below right=-1pt, slab] {35};
\end{tikzpicture}
\caption{需求增加：价格$\uparrow$，数量$\uparrow$}
\end{subfigure}\hfill
\begin{subfigure}{.48\textwidth}\centering
\begin{tikzpicture}[x=.072cm, y=.072cm]
  \econaxes{80}{58}{$Q$}{$P$}
  \draw[dem] (0,45) -- (70,10) node[below, slab, mjAccent] {$D$};
  \draw[sup] (0,0) -- (50,50) node[above, slab, mjSecond] {$S_1$};
  \draw[sup2] (0,15) -- (38,53) node[above, slab, mjSecond] {$S_2$};
  \draw[shiftarrow] (44,44) -- (36.5,51.5);
  \node[eqpt] at (30,30) {}; \node[eqpt] at (20,35) {};
  \draw[guide] (0,30) node[left, slab] {30} -- (30,30) -- (30,0) node[below, slab] {30};
  \draw[guide] (0,35) node[left, slab] {35} -- (20,35) -- (20,0) node[below, slab] {20};
\end{tikzpicture}
\caption{供给减少：价格$\uparrow$，数量$\downarrow$}
\end{subfigure}
\caption{曲线的整体移动。虚线是移动后的曲线。}
\label{fig:es-shift}
\end{figure}

\noindent\begin{minipage}{\linewidth}
\setlength{\parskip}{.35\baselineskip}
四种基本情形可以合成一张表：
\begin{center}
\small
\begin{tabular}{lcc|lcc}
  \toprule
  & 价格 & 数量 & & 价格 & 数量 \\
  \midrule
  需求增加（$D$ 右移） & $\uparrow$ & $\uparrow$ & 供给增加（$S$ 右移） & $\downarrow$ & $\uparrow$ \\
  需求减少（$D$ 左移） & $\downarrow$ & $\downarrow$ & 供给减少（$S$ 左移） & $\uparrow$ & $\downarrow$ \\
  \bottomrule
\end{tabular}
\end{center}
记法：\textbf{需求动，价量同向；供给动，价量反向。}
\end{minipage}

\begin{example}[两条线同时动]
夏天到了（奶茶需求增加），同时牛奶涨价（供给减少）。两件事都使价格上升，所以价格\textbf{一定}上升；
但需求增加使数量上升，供给减少使数量下降，数量的变化\textbf{不确定}，要看哪条线移动得多。
答题时应写“价格上升，数量变化方向取决于两条曲线移动的相对幅度”，而不是硬猜一个方向。
\end{example}

\begin{pitfall}
\begin{itemize}
  \item 区分“需求”和“需求量”。“价格上涨使需求下降”是错误说法：价格上涨使\textbf{需求量}下降，这是沿着曲线移动，需求曲线本身没动。
  \item 不要循环论证：“需求增加使价格上升，价格上升又使需求减少，所以……”是错的。正确做法是直接找新交点。
\end{itemize}
\end{pitfall}

\begin{oneline}
需求动，价量同向；供给动，价量反向。分析时永远先问：是哪条曲线在动？
\end{oneline}

\begin{exercises}[练一练]
  \item 某水果的需求为 $Q_d=120-4P$，供给为 $Q_s=2P$。求均衡价格和数量。如果价格被规定为 15 元，会出现过剩还是短缺？差多少？
  \item 判断下列事件使电动汽车市场的哪条曲线怎样移动，均衡价格和数量如何变化：
        (a) 电池成本大幅下降；(b) 汽油价格大幅上涨；(c) 两件事同时发生。
  \item 有人说：“这几年某城市房价一直在涨，成交量也在增加，可见价格上涨并没有减少人们买房的需求，需求曲线向下倾斜的说法不成立。”指出这个推理的错误。
\end{exercises}

% =============================================================================
\lecture[number=2, short=弹性与剩余]{弹性与剩余：谁更在乎价格，谁占了便宜}
% =============================================================================

\begin{keypoints}
\begin{itemize}
  \item \textbf{弹性}衡量“对价格有多敏感”：价格变化 1\%，数量变化百分之几。用百分比而不用斜率，是为了不受单位影响。
  \item 需求缺乏弹性时，降价会使总收入减少（“谷贱伤农”）；需求富有弹性时，降价会使总收入增加（“薄利多销”）。
  \item \textbf{消费者剩余} = 愿意付的 $-$ 实际付的；\textbf{生产者剩余} = 实际得到的 $-$ 最低肯接受的。在供求图上它们都是三角形的面积。
  \item 没有干预的竞争市场在均衡处成交，这时两种剩余之和最大。
\end{itemize}
\end{keypoints}

\section{弹性：对价格有多敏感}

食盐涨价一倍，你家买盐的量几乎不变；某个牌子的矿泉水涨价一倍，很多人就换别的牌子了。同样的涨价，数量的反应差别很大。
图~\ref{fig:es-elastic-two} 画出了两条需求曲线：价格同样从 30 元涨到 38 元，陡的那条只少卖 2 单位，平的那条少卖 16 单位。

\begin{figure}[!htbp]
\centering
\begin{subfigure}{.48\textwidth}\centering
\begin{tikzpicture}[x=.075cm, y=.075cm]
  \econaxes{70}{52}{$Q$}{$P$}
  \draw[dem] (25.5,48) -- (36.5,4) node[right, slab, mjAccent] {陡：缺乏弹性};
  \draw[guide] (0,38) node[left, slab] {38} -- (28,38) -- (28,0);
  \draw[guide] (0,30) node[left, slab] {30} -- (30,30) -- (30,0);
  \node[eqpt] at (30,30) {}; \node[eqpt] at (28,38) {};
  \draw[decorate, decoration={brace, amplitude=3pt, mirror}, mjOchre, line width=.7pt] (28,-4) -- (30,-4)
    node[midway, below=4pt, slab, mjOchre] {只少 2};
\end{tikzpicture}
\caption{食盐、常用药品}
\end{subfigure}\hfill
\begin{subfigure}{.48\textwidth}\centering
\begin{tikzpicture}[x=.075cm, y=.075cm]
  \econaxes{70}{52}{$Q$}{$P$}
  \draw[dem] (0,45) -- (66,12) node[below, slab, mjAccent] {平：富有弹性};
  \draw[guide] (0,38) node[left, slab] {38} -- (14,38) -- (14,0);
  \draw[guide] (0,30) node[left, slab] {30} -- (30,30) -- (30,0);
  \node[eqpt] at (30,30) {}; \node[eqpt] at (14,38) {};
  \draw[decorate, decoration={brace, amplitude=3pt, mirror}, mjOchre, line width=.7pt] (14,-4) -- (30,-4)
    node[midway, below=4pt, slab, mjOchre] {少 16};
\end{tikzpicture}
\caption{某个牌子的矿泉水}
\end{subfigure}
\caption{同样涨价 8 元，两种商品的需求量反应差别很大}
\label{fig:es-elastic-two}
\end{figure}

\begin{definition}[需求价格弹性]
\[
  \text{弹性}\ |E|=\frac{\text{需求量变化的百分比}}{\text{价格变化的百分比}}\quad\text{（取绝对值）}.
\]
$|E|>1$ 叫\textbf{富有弹性}（数量反应比价格变化更剧烈），$|E|<1$ 叫\textbf{缺乏弹性}，$|E|=1$ 叫单位弹性。
\end{definition}

为什么用百分比，而不直接用斜率？因为斜率和单位有关：价格用“元”还是“分”，数量用“杯”还是“箱”，斜率会差几十上百倍；百分比之比则不受单位影响，不同商品之间才能比较。

\paragraph{同一条直线上，弹性处处不同}
以第 1 讲的需求 $Q=90-2P$ 为例。
\begin{itemize}
  \item 在 $P=30$、$Q=30$ 处：价格涨 1 元（涨 $3.3\%$），数量从 30 降到 28（降 $6.7\%$），$|E|\approx2$，富有弹性。
  \item 在 $P=10$、$Q=70$ 处：价格同样涨 1 元（涨 $10\%$），数量从 70 降到 68（只降 $2.9\%$），$|E|\approx0.29$，缺乏弹性。
\end{itemize}
同样是“涨 1 元少卖 2 杯”，在高价处这 2 杯占比大，在低价处占比小。对直线需求 $Q=a-bP$，某一点的弹性是 $|E|=b\cdot\dfrac PQ$，
直线的中点处 $|E|=1$，中点以上富有弹性，以下缺乏弹性（图~\ref{fig:es-elastic-line}）。

\begin{figure}[!htbp]
\centering
\begin{tikzpicture}[x=.085cm, y=.085cm]
  \econaxes{98}{50}{$Q$}{$P$}
  \draw[dem, line width=2.4pt, mjAccent] (0,45) -- (45,22.5);
  \draw[dem, line width=2.4pt, mjAccent!40] (45,22.5) -- (90,0);
  \node[eqpt] at (45,22.5) {};
  \node[slab, anchor=south west] at (46,23) {中点：$|E|=1$};
  \node[slab, mjAccent, anchor=south west] at (12,40.5) {上半段：$|E|>1$，富有弹性};
  \node[slab, mjAccent!70, anchor=south west] at (62,12) {下半段：$|E|<1$，缺乏弹性};
  \node[eqpt] at (30,30) {}; \node[slab, anchor=north east] at (29.5,29.5) {$|E|=2$};
  \node[eqpt] at (70,10) {}; \node[slab, anchor=north east] at (69.5,9.5) {$|E|\approx0.29$};
  \draw[guide] (0,22.5) node[left, slab] {22.5} -- (45,22.5) -- (45,0) node[below, slab] {45};
  \node[slab, left] at (0,45) {45}; \node[slab, below] at (90,0) {90};
\end{tikzpicture}
\caption{直线需求 $Q=90-2P$ 上弹性的变化}
\label{fig:es-elastic-line}
\end{figure}

\paragraph{什么决定弹性}
\begin{itemize}
  \item \textbf{替代品多不多}：某个牌子的矿泉水有很多替代品，弹性大；食盐几乎没有替代品，弹性小。
  \item \textbf{是必需品还是可有可无}：药品、粮食弹性小，旅游、奢侈品弹性大。
  \item \textbf{占收入的比重}：铅笔涨价一倍你不在乎，房租涨价一倍你会想办法。
  \item \textbf{时间长短}：油价上涨，短期内只能照样开车（弹性小）；长期可以换小排量的车、搬家（弹性大）。
\end{itemize}

\section{弹性与总收入：谷贱伤农}

卖家的总收入等于价格乘数量，$P\times Q$，在图上就是一个矩形的面积。降价时，单价变低（矩形变矮），卖出更多（矩形变宽），总收入是升是降取决于弹性。

粮食是典型的缺乏弹性的商品：粮价降一半，人们也不会多吃一倍的饭。图~\ref{fig:es-grain} 中，需求曲线为 $P=6-0.04Q$，粮食收获后产量固定，供给是一条竖线。
正常年景产量 100，价格 2，农民总收入 200；丰收年产量增加到 120（多 20\%），价格却跌到 1.2（跌 40\%），总收入只剩 144。
多出来的粮食只能靠大幅降价才卖得出去，这就是“谷贱伤农”。

\begin{figure}[!htbp]
\centering
\begin{subfigure}{.48\textwidth}\centering
\begin{tikzpicture}[x=.031cm, y=.75cm]
  \fill[taxfill] (0,0) rectangle (100,2);
  \econaxes{160}{5.2}{产量 $Q$}{价格 $P$}
  \draw[dem] (25,5) -- (145,0.2) node[above right=-2pt, slab, mjAccent] {$D$};
  \draw[sup] (100,0) -- (100,4.6) node[above, slab, mjSecond] {$S_1$};
  \node[eqpt] at (100,2) {};
  \node[slab, left] at (0,2) {2};
  \node[slab, below] at (100,0) {100};
  \node[lab] at (50,1) {收入 200};
\end{tikzpicture}
\caption{正常年景}
\end{subfigure}\hfill
\begin{subfigure}{.48\textwidth}\centering
\begin{tikzpicture}[x=.031cm, y=.75cm]
  \fill[taxfill] (0,0) rectangle (120,1.2);
  \econaxes{160}{5.2}{产量 $Q$}{价格 $P$}
  \draw[dem] (25,5) -- (145,0.2) node[above right=-2pt, slab, mjAccent] {$D$};
  \draw[sup] (120,0) -- (120,4.6) node[above, slab, mjSecond] {$S_2$};
  \draw[guide] (100,0) -- (100,2) -- (0,2);
  \node[eqpt] at (120,1.2) {};
  \node[slab, left] at (0,1.2) {1.2};
  \node[slab, below] at (120,0) {120};
  \node[lab] at (55,.6) {收入 144};
\end{tikzpicture}
\caption{丰收年}
\end{subfigure}
\caption{谷贱伤农：需求缺乏弹性时，产量增加反而使总收入减少}
\label{fig:es-grain}
\end{figure}

一般规律：
\begin{center}
\small
\begin{tabular}{lll}
  \toprule
  需求 & 降价时总收入 & 例子 \\
  \midrule
  缺乏弹性（$|E|<1$） & 减少 & 谷贱伤农；药品、门票涨价能增收 \\
  富有弹性（$|E|>1$） & 增加 & 薄利多销；促销打折 \\
  \bottomrule
\end{tabular}
\end{center}

\section{消费者剩余}

假设有 5 位顾客，各自最多愿意为一杯奶茶付 50、40、30、20、10 元。现在价格是 25 元：
\begin{itemize}
  \item 前三位会买。第一位本来肯付 50 元，只付了 25 元，等于“省下”25 元；第二位省下 15 元；第三位省下 5 元。
  \item 后两位不买，什么也没省下，也没损失。
\end{itemize}
这些“省下的钱”加起来是 $25+15+5=45$ 元，叫\textbf{消费者剩余}。在图~\ref{fig:es-stair} 中，它就是台阶高出价格线的那几块面积。

\begin{figure}[!htbp]
\centering
\begin{tikzpicture}[x=1.15cm, y=.068cm]
  \foreach \i/\h/\n in {0/50/甲, 1/40/乙, 2/30/丙, 3/20/丁, 4/10/戊}{
    \draw[mj thin, fill=mjAccent!6] (\i,0) rectangle (\i+1,\h);
    \node[slab, below] at (\i+.5,0) {\n};
    \node[slab, above] at (\i+.5,\h) {\h};
  }
  \foreach \i/\h in {0/50, 1/40, 2/30}
    \fill[csfill] (\i,25) rectangle (\i+1,\h);
  \foreach \i/\h in {0/50, 1/40, 2/30}
    \draw[mj thin] (\i,0) rectangle (\i+1,\h);
  \draw[mj axis] (0,0) -- (5.8,0) node[right, lab] {顾客};
  \draw[mj axis] (0,0) -- (0,58) node[above, lab] {最多愿付（元）};
  \draw[mjSecond, line width=1pt] (0,25) -- (5.5,25) node[right, slab, mjSecond] {价格 25};
  \node[slab] at (.5,37.5) {25}; \node[slab] at (1.5,32.5) {15}; \node[slab] at (2.5,27.5) {5};
\end{tikzpicture}
\caption{消费者剩余：每位买家“愿意付的”减去“实际付的”，加起来 $25+15+5=45$}
\label{fig:es-stair}
\end{figure}

顾客很多时，台阶变得很细，连成一条光滑的需求曲线，消费者剩余就是\textbf{需求曲线以下、价格线以上}的面积。直线需求时，它是一个三角形。

\section{生产者剩余}

卖方这边完全对称。供给曲线竖着读，是每一杯的最低要价（成本）。价格 30 元时，成本 10 元的那一杯净赚 20 元，成本 25 元的那一杯净赚 5 元……
把这些“多得到的”加起来，就是\textbf{生产者剩余}：\textbf{价格线以下、供给曲线以上}的面积。

图~\ref{fig:es-surplus}(a) 画出了第 1 讲的市场在均衡时的两块剩余：
\[
  \text{消费者剩余}=\tfrac12\times30\times(45-30)=225,\qquad
  \text{生产者剩余}=\tfrac12\times30\times30=450 .
\]
两者之和 675 叫\textbf{总剩余}，表示这个市场的交易给买卖双方一共带来了多少好处。

\section{为什么说均衡“有效率”}

假设由于某种原因，每天只能卖 20 杯（图~\ref{fig:es-surplus}(b)）。第 21 到第 30 杯本来是“值得做”的：
买家愿意付的（需求曲线的高度）高于做它的成本（供给曲线的高度）。不做这 10 杯，中间那块三角形的好处就消失了，面积是
$\frac12\times(35-20)\times(30-20)=75$。

反过来，如果卖到 30 杯以上，多出的那些杯成本高于买家愿意付的价格，也是浪费。所以\textbf{均衡数量使总剩余最大}。
这块因为成交量偏离均衡而消失的三角形，叫\textbf{无谓损失}，第 3 讲会反复用到它。

\begin{figure}[!htbp]
\centering
\begin{subfigure}{.48\textwidth}\centering
\begin{tikzpicture}[x=.075cm, y=.075cm]
  \fill[csfill] (0,45) -- (30,30) -- (0,30) -- cycle;
  \fill[psfill] (0,0) -- (30,30) -- (0,30) -- cycle;
  \econaxes{72}{52}{$Q$}{$P$}
  \draw[dem] (0,45) -- (66,12) node[below, slab, mjAccent] {$D$};
  \draw[sup] (0,0) -- (46,46) node[above, slab, mjSecond] {$S$};
  \draw[guide] (30,30) -- (30,0) node[below, slab] {30};
  \node[slab, left] at (0,30) {30}; \node[slab, left] at (0,45) {45};
  \node[eqpt] at (30,30) {};
  \node[slab, mjAccent] at (10,35) {225};
  \node[slab, mjSecond] at (10,22) {450};
\end{tikzpicture}
\caption{蓝：消费者剩余；红：生产者剩余}
\end{subfigure}\hfill
\begin{subfigure}{.48\textwidth}\centering
\begin{tikzpicture}[x=.075cm, y=.075cm]
  \fill[mjLight] (0,45) -- (20,35) -- (20,20) -- (0,0) -- cycle;
  \fill[dwlfill] (20,35) -- (30,30) -- (20,20) -- cycle;
  \econaxes{72}{52}{$Q$}{$P$}
  \draw[dem] (0,45) -- (66,12) node[below, slab, mjAccent] {$D$};
  \draw[sup] (0,0) -- (46,46) node[above, slab, mjSecond] {$S$};
  \draw[guide] (20,35) -- (20,0) node[below, slab] {20};
  \draw[guide] (30,30) -- (30,0) node[below, slab] {30};
  \draw[guide] (0,35) node[left, slab] {35} -- (20,35);
  \draw[guide] (0,20) node[left, slab] {20} -- (20,20);
  \node[slab] at (8,22) {600};
  \node[slab, anchor=west, mjPine] at (36,8) {损失 75};
  \draw[mjPine, line width=.4pt] (36,9) -- (24,27);
\end{tikzpicture}
\caption{只卖 20 杯：总剩余只剩 600}
\end{subfigure}
\caption{总剩余在均衡处最大。偏离均衡，就会损失一块三角形。}
\label{fig:es-surplus}
\end{figure}

\begin{pitfall}
\begin{itemize}
  \item 剩余是“多得到的好处”，不是花了多少钱、赚了多少钱。消费者剩余不是消费者的支出，生产者剩余也不等于利润（固定成本没有扣除）。
  \item “总剩余最大”说的是\textbf{效率}，不是\textbf{公平}。蛋糕最大，不代表分得合理；这是两个不同的问题。
\end{itemize}
\end{pitfall}

\begin{oneline}
弹性看百分比，剩余看面积；交易创造的好处是需求曲线和供给曲线之间的那块面积。
\end{oneline}

\begin{exercises}[练一练]
  \item 需求 $Q_d=120-4P$，供给 $Q_s=2P$（与练习 1.1 相同）。求均衡时的消费者剩余和生产者剩余。
  \item 某景区门票 100 元时每天 2000 名游客，涨到 120 元后每天 1800 名。用“变化量 $\div$ 原值”估算需求弹性。涨价后门票总收入是增加还是减少？
  \item 同一航班，提前一个月订票和起飞前一天订票，价格常常差很多。用弹性解释航空公司为什么这样定价。
\end{exercises}

% =============================================================================
\lecture[number=3, short=税收]{税收：谁真正掏了钱}
% =============================================================================

\begin{keypoints}
\begin{itemize}
  \item 每件征税 $t$ 元，相当于在“买方付的价格”和“卖方得到的价格”之间打进一个宽 $t$ 的\textbf{楔子}。
  \item 法律上向谁征税不影响结果：向买方征和向卖方征，最终的价格和成交量完全一样。
  \item 税负更多地落在\textbf{弹性小}（对价格不敏感、“跑不掉”）的一方。
  \item 征税使成交量减少，损失一块三角形（无谓损失）。每件的税额翻倍，这块三角形变成 4 倍。
\end{itemize}
\end{keypoints}

\section{征税以后发生了什么}

仍用第 1 讲的奶茶市场：需求 $Q=90-2P$，供给 $Q=P$，原来的均衡是 $P=30$，$Q=30$。
现在政府规定，卖方每卖一杯要交 6 元税。

卖方关心的是\textbf{到手价}。标价 $P$ 时，卖方到手只有 $P-6$，所以供给变为 $Q=P-6$。新的均衡满足
\[
  90-2P=P-6 \quad\Longrightarrow\quad P=32,\qquad Q=26 .
\]
\begin{itemize}
  \item 买方每杯付 32 元，比原来多付 \textbf{2 元}；
  \item 卖方每杯到手 $32-6=26$ 元，比原来少得 \textbf{4 元}；
  \item 政府每杯收 6 元，共收税 $6\times26=156$ 元。
\end{itemize}
6 元的税，买方负担了 2 元，卖方负担了 4 元。虽然税是卖方交的，买方也分担了一部分。

\begin{figure}[!htbp]
\centering
\begin{tikzpicture}[x=.13cm, y=.094cm]
  \fill[csfill] (0,45) -- (26,32) -- (0,32) -- cycle;
  \fill[taxfill] (0,26) rectangle (26,32);
  \fill[psfill] (0,0) -- (26,26) -- (0,26) -- cycle;
  \fill[dwlfill] (26,32) -- (30,30) -- (26,26) -- cycle;
  \econaxes{72}{52}{数量 $Q$}{价格 $P$}
  \draw[dem] (0,45) -- (66,12) node[below, slab, mjAccent] {需求 $D$};
  \draw[sup] (0,0) -- (47,47) node[right, slab, mjSecond] {供给 $S$};
  \draw[sup2] (0,6) -- (41,47) node[above, slab, mjSecond] {$S+6$（含税）};
  \draw[shiftarrow] (40,40) -- (36,46) node[pos=.1, right=3pt, slab] {上移 6};
  \draw[guide] (0,30) -- (30,30) -- (30,0) node[below, slab] {30};
  \draw[guide] (26,26) -- (26,0) node[below, slab] {26};
  \node[slab, left] (lb) at (-2,37) {买方付 32}; \draw[mjInk!60, line width=.3pt] (lb.east) -- (0,32);
  \node[slab, left] at (-2,30) {原价 30};       \draw[mjInk!60, line width=.3pt] (-2,30) -- (0,30);
  \node[slab, left] (ls) at (-2,21) {卖方得 26}; \draw[mjInk!60, line width=.3pt] (ls.east) -- (0,26);
  \node[slab, left] at (0,45) {45};
  \node[eqpt] at (26,32) {}; \node[eqpt] at (26,26) {}; \node[circle, draw=mjInk, fill=white, inner sep=1.3pt] at (30,30) {};
  \node[slab] at (13,29) {税收 156};
  \node[slab, mjAccent] at (7,35.5) {消费者剩余};
  \node[slab, mjSecond] at (8,19) {生产者剩余};
  \node[slab, anchor=west, mjPine] at (50,30) {无谓损失 12};
  \draw[mjPine, line width=.4pt] (50,30) -- (28.5,29.5);
\end{tikzpicture}
\caption{对卖方每件征税 6 元。浅棕色矩形是税收，绿色三角形是无谓损失。}
\label{fig:es-tax}
\end{figure}

图~\ref{fig:es-tax} 是这件事的全貌：
\begin{itemize}
  \item 含税供给曲线比原来\textbf{高出 6 元}：要让卖方到手同样多，标价必须高 6 元；
  \item 买方价格 32 和卖方价格 26 之间隔着 6 元，就是“楔子”；
  \item 原来的总剩余（675）现在分成了四块：消费者剩余 169、生产者剩余 338、税收 156、无谓损失 12。$169+338+156+12=675$。
\end{itemize}

\section{向买方征税，结果完全一样}

如果改为买方每买一杯交 6 元税呢？买方付的总代价是“标价 $+6$”，愿意买的量变为 $Q=90-2(P+6)$，需求曲线\textbf{向下}移 6 元。新均衡：
\[
  90-2(P+6)=P \quad\Longrightarrow\quad P=26 .
\]
这里的 $P=26$ 是标价，也就是卖方到手的钱；买方另付 6 元税，总共付 32 元。和向卖方征税的结果\textbf{一模一样}（图~\ref{fig:es-tax-who}）。

\begin{figure}[!htbp]
\centering
\begin{subfigure}{.48\textwidth}\centering
\begin{tikzpicture}[x=.075cm, y=.075cm]
  \econaxes{72}{52}{$Q$}{$P$}
  \draw[dem] (0,45) -- (66,12) node[below, slab, mjAccent] {$D$};
  \draw[sup] (0,0) -- (46,46) node[right, slab, mjSecond] {$S$};
  \draw[sup2] (0,6) -- (40,46) node[above, slab, mjSecond] {$S+6$};
  \draw[guide] (0,32) node[left, slab] {32} -- (26,32) -- (26,0) node[below, slab] {26};
  \draw[guide] (0,26) node[left, slab] {26} -- (26,26);
  \node[eqpt] at (26,32) {}; \node[eqpt] at (26,26) {};
\end{tikzpicture}
\caption{向卖方征税：供给曲线上移 6}
\end{subfigure}\hfill
\begin{subfigure}{.48\textwidth}\centering
\begin{tikzpicture}[x=.075cm, y=.075cm]
  \econaxes{72}{52}{$Q$}{$P$}
  \draw[dem] (0,45) -- (66,12) node[above right=-2pt, slab, mjAccent] {$D$};
  \draw[dem2] (0,39) -- (66,6) node[below, slab, mjAccent] {$D-6$};
  \draw[sup] (0,0) -- (46,46) node[right, slab, mjSecond] {$S$};
  \draw[guide] (0,32) node[left, slab] {32} -- (26,32) -- (26,0) node[below, slab] {26};
  \draw[guide] (0,26) node[left, slab] {26} -- (26,26);
  \node[eqpt] at (26,32) {}; \node[eqpt] at (26,26) {};
\end{tikzpicture}
\caption{向买方征税：需求曲线下移 6}
\end{subfigure}
\caption{两种征税方式得到同样的买方价格 32、卖方价格 26、成交量 26}
\label{fig:es-tax-who}
\end{figure}

道理很简单：真正起作用的只有一件事——\textbf{买方付的比卖方得的多出 $t$ 元}。至于这 $t$ 元由谁递给税务局，只是手续问题。

\section{谁负担得多：看谁“跑得掉”}

一般地，设需求 $Q=a-bP$，供给 $Q=c+dP$（$b,d>0$），每件征税 $t$。设买方价格比原来\textbf{上升} $x$，卖方价格比原来\textbf{下降} $y$。
\begin{itemize}
  \item 楔子的宽度是 $t$：$x+y=t$；
  \item 买方因涨价少买 $bx$，卖方因降价少卖 $dy$，新的成交量对双方相同，所以 $bx=dy$。
\end{itemize}
解得
\[
  \boxed{\ x=\frac{d}{b+d}\,t,\qquad y=\frac{b}{b+d}\,t\ } .
\]
检验：$b=2$，$d=1$，$t=6$ 时，$x=2$，$y=4$，和上一节一致。

\textbf{怎么理解}：$b$ 大表示买方对价格敏感，稍一涨价就不买了（容易“跑掉”），于是 $x$ 小，买方负担少。
所以\textbf{谁对价格越不敏感，谁负担越多}。严格地说，比较的是弹性：在原均衡点，需求弹性是 $b\cdot\frac PQ$，供给弹性是 $d\cdot\frac PQ$，
两者之比正好是 $b:d$，所以 $x:y=d:b$ 就是“供给弹性 : 需求弹性”。重要的是两边弹性的\textbf{相对}大小。图~\ref{fig:es-incidence} 画出了两种情形。

\begin{figure}[!htbp]
\centering
\begin{subfigure}{.48\textwidth}\centering
\begin{tikzpicture}[x=.075cm, y=.066cm]
  \econaxes{66}{56}{$Q$}{$P$}
  \draw[dem] (22,49.2) node[left, slab, mjAccent] {$D$（陡）} -- (40,6);
  \draw[sup] (0,12) -- (60,48) node[right, slab, mjSecond] {$S$};
  \draw[sup2] (0,27) -- (41.7,52) node[right, slab, mjSecond] {$S+t$};
  \draw[guide] (0,42) node[left, slab] {$P_{\text{买}}$} -- (25,42);
  \draw[guide] (0,30) node[left, slab, anchor=south east] {$P^*$} -- (30,30);
  \draw[guide] (0,27) node[left, slab, anchor=north east] {$P_{\text{卖}}$} -- (25,27);
  \draw[mjAccent, line width=3pt] (25,30) -- (25,42);
  \draw[mjSecond, line width=3pt] (25,27) -- (25,30);
  \node[slab, mjAccent, anchor=west] at (44,10) {买方负担大};
\end{tikzpicture}
\caption{需求比供给更缺乏弹性：买方负担大部分}
\end{subfigure}\hfill
\begin{subfigure}{.48\textwidth}\centering
\begin{tikzpicture}[x=.075cm, y=.066cm]
  \econaxes{66}{56}{$Q$}{$P$}
  \draw[dem] (0,48) -- (65,9) node[below, slab, mjAccent] {$D$};
  \draw[sup] (19,3.6) -- (38,49.2) node[right, slab, mjSecond] {$S$（陡）};
  \draw[sup2] (12.5,3) -- (32.9,52) node[left, slab, mjSecond] {$S+t$};
  \draw[guide] (0,33) node[left, slab, anchor=south east] {$P_{\text{买}}$} -- (25,33);
  \draw[guide] (0,30) node[left, slab, anchor=north east] {$P^*$} -- (30,30);
  \draw[guide] (0,18) node[left, slab] {$P_{\text{卖}}$} -- (25,18);
  \draw[mjAccent, line width=3pt] (25,30) -- (25,33);
  \draw[mjSecond, line width=3pt] (25,18) -- (25,30);
  \node[slab, mjSecond, anchor=west] at (42,44) {卖方负担大};
\end{tikzpicture}
\caption{需求比供给更富有弹性：卖方负担大部分}
\end{subfigure}
\caption{税负的分配。粗线中，蓝色段是买方多付的，红色段是卖方少得的，两段加起来正好是 $t$。}
\label{fig:es-incidence}
\end{figure}

两个常被引用的例子：
\begin{itemize}
  \item \textbf{香烟税}：吸烟者对价格不敏感（成瘾），税负主要落在吸烟者身上。
  \item \textbf{奢侈品税}：美国 1990 年对游艇等奢侈品加征消费税。富人很容易把钱花在别处或到国外买（弹性大），游艇销量大跌，
        税负很大程度上落到了造船厂和工人身上；对游艇等的这部分税在 1993 年就被废除了。这是教科书里常用的例子。
\end{itemize}

\section{无谓损失：税收的“隐形成本”}

税收本身不是社会的损失：买卖双方少了的那部分钱，政府收走了，可以用来修路、办学校。真正的损失是图~\ref{fig:es-tax} 中的\textbf{绿色三角形}：
因为加了税，第 27 到第 30 杯不再成交了，而这几杯本来是买家觉得值、卖家也愿意做的。在例子中，它的面积是
\[
  \text{无谓损失}=\tfrac12\times6\times(30-26)=12 .
\]

一般地，成交量减少 $bx=\dfrac{bd}{b+d}\,t$，所以
\[
  \text{无谓损失}=\frac12\,t\cdot\frac{bd}{b+d}\,t=\frac{bd}{2(b+d)}\,t^2 .
\]
无谓损失和\textbf{每件税额的平方}成正比：税额翻倍，损失变成 4 倍。图~\ref{fig:es-laffer} 画出了例子中的税收和无谓损失随每杯税额 $t$ 的变化：
税收 $R(t)=t\left(30-\frac23t\right)$，无谓损失 $\frac13t^2$。

\begin{figure}[!htbp]
\centering
\begin{tikzpicture}
\begin{axis}[mojing, width=9.5cm, height=5.2cm, xmin=0, xmax=48, ymin=0, ymax=420,
    xlabel={每杯税额 $t$（元）}, ylabel={元}, xtick={6,22.5,30,45}, xticklabels={6,22.5,30,45},
    ytick={156,337.5}, yticklabels={156,337.5}, legend style={at={(1.03,.5)}, anchor=west}]
  \addplot[mjOchre, domain=0:45, samples=60, line width=1pt] {x*(30-2*x/3)};
  \addplot[mjPine, domain=0:34.6, samples=60, line width=1pt] {x^2/3};
  \legend{税收,无谓损失}
  \draw[guide] (axis cs:22.5,0) -- (axis cs:22.5,337.5) -- (axis cs:0,337.5);
  \draw[guide] (axis cs:6,0) -- (axis cs:6,156) -- (axis cs:0,156);
  \node[eqpt] at (axis cs:22.5,337.5) {};
  \node[eqpt] at (axis cs:6,156) {};
  \node[eqpt] at (axis cs:30,300) {};
  \node[slab, anchor=west] at (axis cs:31,300) {$t=30$ 时两者相等};
\end{axis}
\end{tikzpicture}
\caption{每杯税额提高时，税收先增后减，无谓损失却一直加速增长}
\label{fig:es-laffer}
\end{figure}

图上有两个值得记住的现象：
\begin{itemize}
  \item 税收在 $t=22.5$ 时最多。税额再高，成交量萎缩得太厉害，税收反而减少；$t=45$ 时没人交易，税收为零（这就是“拉弗曲线”的想法）。
  \item 无谓损失增长得比税收快。所以经济学家通常认为：\textbf{税基宽、税率低}，比税基窄、税率高的效率损失小。
\end{itemize}

\begin{remark}
\textbf{补贴}可以看作“负的税”：买方价格下降、卖方价格上升，成交量超过均衡量。多生产的那部分成本高于买家的评价，同样会产生无谓损失。
不过，如果市场本身有问题（例如第 4 讲的外部性），税收和补贴反而能把成交量拉回到合适的位置。
\end{remark}

\begin{pitfall}
\begin{itemize}
  \item “税是向卖方征的，所以卖方负担”——错。税负分配只取决于供求两边的弹性，与法律上谁交税无关。
  \item “政府收的税就是社会的损失”——错。税收是从买卖双方转给政府的，只是换了主人；真正消失的只有无谓损失那块三角形。
  \item 算税负时，买方多付的 $x$ 和卖方少得的 $y$ 加起来一定等于 $t$，可以用来检查计算。
\end{itemize}
\end{pitfall}

\begin{oneline}
税是一个楔子：谁对价格越不敏感，谁负担越多；损失的是楔子挤掉的那块三角形。
\end{oneline}

\begin{exercises}[练一练]
  \item 需求 $Q_d=120-4P$，供给 $Q_s=2P$。政府对\textbf{买方}每件征税 3 元。求买方付出的总价格、卖方到手的价格、成交量、税收、双方各负担多少，以及无谓损失。
  \item 证明：在直线供求下，每件税额从 $t$ 提高到 $2t$（成交量仍为正）时，无谓损失变为原来的 4 倍，而税收不到原来的 2 倍。
  \item 有人建议：“为了让富人多交税，应该对豪华跑车征收重税。”用本讲的知识分析，这项税的负担可能落在谁身上？还需要知道什么才能判断？
\end{exercises}

% =============================================================================
\lecture[number=4, short=价格管制与外部性]{价格管制与外部性：政府什么时候该出手}
% =============================================================================

\begin{keypoints}
\begin{itemize}
  \item \textbf{最高限价}（低于均衡价）造成短缺：排队、托关系、黑市、质量下降。
  \item \textbf{最低限价}（高于均衡价）造成过剩：卖不掉的商品、找不到工作的人。
  \item \textbf{外部性}：交易影响了没有参与交易的第三方。污染使社会成本高于企业自己的成本，市场会生产得太多。
  \item 对每件产品征收等于外部成本的税（\textbf{庇古税}），能把产量拉回到对社会最好的水平。这时税收不但不造成损失，反而消除了损失。
\end{itemize}
\end{keypoints}

\section{最高限价：便宜了，却买不到}

政府有时担心价格太高，规定“不许超过某个价格”，这叫\textbf{最高限价}。只有当限价\textbf{低于}均衡价格时，它才起作用。

还是奶茶市场（需求 $Q=90-2P$，供给 $Q=P$）。若规定每杯最多卖 24 元（图~\ref{fig:es-control}(a)）：
\begin{itemize}
  \item 顾客想买 42 杯，店家只肯做 24 杯，\textbf{短缺} 18 杯；
  \item 成交量由少的一方决定，只成交 24 杯，比均衡少 6 杯。少成交的这 6 杯造成无谓损失（图中绿色三角形），面积为 $\frac12\times9\times6=27$；
  \item 谁能买到这 24 杯？只能靠排队、先到先得、关系或抽签。排队的时间也是成本；有人会高价倒卖（黑市）；店家可能偷偷减料（质量下降）。
\end{itemize}
节假日的热门车票、门票“一票难求”，背后就是价格被压在均衡价之下，于是只好用排队、抢票来分配。

\section{最低限价：贵了，却卖不掉}

反过来，规定“不许低于某个价格”叫\textbf{最低限价}，只有\textbf{高于}均衡价格时才起作用。若规定每杯最低 36 元（图~\ref{fig:es-control}(b)），
店家想卖 36 杯，顾客只买 18 杯，\textbf{过剩} 18 杯。

最低限价的例子：
\begin{itemize}
  \item \textbf{粮食最低收购价}：为保护农民收入，政府常常要按这个价格把卖不掉的粮食收购下来，并承担储存的成本。
  \item \textbf{最低工资}：在劳动市场里，“价格”是工资，“过剩的劳动”就是失业。按这个简单模型，最低工资定得高于市场工资时，会使一部分人找不到工作。
        但现实中的实证研究结论并不一致，影响大小取决于最低工资比市场工资高多少，以及雇主在劳动市场中的地位等。模型告诉我们应当关注哪种力量，影响多大还要看数据。
\end{itemize}

\begin{figure}[!htbp]
\centering
\begin{subfigure}{.48\textwidth}\centering
\begin{tikzpicture}[x=.075cm, y=.075cm]
  \fill[dwlfill] (24,33) -- (30,30) -- (24,24) -- cycle;
  \econaxes{72}{52}{$Q$}{$P$}
  \draw[dem] (0,45) -- (66,12) node[below, slab, mjAccent] {$D$};
  \draw[sup] (0,0) -- (46,46) node[right, slab, mjSecond] {$S$};
  \draw[mjOchre, line width=1.1pt] (0,24) -- (60,24) node[above, slab, mjOchre] {限价 24};
  \draw[guide] (24,24) -- (24,0) node[below, slab] {24};
  \draw[guide] (42,24) -- (42,0) node[below, slab] {42};
  \draw[decorate, decoration={brace, amplitude=4pt, mirror}, mjOchre, line width=.7pt] (24,21.5) -- (42,21.5)
    node[midway, below=4pt, slab, mjOchre] {短缺 18};
  \node[circle, draw=mjInk, fill=white, inner sep=1.3pt] at (30,30) {};
\end{tikzpicture}
\caption{最高限价 24 元：短缺}
\end{subfigure}\hfill
\begin{subfigure}{.48\textwidth}\centering
\begin{tikzpicture}[x=.075cm, y=.075cm]
  \fill[dwlfill] (18,36) -- (30,30) -- (18,18) -- cycle;
  \econaxes{72}{52}{$Q$}{$P$}
  \draw[dem] (0,45) -- (66,12) node[below, slab, mjAccent] {$D$};
  \draw[sup] (0,0) -- (46,46) node[right, slab, mjSecond] {$S$};
  \draw[mjOchre, line width=1.1pt] (0,36) -- (60,36) node[above, slab, mjOchre] {限价 36};
  \draw[guide] (18,36) -- (18,0) node[below, slab] {18};
  \draw[guide] (36,36) -- (36,0) node[below, slab] {36};
  \draw[decorate, decoration={brace, amplitude=4pt}, mjOchre, line width=.7pt] (18,38.5) -- (36,38.5)
    node[midway, above=4pt, slab, mjOchre] {过剩 18};
  \node[circle, draw=mjInk, fill=white, inner sep=1.3pt] at (30,30) {};
\end{tikzpicture}
\caption{最低限价 36 元：过剩}
\end{subfigure}
\caption{价格管制。空心点是原来的均衡，绿色三角形是成交量减少造成的无谓损失。}
\label{fig:es-control}
\end{figure}

\begin{pitfall}
限价只有在“有约束”时才起作用：最高限价定在均衡价\textbf{以上}、最低限价定在均衡价\textbf{以下}时，市场照常在均衡处成交，什么也不会发生。答题时先判断限价在均衡价的哪一侧。
\end{pitfall}

\section{外部性：没有写进价格的成本}

\begin{definition}[外部性]
一项经济活动给\textbf{没有参与交易}的人带来了成本或好处，而这些成本或好处没有反映在价格里，就叫外部性。带来成本的叫负外部性，带来好处的叫正外部性。
\end{definition}

\begin{itemize}
  \item 负外部性：工厂排污、深夜装修的噪声、开车造成的拥堵。
  \item 正外部性：接种疫苗（也保护了身边的人）、基础科学研究、养蜂人的蜜蜂给邻近果园授粉。
\end{itemize}

\paragraph{为什么市场会生产过多}
假设奶茶店每做一杯奶茶，会给周围居民带来 6 元的垃圾和噪声成本，但店家不用为此付钱。店家只按自己的成本（供给曲线 $P=Q$）决定产量，
可对整个社会来说，每一杯的真实成本是 $Q+6$，这条线叫\textbf{社会成本曲线}（图~\ref{fig:es-extern}）。
\begin{itemize}
  \item 市场均衡：需求与店家成本相交，$Q=30$。
  \item 对社会最好的产量：需求与社会成本相交，$45-\frac12Q=Q+6$，$Q=26$。
\end{itemize}
第 27 到第 30 杯，买家愿意付的钱低于它们的社会成本，却还是被生产出来了。多生产造成的损失是图中的绿色三角形，面积
$\frac12\times(36-30)\times(30-26)=12$。

\begin{figure}[!htbp]
\centering
\begin{tikzpicture}[x=.12cm, y=.1cm]
  \fill[dwlfill] (26,32) -- (30,36) -- (30,30) -- cycle;
  \econaxes{72}{52}{数量 $Q$}{价格 $P$}
  \draw[dem] (0,45) -- (66,12) node[below, slab, mjAccent] {需求 $D$};
  \draw[sup] (0,0) -- (47,47) node[right, slab, mjSecond] {店家自己的成本};
  \draw[sup2] (0,6) -- (41,47) node[above, slab, mjSecond] {社会成本（多 6 元）};
  \draw[guide] (26,32) -- (26,0) node[below, slab] {26};
  \draw[guide] (30,36) -- (30,0) node[below, slab] {30};
  \node[eqpt] at (30,30) {}; \node[eqpt] at (26,32) {};
  \node[slab, anchor=west] (me) at (40,14) {市场均衡}; \draw[mjInk!60, line width=.3pt] (me.west) -- (30.6,29.4);
  \node[slab, anchor=east] (so) at (14,49) {社会最优}; \draw[mjInk!60, line width=.3pt] (so.east) -- (25.6,32.6);
  \node[slab, anchor=west, mjPine] at (50,31) {多生产的损失 12};
  \draw[mjPine, line width=.4pt] (50,31) -- (29,32.5);
\end{tikzpicture}
\caption{负外部性：市场的产量（30）超过对社会最好的产量（26）}
\label{fig:es-extern}
\end{figure}

\paragraph{庇古税：让成本回到价格里}
如果政府对每杯征收 6 元税，店家的成本曲线正好上移到社会成本曲线，市场自己就会停在 $Q=26$。
请注意：这和第 3 讲的数字一模一样，同样是每杯 6 元的税、成交量从 30 降到 26。区别在于：
\begin{itemize}
  \item 第 3 讲的市场\textbf{没有}外部性，原来的 30 杯本来就是最好的，税造成了 12 的损失；
  \item 这里的市场\textbf{有}外部性，原来的 30 杯太多了，税\textbf{消除}了 12 的损失。
\end{itemize}
同一种政策是好是坏，要看市场原来有没有问题。按外部成本征收的税叫\textbf{庇古税}，碳税、排污费都是这个思路。

\paragraph{其他解决办法}
\begin{itemize}
  \item \textbf{排污许可交易}：政府定下排放总量，把许可证发给企业，允许买卖。减排成本低的企业多减排、卖出许可证，总量控制住了，减排的总成本也最低。碳排放权交易就是这样设计的。
  \item \textbf{明确产权、自行协商}：如果法律明确规定居民有权不受噪声干扰，店家要么赔偿，要么减产；只要协商的成本很低，双方自己就能谈出有效率的结果（经济学家科斯的思想）。
  \item \textbf{补贴正外部性}：疫苗、基础研究的社会好处大于个人好处，个人自愿付费时会做得太少，所以常由政府补贴或免费提供。
\end{itemize}

\section{公共物品与公地悲剧（了解）}

物品可以按两个问题分类：能不能阻止不付钱的人使用（\textbf{排他性}）？一个人用了，会不会减少别人能用的（\textbf{竞争性}）？

\begin{figure}[!htbp]
\centering
\begin{tikzpicture}[x=4.2cm, y=1.6cm]
  \fill[mjSecond!8] (0,0) rectangle (1,1);
  \fill[mjAccent!8] (1,0) rectangle (2,1);
  \draw[mj thin] (0,0) rectangle (2,2);
  \draw[mj thin] (1,0) -- (1,2) (0,1) -- (2,1);
  \node[lab, above] at (.5,2) {有竞争性}; \node[lab, above] at (1.5,2) {无竞争性};
  \node[lab, left] at (0,1.5) {有排他性}; \node[lab, left] at (0,.5) {无排他性};
  \node[align=center, font=\small] at (.5,1.5) {\textbf{私人物品}\\ 奶茶、衣服};
  \node[align=center, font=\small] at (1.5,1.5) {\textbf{俱乐部物品}\\ 视频会员、收费公路};
  \node[align=center, font=\small] at (.5,.5) {\textbf{公共资源}\\ 公海的鱼、草场};
  \node[align=center, font=\small] at (1.5,.5) {\textbf{公共物品}\\ 国防、路灯、天气预报};
\end{tikzpicture}
\caption{四类物品。下面一行没有排他性，市场难以自己解决。}
\label{fig:es-goods}
\end{figure}

\begin{itemize}
  \item \textbf{公共物品}：不付钱也能用，人人都想“搭便车”，私人不愿提供，所以通常由政府用税收提供。
  \item \textbf{公共资源}：谁都能用，但用一点就少一点。每个牧民都想多放一只羊，好处归自己，草场退化的代价大家分担，最后草场被过度放牧，这叫\textbf{公地悲剧}。
        它和第 5 讲的“囚徒困境”是同一种结构：每个人的最优选择，合起来对大家都不好。
\end{itemize}

\begin{oneline}
限价制造短缺或过剩；外部性让价格“说错了话”，好的税收能把它纠正过来。
\end{oneline}

\begin{exercises}[练一练]
  \item 需求 $Q_d=120-4P$，供给 $Q_s=2P$。(1) 规定最低价 25 元，会过剩多少？(2) 若政府承诺按 25 元收购全部剩余，要花多少钱？
  \item 某工厂每生产 1 吨产品，会给下游渔民造成 50 元损失。(1) 不加干预时，为什么产量会过多？(2) 为什么每吨征税 50 元能纠正这一点？
        (3) 若法律规定渔民有权要求工厂赔偿全部损失，结果会怎样？
  \item 接种流感疫苗有正外部性。如果完全由个人自愿付费接种，接种人数会比对社会最好的人数多还是少？为什么？
\end{exercises}

% =============================================================================
\lecture[number=5, short=博弈论入门]{博弈论入门：当你的最优选择取决于别人}
% =============================================================================

% 2×2 收益表：\gm{行1}{行2}{列1}{列2}{左上}{右上}{左下}{右下}
\newcommand{\gm}[8]{%
  \draw[mj thin] (0,0) rectangle (2,2); \draw[mj thin] (1,0) -- (1,2) (0,1) -- (2,1);
  \node[slab, left, mjAccent] at (0,1.5) {#1}; \node[slab, left, mjAccent] at (0,.5) {#2};
  \node[slab, above, mjSecond] at (.5,2) {#3}; \node[slab, above, mjSecond] at (1.5,2) {#4};
  \node[font=\small] at (.5,1.5) {#5}; \node[font=\small] at (1.5,1.5) {#6};
  \node[font=\small] at (.5,.5) {#7}; \node[font=\small] at (1.5,.5) {#8};}
\tikzset{
  devA/.style={-{Stealth[length=2mm,width=1.5mm]}, draw=mjAccent, line width=.9pt},
  devB/.style={-{Stealth[length=2mm,width=1.5mm]}, draw=mjSecond, line width=.9pt},
  nebox/.style={draw=mjPine, line width=1.4pt, rounded corners=2pt},
}

\begin{keypoints}
\begin{itemize}
  \item 一个博弈由三样东西组成：参与人、各自可选的策略、每种策略组合下各自的收益。
  \item \textbf{纳什均衡}：给定别人的选择，谁也不想单独改变自己的选择。用“划线法”或“偏离箭头”来找。
  \item \textbf{囚徒困境}：每个人都有不论对方怎么做都更好的选择（背叛），结果却对大家都不好。
  \item 博弈重复进行时，未来的惩罚可以维持合作：冷酷策略要求贴现因子 $\delta\ge\dfrac{T-R}{T-P}$。
  \item 先后行动的博弈用\textbf{逆向归纳}（从最后一步往回推）；不可信的威胁吓不住人。
\end{itemize}
\end{keypoints}

\section{怎么读收益表}

前几讲里，买家、卖家各自面对价格做决定，不用考虑别人会怎么做。但很多场合，你的最好选择取决于别人的选择：两家企业定价、两国谈判、班上同学要不要都去补课。
研究这类问题的工具叫\textbf{博弈论}。

经典例子是\textbf{囚徒困境}。两名嫌疑人甲、乙被分开审讯，各自可以选择“抵赖”或“坦白”：
\begin{itemize}
  \item 两人都抵赖：证据不足，各判 1 年；
  \item 两人都坦白：各判 5 年；
  \item 一人坦白、一人抵赖：坦白的人立功释放，抵赖的人判 10 年。
\end{itemize}
把“判几年”写成负数作为收益（越大越好），就得到图~\ref{fig:es-pd}。规定：\textbf{行}是甲的选择，\textbf{列}是乙的选择，每格括号里\textbf{先写甲的收益，后写乙的}。

\begin{figure}[!htbp]
\centering
\begin{tikzpicture}[x=2.6cm, y=1.2cm]
  \gm{抵赖}{坦白}{抵赖}{坦白}{\pay{-1}{-1}}{\pay{-10}{0}}{\pay{0}{-10}}{\pay{-5}{-5}}
  \node[lab, mjAccent] at (-.75,1) {甲};
  \node[lab, mjSecond] at (1,2.65) {乙};
  \draw[mjAccent, line width=.4pt] (2.15,1.6) -- (2.45,1.85) node[right, slab, mjAccent] {甲的收益（第一个数）};
  \draw[mjSecond, line width=.4pt] (2.15,1.4) -- (2.45,1.2) node[right, slab, mjSecond] {乙的收益（第二个数）};
\end{tikzpicture}
\caption{囚徒困境的收益表（收益 $=-$ 刑期）}
\label{fig:es-pd}
\end{figure}

\section{占优策略：不管你怎么做，我都这样做}

站在甲的角度想：
\begin{itemize}
  \item 如果乙抵赖：我抵赖得 $-1$，坦白得 $0$，坦白更好；
  \item 如果乙坦白：我抵赖得 $-10$，坦白得 $-5$，还是坦白更好。
\end{itemize}
无论乙怎么做，坦白都更好。这样的策略叫\textbf{占优策略}。乙的处境完全对称，也会坦白。结果两人都坦白，各判 5 年；可如果两人都抵赖，只各判 1 年。

\textbf{每个人都做了对自己最好的选择，结果却对两人都不好。}这就是囚徒困境：个人理性和集体理性发生了冲突。
现实中，同样的结构到处都是：
\begin{itemize}
  \item 两家企业打价格战：都维持高价对双方最好，但各自都想降价抢市场；
  \item 国家间的军备竞赛、企业间的广告大战；
  \item 第 4 讲的公地悲剧：每个牧民都多放羊；
  \item 考前大家都去补课：每个人都怕落后，结果相对名次没变，大家都更累。
\end{itemize}

\section{纳什均衡与划线法}

不是每个博弈都有占优策略，更一般的概念是：

\begin{definition}[纳什均衡]
一组策略，使得\textbf{在别人的策略给定时，任何一个人单独改变自己的策略都不能使自己的收益增加}。
\end{definition}

找纳什均衡的\textbf{划线法}：
\begin{enumerate}
  \item 对乙的每一个选择（每一列），找出甲收益最大的那格，在\textcolor{mjAccent}{甲的收益}下面划线；
  \item 对甲的每一个选择（每一行），找出乙收益最大的那格，在\textcolor{mjSecond}{乙的收益}下面划线；
  \item 两个数都被划线的格子，就是纳什均衡。
\end{enumerate}

图~\ref{fig:es-underline}(b) 是一个\textbf{协调博弈}：甲、乙约好周末见面，但忘了约在哪里，手机也联系不上。甲更想去图书馆，乙更想去咖啡馆，但两人都最怕走散（收益为 0）。
划线后有\textbf{两个}纳什均衡：都去图书馆，或者都去咖啡馆。一旦两人到了同一个地方，谁也不想单独换地方。

\begin{figure}[!htbp]
\centering
\begin{subfigure}{.48\textwidth}\centering
\begin{tikzpicture}[x=1.9cm, y=1.05cm]
  \gm{抵赖}{坦白}{抵赖}{坦白}{\pay{-1}{-1}}{\upayB{-10}{0}}{\upayA{0}{-10}}{\upay{-5}{-5}}
  \draw[nebox] (1.05,.07) rectangle (1.95,.93);
  \node[lab, mjAccent] at (-.75,1) {甲}; \node[lab, mjSecond] at (1,2.6) {乙};
\end{tikzpicture}
\caption{囚徒困境：唯一均衡（坦白，坦白）}
\end{subfigure}\hfill
\begin{subfigure}{.48\textwidth}\centering
\begin{tikzpicture}[x=1.9cm, y=1.05cm]
  \gm{图书馆}{咖啡馆}{图书馆}{咖啡馆}{\upay{3}{2}}{\pay{0}{0}}{\pay{0}{0}}{\upay{2}{3}}
  \draw[nebox] (.05,1.07) rectangle (.95,1.93);
  \draw[nebox] (1.05,.07) rectangle (1.95,.93);
  \node[lab, mjAccent] at (-.85,1) {甲}; \node[lab, mjSecond] at (1,2.6) {乙};
\end{tikzpicture}
\caption{协调博弈：两个均衡}
\end{subfigure}
\caption{划线法。绿框是两个数都被划线的格子，即纳什均衡。}
\label{fig:es-underline}
\end{figure}

\paragraph{另一种看法：偏离箭头}
在每个格子里问：有没有人想单独换？如果甲想换行，就画一个从这格指向另一行的\textcolor{mjAccent}{蓝箭头}；如果乙想换列，就画\textcolor{mjSecond}{红箭头}。
\textbf{没有箭头射出的格子就是纳什均衡}（图~\ref{fig:es-arrows}）。
\begin{itemize}
  \item 囚徒困境：所有箭头最后都流向（坦白，坦白）。
  \item 协调博弈：箭头流向两个角，两个均衡。
  \item 猜硬币：甲、乙同时出硬币的正面或反面，相同则甲赢 1 分，不同则乙赢 1 分。箭头绕成一圈，每一格都有人想换，\textbf{没有}这种“确定选一个”的纳什均衡。
\end{itemize}
猜硬币中，最好的办法是让对方猜不到：各以 $\frac12$ 的概率随机出正反。这叫\textbf{混合策略}，足球点球时罚球员和守门员的选择就是这样。

\begin{figure}[!htbp]
\centering
\begin{subfigure}{.32\textwidth}\centering
\begin{tikzpicture}[x=1.45cm, y=1cm, every node/.append style={font=\footnotesize}]
  \gm{抵赖}{坦白}{抵赖}{坦白}{\pay{-1}{-1}}{\pay{-10}{0}}{\pay{0}{-10}}{\pay{-5}{-5}}
  \draw[devB] (.72,1.83) -- (1.28,1.83);
  \draw[devB] (.72,.83) -- (1.28,.83);
  \draw[devA] (.1,1.3) -- (.1,.7);
  \draw[devA] (1.1,1.3) -- (1.1,.7);
  \draw[nebox] (1.04,.05) rectangle (1.96,.95);
\end{tikzpicture}
\caption{囚徒困境}
\end{subfigure}\hfill
\begin{subfigure}{.32\textwidth}\centering
\begin{tikzpicture}[x=1.45cm, y=1cm, every node/.append style={font=\footnotesize}]
  \gm{图书馆}{咖啡馆}{图}{咖}{\pay{3}{2}}{\pay{0}{0}}{\pay{0}{0}}{\pay{2}{3}}
  \draw[devB] (1.28,1.83) -- (.72,1.83);
  \draw[devB] (.72,.83) -- (1.28,.83);
  \draw[devA] (.1,.7) -- (.1,1.3);
  \draw[devA] (1.1,1.3) -- (1.1,.7);
  \draw[nebox] (.04,1.05) rectangle (.96,1.95);
  \draw[nebox] (1.04,.05) rectangle (1.96,.95);
\end{tikzpicture}
\caption{协调博弈}
\end{subfigure}\hfill
\begin{subfigure}{.32\textwidth}\centering
\begin{tikzpicture}[x=1.45cm, y=1cm, every node/.append style={font=\footnotesize}]
  \gm{正}{反}{正}{反}{\pay{1}{-1}}{\pay{-1}{1}}{\pay{-1}{1}}{\pay{1}{-1}}
  \draw[devB] (.72,1.83) -- (1.28,1.83);
  \draw[devB] (1.28,.83) -- (.72,.83);
  \draw[devA] (.1,.7) -- (.1,1.3);
  \draw[devA] (1.1,1.3) -- (1.1,.7);
\end{tikzpicture}
\caption{猜硬币}
\end{subfigure}
\caption{偏离箭头：蓝箭头表示甲想换行，红箭头表示乙想换列。没有箭头射出的格子（绿框）是纳什均衡。}
\label{fig:es-arrows}
\end{figure}

\begin{pitfall}
\begin{itemize}
  \item 纳什均衡\textbf{不一定}是对大家最好的结果（囚徒困境），也\textbf{不一定}唯一（协调博弈）。
  \item “占优策略”是“不论对方怎么做都更好”的策略，不是“结果最好”的策略。
  \item 检查纳什均衡时，每次只让\textbf{一个人}改变，另一个人的选择保持不动。
\end{itemize}
\end{pitfall}

\section{重复博弈：未来的影子}

如果两家企业不是只打一次交道，而是年年都在同一个市场里竞争呢？今天背叛，明天对方可能报复。未来的报复能不能让双方今天就合作？

\paragraph{收益的记号}
把囚徒困境的四种收益记为：
\begin{center}
\small
\begin{tabular}{llc}
  \toprule
  记号 & 含义 & 大小关系 \\
  \midrule
  $T$（诱惑） & 我背叛、对方合作时我的收益 & 最大 \\
  $R$（奖励） & 双方合作时我的收益 & 第二 \\
  $P$（惩罚） & 双方背叛时我的收益 & 第三 \\
  $S$（吃亏） & 我合作、对方背叛时我的收益 & 最小 \\
  \bottomrule
\end{tabular}
\end{center}
下面用 $T=4$，$R=3$，$P=1$，$S=0$ 作例子。

\paragraph{贴现因子}
明天的 1 分收益，今天看来只值 $\delta$ 分（$0<\delta<1$），后天的只值 $\delta^2$ 分，依此类推。$\delta$ 叫\textbf{贴现因子}，有两种理解：
\begin{itemize}
  \item \textbf{耐心}：钱可以生利息，明天的钱不如今天的钱值钱。年利率为 $r$ 时，$\delta=\frac1{1+r}$（见第 6 讲）；
  \item \textbf{博弈继续的概率}：每一轮结束后，有 $\delta$ 的概率还有下一轮。
\end{itemize}
$\delta$ 越大，越看重未来。

\paragraph{冷酷策略}
“先合作；只要对方背叛过一次，从此永远背叛。”双方都用这个策略时，会一直合作下去。问题是：会不会有人想偷偷背叛一次？

\begin{figure}[!htbp]
\centering
\begin{tikzpicture}[x=1.25cm, y=.42cm]
  % 一直合作
  \node[lab, anchor=east] at (-.2,6.4) {一直合作};
  \foreach \t in {1,...,6} \fill[mjAccent!55] (\t-.35,5) rectangle (\t+.35,8);
  \foreach \t in {1,...,6} \node[slab] at (\t,8.6) {3};
  % 背叛一次
  \node[lab, anchor=east] at (-.2,1.4) {第 1 期背叛};
  \fill[mjSecond!65] (1-.35,0) rectangle (1+.35,4);
  \node[slab] at (1,4.6) {4};
  \foreach \t in {2,...,6} \fill[mjSecond!30] (\t-.35,0) rectangle (\t+.35,1);
  \foreach \t in {2,...,6} \node[slab] at (\t,1.6) {1};
  \draw[mj thin] (.4,0) -- (6.9,0); \draw[mj thin] (.4,5) -- (6.9,5);
  \node[slab] at (7.45,6.4) {$\cdots$}; \node[slab] at (7.45,1) {$\cdots$};
  \foreach \t/\w in {1/1, 2/\delta, 3/\delta^2, 4/\delta^3, 5/\delta^4, 6/\delta^5}
    \node[slab] at (\t,-1) {$\w$};
  \node[slab, anchor=east] at (-.2,-1) {折算系数};
  \node[slab, mjSecond] at (1,-2.5) {多得 1};
  \foreach \t in {2,...,6} \node[slab, mjAccent] at (\t,-2.5) {少得 2};
  \node[slab, anchor=east] at (-.2,-2.5) {背叛比合作};
\end{tikzpicture}
\caption{背叛一次：当期多得 $T-R=1$，此后每期少得 $R-P=2$}
\label{fig:es-repeat}
\end{figure}

图~\ref{fig:es-repeat} 比较了两条路：
\begin{itemize}
  \item 一直合作：$R+\delta R+\delta^2R+\cdots=\dfrac{R}{1-\delta}$；
  \item 第 1 期背叛：当期得 $T$，此后对方永远背叛，自己最好也只能每期得 $P$，合计
        $T+\delta P+\delta^2P+\cdots=T+\dfrac{\delta P}{1-\delta}$。
\end{itemize}
不背叛的条件是
\[
  \frac{R}{1-\delta}\ \ge\ T+\frac{\delta P}{1-\delta}
  \quad\Longleftrightarrow\quad R\ge(1-\delta)T+\delta P
  \quad\Longleftrightarrow\quad \boxed{\ \delta\ \ge\ \frac{T-R}{T-P}\ } .
\]
例子中 $\delta\ge\frac{4-3}{4-1}=\frac13$。

\textbf{怎么理解}：背叛的好处是一次性的（$T-R$），代价却是以后每一期都少得 $R-P$。只要足够看重未来，未来损失的总和就会超过眼前的好处。
从公式还能读出：诱惑 $T$ 越大，要求的 $\delta$ 越大；惩罚 $P$ 越重（越小），合作越容易维持。

\paragraph{几点补充}
\begin{itemize}
  \item \textbf{有限次重复不行}：如果大家都知道博弈恰好进行 10 轮，第 10 轮就是一次性博弈，双方都会背叛；那么第 9 轮的合作也换不来未来的回报……一路倒推，每一轮都背叛。
        合作需要“没有确定的终点”。
  \item \textbf{以牙还牙}：“先合作，之后每一轮都照搬对方上一轮的做法”。在政治学家阿克塞尔罗德 1980 年前后组织的计算机博弈比赛中，这个简单策略总成绩最好。
        它的特点是：善良（不先背叛）、会报复、肯原谅、容易被对方看懂。
  \item \textbf{现实例子}：少数几家企业长期共处一个市场，容易形成“心照不宣”的高价，所以反垄断机构特别留意这种默契；
        老店重视回头客，商誉值钱；人与人之间的信任，也常常建立在“以后还要打交道”之上。
\end{itemize}

\section{先后行动：博弈树与逆向归纳}

有些博弈是一方先动、另一方看到后再动。例如：一家新企业考虑是否进入某个市场，原来独占市场的企业（在位者）看到后，再决定是打价格战还是容忍。
这样的博弈用\textbf{博弈树}表示（图~\ref{fig:es-tree}），叶子上的括号先写进入者的收益，后写在位者的。

\begin{figure}[!htbp]
\centering
\begin{tikzpicture}[
    dec/.style={circle, draw=mjInk, fill=white, line width=.6pt, inner sep=2.2pt},
    leaf/.style={font=\small},
    br/.style={draw=mjInk!60, line width=.7pt},
    pick/.style={draw=mjPine, line width=2.2pt}]
  \node[dec, label={[lab]above:进入者}] (A) at (0,0) {};
  \node[dec, label={[lab]right:在位者}] (B) at (2.6,-1.6) {};
  \node[leaf] (L0) at (-2.6,-1.9) {\pay{0}{10}};
  \node[leaf] (L1) at (0.8,-3.6) {\pay{-2}{2}};
  \node[leaf] (L2) at (4.4,-3.6) {\pay{3}{5}};
  \draw[br] (A) -- node[above left, slab] {不进入} (L0);
  \draw[pick] (A) -- node[above right, slab] {进入} (B);
  \draw[br] (B) -- node[left, slab] {打价格战} (L1);
  \draw[pick] (B) -- node[right, slab] {容忍} (L2);
  \node[slab, anchor=west, mjPine, align=left] at (5.2,-1.2) {第一步：在位者比较 5 和 2，选“容忍”\\ 第二步：进入者知道这一点，比较 3 和 0，选“进入”};
\end{tikzpicture}
\caption{进入博弈。粗绿线是逆向归纳选出的路径。}
\label{fig:es-tree}
\end{figure}

\textbf{逆向归纳}：从最后一步往回推。
\begin{enumerate}
  \item 如果新企业已经进入了，在位者打价格战得 2，容忍得 5，它会容忍；
  \item 新企业预见到这一点：进入得 3，不进入得 0，所以进入。
\end{enumerate}
结果是（进入，容忍）。在位者事前可以放话“你敢进来我就打价格战”，但这是\textbf{不可信的威胁}：真到那一步，打价格战对它自己也不利。

怎样让威胁变得可信？办法是\textbf{事先把自己的退路堵死}。比如在位者预先建好大量产能，一旦有人进入，打价格战反而对它更有利；
“破釜沉舟”也是这个道理：把船凿沉，让后退变得不可能，“死战”的决心就可信了。

\begin{oneline}
找均衡就是找“谁都不想单独改变”的格子；重复让合作有了未来的回报，逆向归纳让空口威胁失效。
\end{oneline}

\begin{exercises}[练一练]
  \item （猎鹿博弈）两名猎人可以合作猎鹿，也可以各自去抓兔子。收益为：都猎鹿 $(4,4)$；甲猎鹿、乙抓兔 $(0,3)$；甲抓兔、乙猎鹿 $(3,0)$；都抓兔 $(3,3)$。
        画出收益表，找出所有纳什均衡。哪个均衡对双方更好？哪个更“保险”？
  \item 囚徒困境中 $T=6$，$R=4$，$P=2$，$S=0$。(1) 冷酷策略维持合作需要 $\delta$ 满足什么条件？
        (2) 若双方不考虑利息，但每一轮结束后有 80\% 的概率进行下一轮，合作能否维持？
  \item 在图~\ref{fig:es-tree} 的进入博弈中，若在位者事先扩建了产能，使“打价格战”时它的收益从 2 变为 6，其他收益不变。用逆向归纳求新的结果，并说明扩建产能起了什么作用。
\end{exercises}

% =============================================================================
\lecture[number=6, short=钱的时间价值]{钱的时间价值：利息、贴现与房贷}
% =============================================================================

\begin{keypoints}
\begin{itemize}
  \item 今天的 1 元比明年的 1 元值钱，因为今天的钱可以生利息。年利率为 $r$ 时，$n$ 年后的 $F$ 元，今天只值 $\dfrac{F}{(1+r)^n}$，这叫\textbf{现值}。
  \item 复利是指数增长。“72 法则”：翻倍所需年数约为 $72\div$（年利率的百分数）。
  \item 每期付 $A$、共付 $n$ 期、每期利率为 $i$，现值是一个等比数列的和：$A\cdot\dfrac{1-(1+i)^{-n}}{i}$；一直付下去，现值是 $\dfrac Ai$。
  \item 比较不同的付款方式，要把所有现金流折算到同一时刻。“总利息少”不等于“更划算”。
\end{itemize}
\end{keypoints}

\section{单利与复利}

把 100 元存起来，年利率 5\%：
\begin{itemize}
  \item \textbf{单利}：每年只按本金 100 元算利息，每年 5 元，30 年后共 $100+5\times30=250$ 元；
  \item \textbf{复利}：每年的利息也加进本金一起生息，$n$ 年后是 $100\times1.05^n$，30 年后约 432 元。
\end{itemize}
单利是直线，复利是指数曲线。开始时两者差不多，时间越长差得越多（图~\ref{fig:es-compound}）。

\paragraph{72 法则}
复利下，钱翻一倍需要多少年？解 $1.05^n=2$，得 $n=\frac{\ln2}{\ln1.05}\approx14.2$ 年。一个好用的近似是
\[
  \text{翻倍年数}\approx\frac{72}{\text{年利率的百分数}},
\]
5\% 时 $72\div5=14.4$ 年，误差很小。这个近似来自 $\ln2\approx0.693$ 以及 $\ln(1+r)\approx r$；用 72 而不用 69，是因为 72 的因数多、好心算，而且在常见的利率下更准一些。

\begin{figure}[!htbp]
\centering
\begin{tikzpicture}
\begin{axis}[mojing, width=9.5cm, height=6cm, xmin=0, xmax=32, ymin=0, ymax=470,
    xlabel={年}, ylabel={元}, xtick={0,10,14.2,20,30}, xticklabels={0,10,14.2,20,30},
    ytick={100,200,250,432}, legend pos=north west]
  \addplot[mjSecond, domain=0:30, samples=40] {100*1.05^x};
  \addplot[mjAccent, domain=0:30, samples=2] {100*(1+0.05*x)};
  \legend{复利,单利}
  \draw[guide] (axis cs:0,200) -- (axis cs:14.2,200) -- (axis cs:14.2,0);
  \node[eqpt] at (axis cs:14.2,200) {};
  \node[slab, anchor=south east] at (axis cs:14,205) {翻倍};
  \node[eqpt] at (axis cs:30,432.2) {}; \node[eqpt] at (axis cs:30,250) {};
\end{axis}
\end{tikzpicture}
\caption{100 元按年利率 5\% 存 30 年：单利 250 元，复利约 432 元}
\label{fig:es-compound}
\end{figure}

\section{现值：把未来的钱搬到今天}

反过来问：3 年后的 1000 元，今天值多少钱？如果今天存 $x$ 元，3 年后正好变成 1000 元，那么 $x(1.05)^3=1000$，$x\approx863.84$。
所以在 5\% 的利率下，“3 年后的 1000 元”和“今天的 863.84 元”是等价的。

\begin{definition}[现值]
利率为 $r$ 时，$n$ 期后的 $F$ 元的\textbf{现值}是 $\dfrac{F}{(1+r)^n}$。把未来的钱换算成现值的过程叫\textbf{贴现}。
\end{definition}

现值的意义在于：\textbf{不同时间点的钱不能直接相加或比较}，换算到同一时刻以后才能比较。第 5 讲重复博弈里的贴现因子 $\delta$，就是 $\frac1{1+r}$：明天的 1 元今天值 $\delta$ 元。

\section{年金：一串相等的付款}

很多付款是一串等额的钱：每月房贷、每年的奖学金、退休金。每期期末付 $A$，共 $n$ 期，每期利率 $i$。把每一笔都贴现到今天（图~\ref{fig:es-timeline}(b)），
记 $v=\frac1{1+i}$：
\[
  \text{现值}=Av+Av^2+\cdots+Av^n=Av\cdot\frac{1-v^n}{1-v}=A\cdot\frac{1-(1+i)^{-n}}{i}.
\]
最后一步用到 $\dfrac{v}{1-v}=\dfrac1i$。

\paragraph{永续年金}
如果一直付下去（$n\to\infty$），$(1+i)^{-n}\to0$，现值是 $\dfrac Ai$。例如年利率 5\% 时，“每年永远给你 100 元”今天只值 $100\div0.05=2000$ 元。
道理很直观：今天存 2000 元，每年的利息正好是 100 元，本金永远不动。

\begin{figure}[!htbp]
\centering
\begin{tikzpicture}[x=1.9cm, y=1cm]
  % (a) 单笔
  \node[lab, anchor=west] at (-.6,4.3) {(a) 单笔：3 年后的 1000 元};
  \draw[mj axis] (0,2.2) -- (4.3,2.2) node[right, slab] {年};
  \foreach \k in {0,...,3} { \draw (\k,2.13) -- (\k,2.27); \node[slab, below] at (\k,2.13) {\k}; }
  \draw[mjSecond, line width=1.2pt, -{Stealth}] (3,2.2) -- (3,3.1) node[above, slab, mjSecond] {1000};
  \draw[mjAccent, line width=1.2pt, -{Stealth}] (0,2.2) -- (0,3.1) node[above, slab, mjAccent] {863.84};
  \draw[mjAccent, line width=.8pt, -{Stealth}] (2.9,2.65) to[bend right=22] node[above, slab, mjAccent] {$\div1.05^3$} (.1,2.65);
  % (b) 年金
  \node[lab, anchor=west] at (-.6,.9) {(b) 年金：每期末付 $A$，共 $n$ 期};
  \draw[mj axis] (0,-.6) -- (5.4,-.6) node[right, slab] {期};
  \foreach \k/\l in {0/0,1/1,2/2,3/3,5/n} { \draw (\k,-.67) -- (\k,-.53); \node[slab, below] at (\k,-.67) {$\l$}; }
  \node[slab] at (4,-.9) {$\cdots$};
  \foreach \k in {1,2,3,5} \draw[mjSecond, line width=1.2pt, -{Stealth}] (\k,-.6) -- (\k,.05) node[above, slab, mjSecond] {$A$};
  \foreach \k/\b in {1/40, 2/30, 3/26, 5/22}
    \draw[mjAccent, line width=.6pt, -{Stealth}] (\k-.03,-1.1) to[bend left=\b] (.03,-1.1);
  \node[slab, mjAccent] at (2.5,-2.75) {现值 $=Av+Av^2+Av^3+\cdots+Av^n$，其中 $v=\frac1{1+i}$};
\end{tikzpicture}
\caption{贴现：把未来每一笔钱都折算到 0 时刻再相加}
\label{fig:es-timeline}
\end{figure}

\section{两种还款方式}

借 10000 元，分 12 个月还清，月利率 1\%。银行通常提供两种方式。

\paragraph{等额本息}
每月还款额相同，记为 $A$。还款的现值之和应等于借到的钱：
\[
  10000=A\cdot\frac{1-1.01^{-12}}{0.01}\quad\Longrightarrow\quad A\approx888.49\ \text{元}.
\]
12 个月共还 $12A\approx10661.85$ 元，总利息约 661.85 元。每月的 888.49 元里，先付当月利息（欠款 $\times1\%$），剩下的才用来还本金。
开始时欠得多，利息占的比例大；越往后，利息越少，还的本金越多。

\paragraph{等额本金}
每月还相同的本金 $10000\div12\approx833.33$ 元，另付当月剩余欠款的利息。第 1 个月欠 10000 元，利息 100 元，共还 933.33 元；
以后欠款每月减少 833.33 元，利息每月减少 8.33 元；第 12 个月只还 841.67 元。总利息是
\[
  100+91.67+83.33+\cdots+8.33=0.01\times\frac{10000\times(12+11+\cdots+1)}{12}=650\ \text{元}.
\]

\begin{figure}[!htbp]
\centering
\begin{tikzpicture}
\begin{axis}[mojing box, ybar stacked, bar width=7pt, width=7.6cm, height=5.6cm,
    ymin=0, ymax=1000, xmin=.3, xmax=12.7, xtick={1,3,6,9,12}, xlabel={月}, ylabel={元},
    title={\small 等额本息}, legend style={at={(.5,1.02)}, anchor=south, legend columns=2, draw=none, yshift=10pt}, enlarge x limits=false]
  \addplot[fill=mjAccent!60, draw=none] coordinates {(1,788.49) (2,796.37) (3,804.34) (4,812.38) (5,820.50) (6,828.71) (7,837.00) (8,845.37) (9,853.82) (10,862.36) (11,870.98) (12,879.69)};
  \addplot[fill=mjOchre!70, draw=none] coordinates {(1,100.00) (2,92.12) (3,84.15) (4,76.11) (5,67.98) (6,59.78) (7,51.49) (8,43.12) (9,34.67) (10,26.13) (11,17.51) (12,8.80)};
\end{axis}
\end{tikzpicture}\hfill
\begin{tikzpicture}
\begin{axis}[mojing box, ybar stacked, bar width=7pt, width=7.6cm, height=5.6cm,
    ymin=0, ymax=1000, xmin=.3, xmax=12.7, xtick={1,3,6,9,12}, xlabel={月}, ylabel={元},
    title={\small 等额本金}, enlarge x limits=false]
  \addplot[fill=mjAccent!60, draw=none] coordinates {(1,833.33) (2,833.33) (3,833.33) (4,833.33) (5,833.33) (6,833.33) (7,833.33) (8,833.33) (9,833.33) (10,833.33) (11,833.33) (12,833.33)};
  \addplot[fill=mjOchre!70, draw=none] coordinates {(1,100.00) (2,91.67) (3,83.33) (4,75.00) (5,66.67) (6,58.33) (7,50.00) (8,41.67) (9,33.33) (10,25.00) (11,16.67) (12,8.33)};
\end{axis}
\end{tikzpicture}
\caption{每月还款的构成：\textcolor{mjAccent}{\textbf{蓝色}}是本金，\textcolor{mjOchre}{\textbf{棕色}}是利息。借 10000 元，12 个月，月利率 1\%。}
\label{fig:es-loan}
\end{figure}

\section{总利息少，就一定更划算吗}

等额本金总利息 650 元，比等额本息的 661.85 元少。原因在于：\textbf{利息是借钱的“租金”}，每个月的利息 $=$ 月初欠款 $\times$ 月利率。
等额本金每月还的本金多、还得早，平均欠款少，租金自然少。两种方式每月的月初欠款加起来分别是 65000 元和约 66185 元，乘以 1\% 正好是两种方式的总利息。

但这并不说明等额本金更划算：
\begin{itemize}
  \item 按 1\% 的月利率贴现，\textbf{两种方式还款的现值都恰好是 10000 元}。站在银行的角度，两者完全等价；
  \item 等额本息让你把钱留在手里更久。如果你能把这些钱投到月收益高于 1\% 的地方，等额本息反而更好；如果你手里的钱只能存银行拿更低的利息，早点还更好；
  \item 等额本金第 1 个月要比等额本息多还约 45 元，前期的还款压力更大。刚工作、收入逐年增加的人，常常更适合等额本息。
\end{itemize}
所以，哪种方式更好，取决于借款人的收入节奏和其他投资机会，不能只看总利息。

\section{识破“免息分期”的真实利率}

买一件 12000 元的商品，分 12 期，每期还本金 1000 元，另收“手续费” $12000\times0.6\%=72$ 元。一年手续费共 864 元，看起来“年费率” $7.2\%$。

可是你欠的钱在逐月减少：最后一个月只欠 1000 元，却仍按 12000 元交手续费（图~\ref{fig:es-fee}）。真实的月利率 $i$ 应该满足“还款的现值等于借到的钱”：
\[
  12000=1072\cdot\frac{1-(1+i)^{-12}}{i},
\]
数值求解得 $i\approx1.09\%$，按 12 个月简单相加折合年利率约 $13\%$（按复利算约 $13.8\%$），接近名义费率的两倍。

\begin{figure}[!htbp]
\centering
\begin{tikzpicture}[x=.85cm, y=.03cm]
  \fill[mjSecond!10] (0,0) rectangle (12,120);
  \foreach \k in {1,...,12} {
    \pgfmathsetmacro{\h}{120-10*(\k-1)}
    \fill[mjAccent!45] (\k-1+.08,0) rectangle (\k-.08,\h);
  }
  \draw[mjSecond, line width=1pt] (0,120) -- (12,120) node[right, slab, mjSecond, align=left] {手续费按\\12000 元算};
  \draw[mj axis] (0,0) -- (12.8,0) node[right, slab] {月};
  \draw[mj axis] (0,0) -- (0,140) node[above, slab] {欠款（百元）};
  \foreach \k in {1,6,12} \node[slab, below] at (\k-.5,0) {\k};
  \node[slab, left] at (0,120) {120};
  \node[slab, left] at (0,10) {10};
  \node[slab, mjAccent!80!black, anchor=west] (ow) at (12.2,35) {实际欠款};
  \draw[mjAccent!80!black, line width=.3pt] (ow.west) -- (11.5,6);
  \node[slab, mjSecond!80!black] at (8.6,95) {已还掉，却仍被收手续费};
\end{tikzpicture}
\caption{分期手续费的陷阱：蓝色是每月实际欠的钱，红色区域是被收了手续费、你却早已还掉的部分}
\label{fig:es-fee}
\end{figure}

\paragraph{名义利率与实际利率}
还有一种“实际”：扣除物价上涨。存款年利率 3\%，物价一年涨 2\%，存款的购买力一年只增加约 1\%。经济学里常说的“实际利率”$\approx$ 名义利率 $-$ 通货膨胀率。

\begin{pitfall}
\begin{itemize}
  \item 不同时间的钱不能直接相加比较，先贴现到同一时刻。
  \item 年金公式里的 $i$ 和 $n$ 必须对应：按月还款就用月利率和月数，不能用年利率配月数。
  \item “总利息少”不等于“更划算”，要按同一个利率比较现值，再考虑现金流和其他投资机会。
\end{itemize}
\end{pitfall}

\begin{oneline}
钱有时间价值：一切比较都先折算到同一时刻；利息多少取决于你欠了多少钱、欠了多久。
\end{oneline}

\begin{exercises}[练一练]
  \item 某奖学金有两种领取方式：现在领 8500 元，或 4 年后一次领 10000 元。年利率 4\% 时，哪种更值钱？
  \item 用等比数列求和公式推导年金现值公式 $A\cdot\dfrac{1-(1+i)^{-n}}{i}$，并说明 $n\to\infty$ 时的极限。
  \item 贷款 12000 元，12 个月，月利率 1\%，按等额本金还款。求总利息和第 1 个月的还款额。
  \item 年利率 3\% 时，用 72 法则估计钱翻倍需要多少年，再用计算器求精确值。
\end{exercises}

% =============================================================================
\lecture[number=7, short=比较优势与贸易]{比较优势：为什么“样样都强”也要分工}
% =============================================================================

\begin{keypoints}
\begin{itemize}
  \item \textbf{机会成本}：为了多得到一样东西，必须放弃的另一样东西的数量。
  \item \textbf{绝对优势}看谁效率高，\textbf{比较优势}看谁的机会成本低。分工应当按比较优势，而不是按绝对优势。
  \item 只要两国的机会成本不同，就存在双方都获益的交换比例，它落在两国的机会成本之间。
  \item 贸易使国家整体获益，但国内会有受损的人。“总收益大于总损失”不等于“每个人都获益”。
\end{itemize}
\end{keypoints}

\section{机会成本}

你用一个晚上做数学题，就不能用这个晚上背单词。这个晚上的“成本”，不是花了多少钱，而是\textbf{放弃了什么}。这叫\textbf{机会成本}。

设甲、乙两国都只生产布和粮，劳动是唯一的投入。每单位劳动能生产：
\begin{center}
\small
\begin{tabular}{lcc|cc}
  \toprule
  & 布 & 粮 & 1 单位粮的机会成本 & 1 单位布的机会成本 \\
  \midrule
  甲国 & 4 & 2 & 2 单位布 & $\frac12$ 单位粮 \\
  乙国 & 1 & 1 & 1 单位布 & 1 单位粮 \\
  \bottomrule
\end{tabular}
\end{center}
\begin{itemize}
  \item 甲国两样都比乙国做得快，在两种产品上都有\textbf{绝对优势}。
  \item 甲国多生产 1 单位粮，要少生产 2 单位布；乙国只要少生产 1 单位布。所以乙国生产粮的机会成本更低，乙国在\textbf{粮}上有比较优势。
  \item 反过来，甲国生产布的机会成本（$\frac12$ 单位粮）低于乙国（1 单位粮），甲国在\textbf{布}上有比较优势。
\end{itemize}
注意两国的机会成本互为倒数，所以\textbf{一个国家不可能在两种产品上都有比较优势}（除非两国机会成本恰好相同，那样就没有分工的好处）。

\section{生产可能性边界}

设甲国有 50 单位劳动，乙国有 60 单位劳动。甲国全部生产布能得 200，全部生产粮能得 100，也可以按比例兼顾，所有可能的组合构成一条线段，
叫\textbf{生产可能性边界}（图~\ref{fig:es-ppf}）。线段的斜率就是机会成本：甲国沿边界每多要 1 单位粮，要少 2 单位布。

没有贸易时，一国只能消费自己生产的东西，消费点只能在边界上或边界内。假设两国各自要 30 单位粮：甲国用 15 单位劳动种粮，剩 35 单位劳动织布 140；乙国用 30 单位劳动种粮，剩 30 单位劳动织布 30。
两国合计：布 170，粮 60。

\begin{figure}[!htbp]
\centering
\begin{subfigure}{.48\textwidth}\centering
\begin{tikzpicture}[x=.036cm, y=.021cm]
  \fill[mjAccent!8] (0,0) -- (0,200) -- (100,0) -- cycle;
  \econaxes{150}{230}{粮}{布}
  \draw[mjAccent, line width=1.1pt] (0,200) -- (100,0);
  \node[slab, left] at (0,200) {200}; \node[slab, below] at (100,0) {100};
  \node[eqpt] at (30,140) {};
  \draw[guide] (0,140) node[left, slab] {140} -- (30,140) -- (30,0) node[below, slab] {30};
  \node[slab, anchor=south west] at (31,141) {自给时的消费};
  \node[slab, mjAccent] at (62,120) {斜率 $-2$};
\end{tikzpicture}
\caption{甲国（50 单位劳动）}
\end{subfigure}\hfill
\begin{subfigure}{.48\textwidth}\centering
\begin{tikzpicture}[x=.064cm, y=.07cm]
  \fill[mjSecond!8] (0,0) -- (0,60) -- (60,0) -- cycle;
  \econaxes{85}{70}{粮}{布}
  \draw[mjSecond, line width=1.1pt] (0,60) -- (60,0);
  \node[slab, left] at (0,60) {60}; \node[slab, below] at (60,0) {60};
  \node[eqpt] at (30,30) {};
  \draw[guide] (0,30) node[left, slab] {30} -- (30,30) -- (30,0) node[below, slab] {30};
  \node[slab, anchor=south west] at (31,31) {自给时的消费};
  \node[slab, mjSecond] at (50,22) {斜率 $-1$};
\end{tikzpicture}
\caption{乙国（60 单位劳动）}
\end{subfigure}
\caption{两国的生产可能性边界（两图刻度不同）。斜率的绝对值是 1 单位粮的机会成本。}
\label{fig:es-ppf}
\end{figure}

\section{分工与交换}

现在按比较优势分工：甲国全部织布，得 200；乙国全部种粮，得 60。两国合计：\textbf{布 200，粮 60}。
和自给时（布 170，粮 60）相比，粮一样多，布多了 30（图~\ref{fig:es-total}）。劳动没有增加，只是换了分工方式，总产量就增加了。

\begin{figure}[!htbp]
\centering
\begin{tikzpicture}
\begin{axis}[mojing box, ybar, bar width=16pt, width=8cm, height=5.2cm, ymin=0, ymax=235,
    symbolic x coords={布,粮}, xtick=data, enlarge x limits=.45, ylabel={两国合计},
    nodes near coords, nodes near coords style={font=\footnotesize},
    legend style={at={(.98,.97)}, anchor=north east, draw=none}, area legend]
  \addplot[fill=mjGray!35, draw=none] coordinates {(布,170) (粮,60)};
  \addplot[fill=mjPine!60, draw=none] coordinates {(布,200) (粮,60)};
  \legend{各自生产,按比较优势分工}
\end{axis}
\end{tikzpicture}
\caption{分工使两国的总产量增加：布多了 30}
\label{fig:es-total}
\end{figure}

多出来的产品怎样分？设两国按“1 单位粮换 1.5 单位布”交换，乙国用 30 单位粮换甲国 45 单位布：
\begin{center}
\small
\begin{tabular}{lcc|cc}
  \toprule
  & \multicolumn{2}{c|}{自给时消费} & \multicolumn{2}{c}{分工并交换后消费} \\
  & 粮 & 布 & 粮 & 布 \\
  \midrule
  甲国 & 30 & 140 & 30 & $200-45=155$ \\
  乙国 & 30 & 30 & $60-30=30$ & 45 \\
  \bottomrule
\end{tabular}
\end{center}
两国的粮都和原来一样，布都多了 15。更重要的是，两国消费的组合都落在了\textbf{自己的生产可能性边界以外}，这是闭关自守时做不到的（图~\ref{fig:es-trade}）。

\begin{figure}[!htbp]
\centering
\begin{subfigure}{.48\textwidth}\centering
\begin{tikzpicture}[x=.036cm, y=.021cm]
  \econaxes{150}{230}{粮}{布}
  \draw[mjAccent, line width=1.1pt] (0,200) -- (100,0);
  \draw[mjPine, line width=1pt, dash pattern=on 4pt off 2pt] (0,200) -- node[pos=.62, anchor=south west, slab, mjPine] {交换线，斜率 $-1.5$} (125,12.5);
  \node[eqpt] at (30,140) {}; \node[slab, anchor=north east] at (29,139) {自给};
  \node[eqpt, fill=mjPine] at (30,155) {}; \node[slab, anchor=south west, mjPine] at (31,156) {贸易后 $(30,155)$};
  \draw[shiftarrow] (30,141) -- (30,153);
  \node[eqpt, fill=mjSecond] at (0,200) {}; \node[slab, anchor=west] at (3,205) {只生产布};
\end{tikzpicture}
\caption{甲国：生产 $(0,200)$，消费 $(30,155)$}
\end{subfigure}\hfill
\begin{subfigure}{.48\textwidth}\centering
\begin{tikzpicture}[x=.064cm, y=.07cm]
  \econaxes{85}{80}{粮}{布}
  \draw[mjSecond, line width=1.1pt] (0,60) -- (60,0);
  \draw[mjPine, line width=1pt, dash pattern=on 4pt off 2pt] (60,0) -- node[pos=.82, anchor=west, slab, mjPine, xshift=3pt] {交换线，斜率 $-1.5$} (13.3,70);
  \node[eqpt] at (30,30) {}; \node[slab, anchor=north east] at (29,29) {自给};
  \node[eqpt, fill=mjPine] at (30,45) {}; \node[slab, anchor=west, mjPine] at (31.5,46) {贸易后 $(30,45)$};
  \draw[shiftarrow] (30,31) -- (30,43.5);
  \node[eqpt, fill=mjSecond] at (60,0) {}; \node[slab, anchor=south west] at (61,1) {只生产粮};
\end{tikzpicture}
\caption{乙国：生产 $(60,0)$，消费 $(30,45)$}
\end{subfigure}
\caption{贸易让两国的消费点都跑到了生产可能性边界之外}
\label{fig:es-trade}
\end{figure}

\section{交换比例的范围}

交换比例不是随便定的。设 1 单位粮换 $p$ 单位布：
\begin{itemize}
  \item 乙国自己生产，放弃 1 单位粮能换来 1 单位布。只有 $p>1$，乙国用粮去换布才比自己织布划算；
  \item 甲国自己生产，1 单位粮要放弃 2 单位布。只有 $p<2$，甲国用布去换粮才比自己种粮划算。
\end{itemize}
所以\textbf{双方都获益的交换比例在两国的机会成本之间}：$1<p<2$（图~\ref{fig:es-tot}）。比例越靠近 2，乙国得到的好处越多；越靠近 1，甲国得到的越多。

\begin{figure}[!htbp]
\centering
\begin{tikzpicture}[x=4cm]
  \fill[mjPine!18] (1,-.18) rectangle (2,.18);
  \draw[mj axis] (0.2,0) -- (2.85,0) node[right, slab] {1 单位粮换多少布};
  \foreach \x in {1,2} { \draw[line width=.8pt] (\x,-.12) -- (\x,.12); }
  \node[slab, below=3pt] at (1,-.12) {1};
  \node[slab, below=3pt] at (2,-.12) {2};
  \node[slab, above=3pt, mjSecond] at (1,.12) {乙国的机会成本};
  \node[slab, above=3pt, mjAccent] at (2,.12) {甲国的机会成本};
  \node[slab, mjPine] at (1.5,-.33) {双方都获益};
  \node[slab, align=center] at (.55,-.65) {乙国宁可自己织布\\（不愿卖粮）};
  \node[slab, align=center] at (2.45,-.65) {甲国宁可自己种粮\\（不愿买粮）};
\end{tikzpicture}
\caption{互利的交换比例落在两国的机会成本之间}
\label{fig:es-tot}
\end{figure}

\section{比较优势从哪里来？谁会受损？}

\paragraph{来源}
\begin{itemize}
  \item \textbf{技术差异}：本讲的模型，两国的生产效率不同（经济学家李嘉图的思路）；
  \item \textbf{资源禀赋}：劳动力多的国家，劳动密集型产品的机会成本低；土地多的国家，农产品的机会成本低；
  \item \textbf{规模经济}：产量越大、单位成本越低时，即使两国原本一样，专门生产一种产品也有好处。
\end{itemize}

\paragraph{模型的假设}
劳动可以在两个行业之间自由、无成本地转移；没有运输成本和关税；充分就业；两国生产的布、粮完全一样；每单位劳动的产量不随产量变化。

\paragraph{分配问题}
分工以后，甲国种粮的人要改行织布。现实中，改行要重新学技能、可能要搬家，有人会失业或收入下降。贸易带来的\textbf{总}收益大于这些人的损失，
理论上获益者可以补偿受损者，使每个人都不比原来差；但这种补偿不会自动发生，需要政策来做，例如转岗培训、失业保险。
所以“贸易对国家有利”和“有人反对贸易”可以同时成立。

\begin{pitfall}
\begin{itemize}
  \item “乙国样样不如甲国，没法和甲国做生意”——错。决定分工的是比较优势，不是绝对优势。
  \item “贸易是零和的，一方赚了另一方就亏了”——错。分工增加了总产量，双方都能获益。
  \item 计算机会成本时要看清表中给的是“每单位劳动的产量”还是“每单位产品需要的劳动”，两种表格的算法正好相反。
\end{itemize}
\end{pitfall}

\begin{oneline}
比较的是机会成本，不是效率；交换比例落在两国机会成本之间，双方都赚；国家赚了，不代表每个人都赚。
\end{oneline}

\begin{exercises}[练一练]
  \item 一小时内，小王能做 6 道数学题或写 3 页作文，小李能做 2 道数学题或写 2 页作文。两人合作完成一份作业（既有数学题又有作文），应当怎样分工？
  \item 在本讲的例子中，如果交换比例是“1 单位粮换 2.5 单位布”，两国还会交换吗？为什么？
  \item 有人说：“某国工资很低，什么产品都比我们便宜，和它做贸易我们只会吃亏。”用比较优势的道理分析这句话。
\end{exercises}

% =============================================================================
\lecture[number=8, short=投票与集体决策]{投票与集体决策：“多数人的意见”存在吗}
% =============================================================================

\begin{keypoints}
\begin{itemize}
  \item 两两多数表决可能出现循环（A 胜 B，B 胜 C，C 又胜 A），这叫\textbf{孔多塞悖论}。出现循环时，谁赢取决于表决的顺序。
  \item 不同的投票规则（相对多数、两轮决选、波达计数……）对同一组选票可能选出不同的赢家。
  \item 人们可能不按真实偏好投票（\textbf{策略性投票}）。
  \item \textbf{阿罗不可能定理}：方案有三个或更多时，没有一种规则能同时满足几条看起来都应该满足的要求。
  \item 如果大家的偏好都是“单峰”的，多数表决不会循环，\textbf{中位选民}的理想方案获胜。
\end{itemize}
\end{keypoints}

\section{两两多数表决与孔多塞悖论}

一个人做选择很简单：选自己最喜欢的。一群人呢？最自然的想法是\textbf{两两比较，多数说了算}。

\begin{definition}[孔多塞赢家]
如果某个方案和其他\textbf{每一个}方案两两表决时都能获得多数，就称它为\textbf{孔多塞赢家}。
\end{definition}

孔多塞赢家看起来是“多数人的意见”最理想的代表。可惜它不一定存在。班委会 5 人要从 A（植物园）、B（科技馆）、C（郊野公园）中选春游地点，各人的排序（从最喜欢到最不喜欢）是：
\begin{center}
\small
\begin{tabular}{lc}
  \toprule
  人数 & 排序 \\
  \midrule
  2 人 & A，B，C \\
  2 人 & B，C，A \\
  1 人 & C，A，B \\
  \bottomrule
\end{tabular}
\end{center}
两两表决：A 对 B，3 比 2，A 胜；B 对 C，4 比 1，B 胜；C 对 A，3 比 2，C 胜。结果是一个圈（图~\ref{fig:es-cycle}）：
\[
  \text{A 胜 B},\quad \text{B 胜 C},\quad \text{C 胜 A}.
\]
每个人自己的排序都是有条理的（喜欢 A 胜过 B、B 胜过 C，就一定喜欢 A 胜过 C，这叫\textbf{传递性}），合起来的“多数意见”却没有传递性。这就是\textbf{孔多塞悖论}。

\begin{figure}[!htbp]
\centering
\begin{tikzpicture}[opt/.style={circle, draw=mjAccent, line width=1pt, fill=mjAccent!8, minimum size=1.15cm, font=\small, align=center},
    win/.style={-{Stealth[length=2.6mm,width=2mm]}, draw=mjSecond, line width=1.1pt}]
  \node[opt] (A) at (90:1.9) {A\\[-2pt]{\scriptsize 植物园}};
  \node[opt] (B) at (210:1.9) {B\\[-2pt]{\scriptsize 科技馆}};
  \node[opt] (C) at (330:1.9) {C\\[-2pt]{\scriptsize 郊野公园}};
  \draw[win] (A) -- node[left=2pt, slab] {3:2} (B);
  \draw[win] (B) -- node[below=2pt, slab] {4:1} (C);
  \draw[win] (C) -- node[right=2pt, slab] {3:2} (A);
  \node[slab, align=left, anchor=west] at (3.2,.3) {箭头从胜者指向败者\\ 三个箭头首尾相接，没有“最好”的方案};
\end{tikzpicture}
\caption{孔多塞循环}
\label{fig:es-cycle}
\end{figure}

\section{议程决定结果}

出现循环时，规则的细节就变得很要紧。设主席安排两轮表决：先在两个方案中表决，胜者再与第三个方案表决。图~\ref{fig:es-agenda} 列出了三种议程：

\begin{figure}[!htbp]
\centering
\begin{tikzpicture}[x=1cm, y=.75cm,
    o/.style={draw=mjInk!60, rounded corners=1.5pt, minimum width=.7cm, minimum height=.5cm, font=\small, fill=white},
    w/.style={draw=mjSecond, line width=1pt, rounded corners=1.5pt, minimum width=.7cm, minimum height=.5cm, font=\small\bfseries, fill=mjSecond!10},
    l/.style={draw=mjInk!50, line width=.6pt}]
  \foreach \dx/\a/\b/\c/\r/\f in {0/A/B/C/A/C, 4.6/B/C/A/B/A, 9.2/A/C/B/C/B} {
    \begin{scope}[xshift=\dx cm]
      \node[o] (a) at (0,2) {\a}; \node[o] (b) at (0,1) {\b}; \node[o] (c) at (0,0) {\c};
      \node[o] (r) at (1.4,1.5) {\r};
      \node[w] (f) at (2.8,.75) {\f};
      \draw[l] (a.east) -- ++(.3,0) |- (r.west); \draw[l] (b.east) -- ++(.3,0) |- (r.west);
      \draw[l] (r.east) -- ++(.3,0) |- (f.west); \draw[l] (c.east) -- ++(1.7,0) |- (f.west);
      \node[slab] at (1.4,-1) {先 \a 对 \b，胜者对 \c};
    \end{scope}}
\end{tikzpicture}
\caption{三种议程，三个不同的赢家：最后才上场的方案获胜}
\label{fig:es-agenda}
\end{figure}

道理很简单：在循环中，每个方案都恰好能赢一个、输一个。最后上场的方案只需打败第一轮的胜者；而它正好能打败的那个方案，主席总可以安排它在第一轮胜出。
所以\textbf{控制议程的人，就能控制结果}。这也是为什么很多组织要预先公开、固定表决程序。

\section{不同的规则，不同的赢家}

除了两两表决，常见的投票规则还有：
\begin{itemize}
  \item \textbf{相对多数制}：每人投一票给最喜欢的方案，得票最多者胜；
  \item \textbf{两轮决选制}：先按相对多数，取前两名再进行一轮二选一；
  \item \textbf{波达计数}：$m$ 个方案时，每人给第 1 名 $m-1$ 分，第 2 名 $m-2$ 分……最后一名 0 分，总分最高者胜。
\end{itemize}

看一个 15 人的例子：
\begin{center}
\small
\begin{tabular}{lc}
  \toprule
  人数 & 排序 \\
  \midrule
  6 人 & A，B，C \\
  5 人 & C，B，A \\
  4 人 & B，C，A \\
  \bottomrule
\end{tabular}
\end{center}
\begin{itemize}
  \item 相对多数：A 得 6 票，C 得 5 票，B 得 4 票，\textbf{A 胜}。
  \item 两轮决选：A、C 进入第二轮。原来投 B 的 4 人更喜欢 C，C 以 $5+4=9$ 比 6 胜，\textbf{C 胜}。
  \item 波达计数：A 得 $6\times2=12$ 分；B 得 $6\times1+5\times1+4\times2=19$ 分；C 得 $5\times2+4\times1=14$ 分，\textbf{B 胜}。
\end{itemize}
同一组选票，三种规则，三个赢家（图~\ref{fig:es-rules}）。再用两两表决检查：B 对 A 是 9 比 6，B 对 C 是 10 比 5，\textbf{B 是孔多塞赢家}；
而相对多数选出的 A，和 B、C 两两表决都会输，是“孔多塞输家”。

\begin{figure}[!htbp]
\centering
\pgfplotsset{rulebar/.style={mojing box, ybar, bar width=13pt, width=4.9cm, height=4.4cm, ymin=0,
    symbolic x coords={A,B,C}, xtick={A,B,C}, bar shift=0pt, enlarge x limits=.3, nodes near coords,
    nodes near coords style={font=\footnotesize}, ytick=\empty, axis line style={draw=mjInk!40}}}
\begin{tikzpicture}
\begin{axis}[rulebar, ymax=8, title={\small 相对多数：A 胜}]
  \addplot[fill=mjGray!35, draw=none] coordinates {(B,4) (C,5)};
  \addplot[fill=mjSecond!70, draw=none] coordinates {(A,6)};
\end{axis}
\end{tikzpicture}\hfill
\begin{tikzpicture}
\begin{axis}[rulebar, ymax=12, title={\small 两轮决选：C 胜}]
  \addplot[fill=mjGray!35, draw=none] coordinates {(A,6)};
  \addplot[fill=mjSecond!70, draw=none] coordinates {(C,9)};
  \addplot[fill=mjGray!12, draw=mjGray!50, dashed, nodes near coords={淘汰}] coordinates {(B,0.6)};
\end{axis}
\end{tikzpicture}\hfill
\begin{tikzpicture}
\begin{axis}[rulebar, ymax=24, title={\small 波达计数：B 胜}]
  \addplot[fill=mjGray!35, draw=none] coordinates {(A,12) (C,14)};
  \addplot[fill=mjSecond!70, draw=none] coordinates {(B,19)};
\end{axis}
\end{tikzpicture}
\caption{同一组选票（6 人 A>B>C，5 人 C>B>A，4 人 B>C>A），三种规则选出三个不同的赢家}
\label{fig:es-rules}
\end{figure}

相对多数制的毛病在这里看得很清楚：B、C 的支持者都不喜欢 A，但他们的票被 B、C 两个相近的方案\textbf{分散}了，于是 A 以少数票胜出。

\section{策略性投票}

上面的例子里，最喜欢 B 的 4 个人（排序 B，C，A）如果在相对多数制下照实投 B，结果是他们最不喜欢的 A 胜出。
如果他们改投 C，C 得 9 票胜出——对他们来说，C 比 A 好。\textbf{不按真实偏好投票，反而得到对自己更好的结果}，这叫\textbf{策略性投票}。
“不想浪费选票”就是这个意思。

在第二节的分轮表决中，同样可能出现策略性投票：委员预见到第二轮的结果，第一轮就故意投给自己不是最喜欢的方案。
数学上已经证明（吉巴德—萨特斯韦特定理）：方案有三个或更多时，任何一种不是“独裁”、每个方案都有可能胜出的规则，都会在某些情况下给人策略性投票的机会。

\section{阿罗不可能定理}

能不能设计一个“完美”的规则？经济学家阿罗（1951）证明了：方案有三个或更多时，不存在一种把个人排序汇总为集体排序的规则，能同时满足下面四条：
\begin{enumerate}
  \item \textbf{无限制}：每个人可以有任意的排序，规则都能给出一个有传递性的集体排序；
  \item \textbf{一致性（帕累托原则）}：如果人人都认为 X 比 Y 好，集体也认为 X 比 Y 好；
  \item \textbf{无关方案独立}：集体对 X、Y 的排序，只取决于每个人对 X、Y 谁好谁坏的看法，不受第三个方案排在哪里的影响；
  \item \textbf{非独裁}：没有哪一个人的排序总是直接成为集体排序。
\end{enumerate}

两两多数表决违反第 1 条（会出现循环）；波达计数违反第 3 条。看一个例子：5 个人中，3 人排序 A，B，C，2 人排序 B，C，A。
波达计数下 A 得 6 分，B 得 7 分，B 排在 A 前面。现在这 3 人把排序改成 A，C，B——他们对 A、B 的看法没变（仍然是 A 比 B 好），只是把 C 往前挪了。
重新计分：A 6 分，B 4 分，C 5 分，A 跑到了 B 前面。没有人改变对 A、B 的看法，集体对 A、B 的排序却反过来了。

阿罗定理的意思\textbf{不是}“民主不可能”，而是：\textbf{每种规则都有代价}，选择规则就是在几条好的性质之间做取舍。

\section{单峰偏好与中位选民}

循环并不总会出现。很多问题的方案可以排在一条线上：班费收多少、空调温度定几度、税率定多高。这时每个人通常有一个最理想的位置，离它越远越不满意。
这样的偏好叫\textbf{单峰偏好}：画出来只有一个“山峰”。

\begin{figure}[!htbp]
\centering
\begin{tikzpicture}[x=.105cm, y=1.6cm]
  \draw[mj axis] (0,0) -- (106,0) node[right, slab] {班费（元）};
  \draw[mj axis] (0,0) -- (0,1.35) node[above, slab] {满意程度};
  \begin{scope}
    \clip (0,0) rectangle (106,1.3);
    \foreach \p/\c in {10/mjAccent, 25/mjAccent, 40/mjPine, 70/mjSecond, 90/mjSecond}
      \draw[\c, line width=.9pt] (\p-40,{1-40/45}) -- (\p,1) -- (\p+40,{1-40/45});
  \end{scope}
  \foreach \p/\c/\n in {10/mjAccent/1, 25/mjAccent/2, 40/mjPine/3, 70/mjSecond/4, 90/mjSecond/5} {
    \node[slab, \c, above] at (\p,1) {\n 号};
    \draw[\c!60, line width=.4pt, dash pattern=on 1.5pt off 1.2pt] (\p,0) -- (\p,1);
    \node[slab, below] at (\p,0) {\p};
  }
  \fill[mjPine] (40,0) circle (2.2pt);
  \node[slab, mjPine, align=center] at (40,-.42) {中位数 40};
\end{tikzpicture}
\caption{5 个人的单峰偏好：各自理想的班费是 10、25、40、70、90 元，中位数是 40 元}
\label{fig:es-single-peak}
\end{figure}

\begin{theorem}[中位选民定理]
如果方案排在一条线上，每个人的偏好都是单峰的，人数为奇数，那么所有人理想位置的\textbf{中位数}是孔多塞赢家：它和任何其他方案两两表决都会胜出。
\end{theorem}

\begin{proof}[理由]
设中位数是 $m$，拿它和任意一个更大的方案 $x>m$ 比较。理想位置不超过 $m$ 的人（超过半数）都在 $m$ 的同一侧，对他们来说 $x$ 比 $m$ 更远，所以都投 $m$，$m$ 胜。
$x<m$ 的情况同理。
\end{proof}

图~\ref{fig:es-single-peak} 中，40 元对 60 元：理想位置在 10、25、40 的三个人都更喜欢 40，40 以 3 比 2 胜。
这个定理解释了一个政治现象：两党竞争时，为了赢得中间的选民，双方的政策主张常常向中间靠拢。

\paragraph{为什么前面会出现循环}
回到第一节的循环例子。把 A、B、C 按某种顺序排成一行，画出三类人对它们的评价（图~\ref{fig:es-not-single}）：总有一类人把\textbf{中间}的方案排在最后，他的偏好“两头高、中间低”，有两个峰。
换任何一种排列都是这样，因为 A、B、C 各是某一类人的最后一名。偏好不是单峰的，中位选民定理就不适用，循环才有了可能。

\begin{figure}[!htbp]
\centering
\begin{tikzpicture}[x=1.5cm, y=.9cm]
  \draw[mj axis] (0,0) -- (3.6,0);
  \draw[mj axis] (0,0) -- (0,3.6) node[above, slab] {排名（越高越喜欢）};
  \foreach \x/\l in {1/A, 2/B, 3/C} \node[slab, below] at (\x,0) {\l};
  \foreach \y/\l in {1/第 3, 2/第 2, 3/第 1} \node[slab, left] at (0,\y) {\l};
  \draw[mjAccent, line width=1pt] (1,3) -- (2,2) -- (3,1);
  \draw[mjPine, line width=1pt] (1,1) -- (2,3) -- (3,2);
  \draw[mjSecond, line width=1.4pt] (1,2) -- (2,1) -- (3,3);
  \foreach \p in {(1,3),(2,2),(3,1)} \fill[mjAccent] \p circle (2pt);
  \foreach \p in {(1,1),(2,3),(3,2)} \fill[mjPine] \p circle (2pt);
  \foreach \p in {(1,2),(2,1),(3,3)} \fill[mjSecond] \p circle (2pt);
  \node[slab, mjAccent, anchor=west] at (3.9,3) {A，B，C：单峰};
  \node[slab, mjPine, anchor=west] at (3.9,2.3) {B，C，A：单峰};
  \node[slab, mjSecond, anchor=west] at (3.9,1.6) {C，A，B：两个峰（中间最差）};
\end{tikzpicture}
\caption{循环例子中，按 A、B、C 排列时，第三类人的偏好不是单峰的}
\label{fig:es-not-single}
\end{figure}

\begin{pitfall}
\begin{itemize}
  \item “多数人的意见”不一定存在：两两表决可能循环，这时说“多数人支持某方案”要看是怎么表决的。
  \item 投票结果是\textbf{偏好和规则共同}决定的，换一种规则可能换一个赢家。分析题里要说清用的是什么规则。
  \item 算波达计数时，注意 $m$ 个方案时第一名是 $m-1$ 分，最后一名是 0 分。
\end{itemize}
\end{pitfall}

\begin{oneline}
集体的选择取决于规则：循环让议程说了算，规则不同赢家不同；偏好单峰时，中位选民说了算。
\end{oneline}

\begin{exercises}[练一练]
  \item 9 人对 A、B、C 的排序为：4 人 A，B，C；3 人 B，C，A；2 人 C，B，A。分别求相对多数制、两轮决选制、波达计数的赢家，并判断是否存在孔多塞赢家。
  \item 举例说明：孔多塞赢家不一定是波达计数的赢家。
  \item 5 名同学按距离排列的出游地点理想位置分别为 2、3、5、9、10 公里，偏好都是单峰的。多数表决会选哪里？如果又加入两位理想位置都是 10 公里的同学呢？
\end{exercises}

% =============================================================================
\lecture[number=9, short=读懂数字]{读懂数字：比率、平均与辛普森悖论}
% =============================================================================

\begin{keypoints}
\begin{itemize}
  \item 看到一个数，先问“\textbf{分母是谁}”。总数多不等于比率高。
  \item 总体的比率是各组比率的\textbf{加权平均}，权重是各组人数的占比。
  \item \textbf{辛普森悖论}：每一组里甲都比乙好，合起来甲却比乙差。原因是两者各组的人数构成不同。
  \item 罕见的病，即使检测很准，检测阳性的人里也可能大多数没病：要考虑\textbf{基础比率}。
  \item 平均数会被极端值拉动，偏斜的数据要看中位数；图表的纵轴起点也会误导人。
\end{itemize}
\end{keypoints}

\section{先问分母}

“某市交通事故中，85\% 的肇事司机是本地人。”这能说明本地人开车更不小心吗？不能。如果该市路上 90\% 的司机是本地人，85\% 反而说明本地司机出事的\textbf{比率}略低。
要比较的是“本地司机中出事的比例”和“外地司机中出事的比例”，而不是“出事司机中本地人的比例”。

同样，“A 市去年发生了 300 起盗窃，B 市只有 100 起”，如果 A 市人口是 B 市的 5 倍，A 市其实更安全。\textbf{比较之前，先把数字换算成同样分母的比率。}

\section{加权平均}

某医院去年做了 500 台同一种手术：轻症病人 100 人，存活 95 人（95\%）；重症病人 400 人，存活 280 人（70\%）。总存活率是
\[
  \frac{95+280}{500}=75\% = \underbrace{\frac{100}{500}}_{\text{轻症占比}}\times95\%+\underbrace{\frac{400}{500}}_{\text{重症占比}}\times70\%.
\]
\textbf{总比率是各组比率的加权平均，权重是各组人数的占比。}可以把它想成跷跷板：各组比率是数轴上的点，人数是点的重量，总比率是平衡点。
重的一端把平衡点拉向自己（图~\ref{fig:es-seesaw}）。

\begin{figure}[!htbp]
\centering
\begin{tikzpicture}[x=.16cm, y=1cm]
  \foreach \yy/\name/\lo/\hi/\wlo/\whi/\avg/\col/\llo/\lhi in
      {2.2/甲医院/70/95/0.8/0.2/75/mjAccent/重症/轻症, 0/乙医院/60/90/0.2/0.8/84/mjSecond/重症/轻症} {
    \draw[mj thin] (55,\yy) -- (100,\yy);
    \foreach \t in {60,70,80,90,100} \draw (\t,\yy-.06) -- (\t,\yy+.06);
    \node[font=\scriptsize, mjGray, right] at (100,\yy) {100\%};
    \node[font=\scriptsize, mjGray, left] at (55,\yy) {};
    \fill[\col!85] (\lo,\yy) circle[radius={0.08cm+0.22cm*\wlo}];
    \fill[\col!85] (\hi,\yy) circle[radius={0.08cm+0.22cm*\whi}];
    \node[slab, \col, above=8pt] at (\lo,\yy) {\llo\ \lo\%};
    \node[slab, \col, above=8pt] at (\hi,\yy) {\lhi\ \hi\%};
    \fill[mjPine] (\avg,\yy-.05) -- ++(-.9,-.36) -- ++(1.8,0) -- cycle;
    \node[slab, mjPine, below=12pt] at (\avg,\yy) {总 \avg\%};
    \node[lab, anchor=east] at (53,\yy) {\name};
  }
\end{tikzpicture}
\caption{加权平均像跷跷板：圆点大小表示人数，三角形是平衡点（总存活率）}
\label{fig:es-seesaw}
\end{figure}

\section{辛普森悖论}

图~\ref{fig:es-seesaw} 下面那条线是另一家医院的数据。把两家放在一起：
\begin{center}
\small
\begin{tabular}{l|cc|cc|c}
  \toprule
  & \multicolumn{2}{c|}{轻症} & \multicolumn{2}{c|}{重症} & 合计 \\
  & 人数 & 存活率 & 人数 & 存活率 & 存活率 \\
  \midrule
  甲医院 & 100 & \textbf{95\%} & 400 & \textbf{70\%} & 75\% \\
  乙医院 & 400 & 90\% & 100 & 60\% & \textbf{84\%} \\
  \bottomrule
\end{tabular}
\end{center}
轻症、重症分开看，甲医院都更好；合起来看，乙医院反而更好。这种“分组的结论和合并的结论方向相反”的现象叫\textbf{辛普森悖论}。

原因就在跷跷板上：甲医院收治的大多是重症病人，平衡点被拉向 70\% 那一端；乙医院大多是轻症病人，平衡点被拉向 90\% 那一端。
\textbf{合并后比较的，其实是两家医院病人构成的差别，而不是医疗水平的差别。}

图~\ref{fig:es-simpson} 把这件事画得更清楚：横轴是重症病人的占比，纵轴是总存活率。每家医院的总存活率随重症占比沿一条直线变化，
甲医院的直线\textbf{处处}高于乙医院；但甲医院实际处在右边（重症占 80\%），乙医院处在左边（重症占 20\%），于是甲的点反而低于乙的点。

\begin{figure}[!htbp]
\centering
\begin{tikzpicture}
\begin{axis}[mojing box, width=9.5cm, height=6.4cm, xmin=0, xmax=1, ymin=55, ymax=100,
    xlabel={重症病人的占比}, ylabel={总存活率（\%）}, xtick={0,.2,.5,.8,1}, xticklabels={0,20\%,50\%,80\%,100\%},
    ytick={60,70,80,90,100}, grid=major, legend pos=north east]
  \addplot[mjAccent, domain=0:1, samples=2, line width=1.1pt] {95-25*x};
  \addplot[mjSecond, domain=0:1, samples=2, line width=1.1pt] {90-30*x};
  \legend{甲医院,乙医院}
  \node[eqpt, fill=mjAccent, minimum size=5pt] at (axis cs:.8,75) {};
  \node[eqpt, fill=mjSecond, minimum size=5pt] at (axis cs:.2,84) {};
  \node[slab, mjAccent, anchor=west] (la) at (axis cs:.5,88) {甲的实际位置 75\%};
  \draw[mjAccent, line width=.4pt] (la.south) -- (axis cs:.79,75.8);
  \node[slab, mjSecond, anchor=west] (lb) at (axis cs:.03,64) {乙的实际位置 84\%};
  \draw[mjSecond, line width=.4pt] (lb.north) -- (axis cs:.2,82.8);
  \draw[guide] (axis cs:.5,55) -- (axis cs:.5,82.5);
  \node[eqpt, fill=mjAccent] at (axis cs:.5,82.5) {};
  \node[eqpt, fill=mjSecond] at (axis cs:.5,75) {};
  \node[slab, anchor=south west] at (axis cs:.51,56) {同样构成下比较};
\end{axis}
\end{tikzpicture}
\caption{辛普森悖论的图示：甲的直线处处在上方，但两家医院在横轴上的位置不同}
\label{fig:es-simpson}
\end{figure}

\paragraph{怎样公平地比较}
\begin{itemize}
  \item \textbf{分组比较}：在轻症内部比、在重症内部比。
  \item \textbf{标准化}：用同一个病人构成去算两家医院的总存活率。比如都按“轻症、重症各一半”计算：甲 $\frac{95\%+70\%}2=82.5\%$，乙 $\frac{90\%+60\%}2=75\%$，甲更好。
        在图~\ref{fig:es-simpson} 中，就是在同一条竖线上比较两条直线的高度。
\end{itemize}
按什么分组也要想一想：分组的变量应该是“在接受治疗之前就已经存在、又会影响结果”的因素，例如病情轻重。如果分组变量本身是治疗造成的结果，分组反而会误导。

\section{基础比率：阳性就是有病吗}

某种病的患病率是 1\%。一种检测方法：病人检测为阳性的概率是 90\%，没病的人被误判为阳性的概率是 5\%。如果你检测为阳性，真的有病的可能有多大？

很多人会说“90\% 左右”。用 10000 个人来算一遍（图~\ref{fig:es-tree-test}）：
\begin{itemize}
  \item 100 人有病，其中 90 人阳性；
  \item 9900 人没病，其中 $9900\times5\%=495$ 人阳性。
\end{itemize}
阳性的人共 $90+495=585$ 人，其中真正有病的只有 90 人，占 $\frac{90}{585}\approx15\%$。

\begin{figure}[!htbp]
\centering
\begin{tikzpicture}[
    box/.style={draw=mjInk!50, rounded corners=2pt, font=\small, inner sep=4pt, fill=white, align=center},
    posbox/.style={box, draw=mjSecond, fill=mjSecond!10, line width=.9pt},
    e/.style={draw=mjInk!50, line width=.6pt}]
  \node[box] (all) at (0,0) {10000 人};
  \node[box] (s) at (-3.2,-1.4) {有病 100 人\\{\footnotesize 患病率 1\%}};
  \node[box] (h) at (3.2,-1.4) {没病 9900 人};
  \node[posbox] (sp) at (-4.6,-3.1) {阳性 90};
  \node[box] (sn) at (-1.8,-3.1) {阴性 10};
  \node[posbox] (hp) at (1.8,-3.1) {阳性 495};
  \node[box] (hn) at (4.6,-3.1) {阴性 9405};
  \draw[e] (all) -- (s); \draw[e] (all) -- (h);
  \draw[e] (s) -- node[left, slab] {90\%} (sp); \draw[e] (s) -- (sn);
  \draw[e] (h) -- node[left, slab] {5\%} (hp); \draw[e] (h) -- (hn);
  \node[slab, mjSecond] at (0,-4.3) {阳性共 $90+495=585$ 人，其中有病的只占 $90/585\approx15\%$};
\end{tikzpicture}
\caption{用“自然频数”算条件概率：把百分比换成具体的人数}
\label{fig:es-tree-test}
\end{figure}

原因是：没病的人太多了（9900 人），哪怕只有 5\% 被误判，人数（495）也远远超过真正的病人（90）。这个“原本有多少人有病”的比例叫\textbf{基础比率}，忽略它就会犯\textbf{基础比率谬误}。
同样的道理适用于很多“筛查”：用人脸识别在人群中找逃犯、用软件判断作业是否抄袭、用检测筛查罕见病……对象越罕见，误报所占的比例就越大。

\begin{remark}
这就是条件概率的贝叶斯公式：
\[
  P(\text{有病}\mid\text{阳性})=\frac{P(\text{有病})\,P(\text{阳性}\mid\text{有病})}{P(\text{阳性})}
  =\frac{0.01\times0.9}{0.01\times0.9+0.99\times0.05}\approx0.15 .
\]
用具体人数来算，和用公式算是一回事，但更不容易出错。
\end{remark}

\section{平均数与中位数}

一家小公司 10 个人，9 人月薪 5000 元，老板月薪 100 万元。平均月薪是 $(9\times5000+1000000)\div10=104500$ 元，可 10 个人里有 9 个人的月薪只是平均数的二十分之一。
这时\textbf{中位数}（排在正中间的数，这里是 5000 元）更能代表“一般人”。

收入、房价、网络视频的播放量，这类数据大多是\textbf{右偏}的：大多数人集中在较低处，少数极高的值拖出一条长尾巴，把平均数拉到中位数的右边（图~\ref{fig:es-skew}）。
所以看到“人均收入”，最好也问一问中位数是多少。

\begin{figure}[!htbp]
\centering
\begin{tikzpicture}
\begin{axis}[mojing, width=10cm, height=5.4cm, xmin=0, xmax=20, ymin=0, ymax=.19,
    xlabel={月收入（千元）}, ylabel={人数比例}, ytick=\empty, xtick={0,5,10,15,20}]
  \addplot[mjAccent, fill=mjAccent!10, domain=.05:20, samples=150, line width=1pt]
    {exp(-(ln(x)-ln(5))^2/(2*0.36))/(x*0.6*sqrt(2*pi))} \closedcycle;
  \draw[mjInk, line width=.7pt] (axis cs:3.49,0) -- (axis cs:3.49,.1592);
  \draw[mjPine, line width=.9pt] (axis cs:5,0) -- (axis cs:5,.133);
  \draw[mjSecond, line width=.9pt] (axis cs:5.99,0) -- (axis cs:5.99,.1062);
  \node[slab, anchor=south west] at (axis cs:3.3,.163) {最多的人在这里};
  \node[slab, mjPine, anchor=west] at (axis cs:7,.14) {中位数：一半人在它以下};
  \draw[mjPine, line width=.4pt] (axis cs:7,.14) -- (axis cs:5.1,.128);
  \node[slab, mjSecond, anchor=west] at (axis cs:8.5,.1) {平均数：被右边的长尾拉高};
  \draw[mjSecond, line width=.4pt] (axis cs:8.5,.1) -- (axis cs:6.1,.095);
\end{axis}
\end{tikzpicture}
\caption{右偏分布：平均数 $>$ 中位数}
\label{fig:es-skew}
\end{figure}

\section{图表的常见陷阱}

同一组数据，换一种画法，给人的印象可以完全不同。图~\ref{fig:es-axis} 中两个品牌的满意度是 96\% 和 98\%：
左图纵轴从 0 开始，两根柱子几乎一样高；右图纵轴从 95\% 开始，B 的柱子看起来是 A 的 3 倍。

\begin{figure}[!htbp]
\centering
\pgfplotsset{axbar/.style={mojing box, ybar, bar width=20pt, width=5.6cm, height=5cm,
    symbolic x coords={A,B}, xtick=data, enlarge x limits=.5, nodes near coords={\pgfmathprintnumber\pgfplotspointmeta\%},
    nodes near coords style={font=\footnotesize}}}
\begin{tikzpicture}
\begin{axis}[axbar, ymin=0, ymax=105, title={\small 纵轴从 0 开始}]
  \addplot[fill=mjAccent!55, draw=none] coordinates {(A,96) (B,98)};
\end{axis}
\end{tikzpicture}\qquad
\begin{tikzpicture}
\begin{axis}[axbar, ymin=95, ymax=98.6, title={\small 纵轴从 95\% 开始}, ytick={95,96,97,98}]
  \addplot[fill=mjSecond!55, draw=none] coordinates {(A,96) (B,98)};
\end{axis}
\end{tikzpicture}
\caption{截断纵轴会夸大差别}
\label{fig:es-axis}
\end{figure}

其他常见陷阱：
\begin{itemize}
  \item \textbf{只给增长率，不给基数}：“销量是去年的 3 倍”，可能只是从 1 件变成 3 件；
  \item \textbf{挑选时间段}：只截取上涨的那几个月；
  \item \textbf{用面积或体积表示数量}：数量翻倍，把图标的长和宽都放大一倍，面积变成 4 倍，看起来像是翻了两番；
  \item \textbf{两个纵轴}：左右两边刻度不同，可以让任意两条曲线看起来“同步”。
\end{itemize}

\begin{pitfall}
\begin{itemize}
  \item 合并的数据和分组的数据结论相反时，不要急着说哪个“对”，先找出各组人数构成的差别。
  \item 算条件概率时，分清“有病的人中阳性的比例”和“阳性的人中有病的比例”，这两个数可能差很多。
\end{itemize}
\end{pitfall}

\begin{oneline}
先问分母，再问权重：总比率是加权平均，构成不同就可能出现辛普森悖论；罕见事件要先看基础比率。
\end{oneline}

\begin{exercises}[练一练]
  \item 1986 年一项肾结石治疗研究常被用来说明辛普森悖论。A 法：小结石 87 人中治愈 81 人，大结石 263 人中治愈 192 人；B 法：小结石 270 人中治愈 234 人，大结石 80 人中治愈 55 人。
        分别计算两种方法在小结石、大结石和全部病人中的治愈率，并解释结果。
  \item 某病的患病率为 0.1\%，检测灵敏度为 99\%（病人阳性的概率），误报率为 1\%（没病的人阳性的概率）。(1) 检测阳性的人真正有病的概率是多少？
        (2) 若阳性者再独立地检测一次，仍为阳性，这时有病的概率是多少？
  \item 某公司招聘广告写着“员工平均月薪 2 万元”。你会追问哪些问题？
\end{exercises}

% =============================================================================
\lecture[number=10, short=抽样调查]{抽样调查：一小部分人为什么能代表全体}
% =============================================================================

\begin{keypoints}
\begin{itemize}
  \item 抽样是用一小部分人推断全体。关键不在人数多，而在\textbf{有没有代表性}。
  \item 误差分两种：\textbf{随机误差}（运气，样本越大越小）和\textbf{系统偏差}（方法有毛病，样本再大也消不掉）。
  \item 估计一个比例时，误差范围大约是 $\dfrac1{\sqrt n}$：1000 人约 $\pm3\%$。误差减半，样本要变成 4 倍。
  \item \textbf{概率抽样}（每个人被抽中的概率已知）是避免偏差的基础：简单随机、分层、整群、系统抽样。
  \item 不回答、问题的措辞、“想显得好一点”都会造成偏差；敏感问题可以用随机化回答。
\end{itemize}
\end{keypoints}

\section{从总体到样本}

想知道全市高中生每天睡几个小时，不可能把每个人都问一遍。我们问其中一部分人，用他们的情况推断全体。几个名词：
\begin{itemize}
  \item \textbf{目标总体}：想了解的全体，例如全市所有高中生；
  \item \textbf{抽样框}：实际用来抽人的名单，例如各校提供的学生名册；
  \item \textbf{样本}：从名单中抽出来的人；
  \item \textbf{回答者}：样本中真正回答了问卷的人。
\end{itemize}
从目标总体到最终的回答者，每一步都可能出偏差（图~\ref{fig:es-frame}）：名单漏掉了一些人（例如借读生），抽样时运气不好，抽中的人不愿回答……

\begin{figure}[!htbp]
\centering
\begin{tikzpicture}
  \fill[mjAccent!7] (0,0) ellipse (4.1 and 2.1);
  \draw[mjAccent, line width=1pt] (0,0) ellipse (4.1 and 2.1);
  \fill[mjOchre!14, opacity=.9] (.9,-.1) ellipse (3.6 and 1.6);
  \draw[mjOchre, line width=1pt, dash pattern=on 4pt off 2pt] (.9,-.1) ellipse (3.6 and 1.6);
  \fill[mjPine!20] (1.2,-.2) circle (1.0);
  \draw[mjPine, line width=1pt] (1.2,-.2) circle (1.0);
  \fill[mjSecond!35] (1.45,-.35) circle (.55);
  \draw[mjSecond, line width=1pt] (1.45,-.35) circle (.55);
  \node[lab, mjAccent] at (-2.6,1.2) {目标总体};
  \node[lab, mjOchre] at (3.1,1.3) {抽样框（名单）};
  \node[slab, mjPine] at (1.2,.55) {样本};
  \node[slab, mjSecond] at (1.45,-.35) {回答者};
  \node[slab, align=center] at (-2.6,-.3) {名单漏掉的人\\（覆盖不足）};
  \node[slab, align=center] at (4.85,-1.95) {名单里多出来的人\\（不属于总体）};
  \draw[mjInk!50, line width=.3pt] (4.1,-1.6) -- (4.0,-1.2);
  \draw[mjInk!50, line width=.3pt] (2.6,.7) -- (5.2,1.5) node[right, slab, align=left] {抽中却没回答\\（无回答）};
\end{tikzpicture}
\caption{目标总体、抽样框、样本与回答者：每一层的差别都可能带来偏差}
\label{fig:es-frame}
\end{figure}

\section{两种误差：准不准与稳不稳}

抽样调查的误差有两种性质完全不同的来源，用打靶来比喻最清楚（图~\ref{fig:es-target}）：
\begin{itemize}
  \item \textbf{系统偏差}：方法本身有毛病，每次都朝同一个方向偏。就像瞄准器歪了，打得再多也打不中靶心。偏差说的是“准不准”。
  \item \textbf{随机误差}：只是运气不同，这次偏高、下次偏低，平均起来在靶心附近。样本越大，随机误差越小。它说的是“稳不稳”。
\end{itemize}

\begin{figure}[!htbp]
\centering
\newcommand{\target}[4]{% x 偏移, y 偏移, 散布, 标题
  \begin{tikzpicture}
    \foreach \r/\c in {1.15/mjLight, .8/white, .45/mjLight} \fill[\c] (0,0) circle (\r);
    \foreach \r in {1.15,.8,.45} \draw[mjInk!45, line width=.4pt] (0,0) circle (\r);
    \fill[mjInk!70] (0,0) circle (1.5pt);
    \foreach \a/\b in {-0.45/1.45, 0.74/-0.42, -0.78/-0.44, 0.09/0.03, 0.27/-1.76, 1.08/-0.01, 0.20/-0.05, -0.86/0.55, 0.34/-0.12, -0.63/0.77}
      \fill[mjSecond] ({#1+#3*\a},{#2+#3*\b}) circle (1.6pt);
    \node[slab, align=center] at (0,-1.55) {#4};
  \end{tikzpicture}}
\target{0}{0}{.1}{偏差小，随机误差小\\ \textbf{又准又稳}}\hfill
\target{0}{0}{.36}{偏差小，随机误差大\\ 准而不稳}\hfill
\target{.5}{.42}{.1}{偏差大，随机误差小\\ 稳而不准}\hfill
\target{.5}{.42}{.36}{偏差大，随机误差大\\ 又不准又不稳}
\caption{偏差与随机误差。红点是多次调查的结果，靶心是真实值。}
\label{fig:es-target}
\end{figure}

\section{样本量有多大用}

估计一个比例（例如支持率）时，如果抽样是随机的，\textbf{95\% 的情况下误差不超过约 $\dfrac1{\sqrt n}$}：
\begin{center}
\small
\begin{tabular}{l*{5}{c}}
  \toprule
  样本量 $n$ & 100 & 400 & 1000 & 2500 & 10000 \\
  误差范围约 & $\pm10\%$ & $\pm5\%$ & $\pm3.2\%$ & $\pm2\%$ & $\pm1\%$ \\
  \bottomrule
\end{tabular}
\end{center}
两个值得记住的事实：
\begin{itemize}
  \item 误差和 $\sqrt n$ 成反比：误差减半，样本要变成 4 倍。样本量到了一两千以后，再增加的作用就不大了。
  \item 只要总体比样本大得多，误差几乎\textbf{和总体有多大无关}：调查一个城市和调查全国，需要的样本量差不多。好比尝一勺汤就知道咸淡，不管锅有多大——前提是汤\textbf{搅匀了}（随机抽样）。
\end{itemize}

但样本量只能减小随机误差，\textbf{对系统偏差无能为力}。图~\ref{fig:es-n-bias} 中，真实支持率是 50\%。无偏的调查，样本越大，结果越集中在 50\% 附近；
有偏的调查，样本越大，结果越集中在一个错误的值（58\%）附近——精度越高，错得越“有把握”。

\begin{figure}[!htbp]
\centering
\begin{tikzpicture}
\begin{axis}[mojing box, xmode=log, width=10cm, height=6cm, xmin=100, xmax=10000, ymin=35, ymax=72,
    xlabel={样本量 $n$（对数刻度）}, ylabel={调查结果（\%）}, xtick={100,1000,10000}, xticklabels={100,1000,10000},
    ytick={40,50,58,65}, legend pos=north east, legend style={fill=white}]
  \addplot[name path=u1, draw=none, domain=100:10000, samples=40] {50+100/sqrt(x)};
  \addplot[name path=l1, draw=none, domain=100:10000, samples=40] {50-100/sqrt(x)};
  \addplot[mjAccent!25] fill between[of=u1 and l1];
  \addplot[name path=u2, draw=none, domain=100:10000, samples=40] {58+100/sqrt(x)};
  \addplot[name path=l2, draw=none, domain=100:10000, samples=40] {58-100/sqrt(x)};
  \addplot[mjSecond!25] fill between[of=u2 and l2];
  \addplot[mjAccent, domain=100:10000, samples=2, line width=1pt] {50};
  \addplot[mjSecond, domain=100:10000, samples=2, line width=1pt, dashed] {58};
  \node[slab, mjAccent, anchor=north east] at (axis cs:9500,48.6) {无偏：收缩到真实值 50\%};
  \node[slab, mjSecond, anchor=south east] at (axis cs:9500,60.4) {有偏：收缩到错误的 58\%};
\end{axis}
\end{tikzpicture}
\caption{样本越大，结果的波动范围越窄；但有偏的调查只会越来越确定地给出错误的答案}
\label{fig:es-n-bias}
\end{figure}

\paragraph{两个著名的失败}
\begin{itemize}
  \item \textbf{1936 年}，美国《文学文摘》杂志收回两百多万份模拟选票，预测兰登将击败罗斯福，结果罗斯福大胜。名单来自电话簿、汽车登记等，在大萧条时期偏向富裕人群；
        寄回选票的人也只是一部分。盖洛普用规模小得多但设计得当的样本预测对了。
  \item \textbf{1948 年}，多家民调机构预测杜威将击败杜鲁门，有报纸甚至提前印出了“杜威击败杜鲁门”的标题，结果杜鲁门获胜。
        原因之一是当时使用\textbf{配额抽样}：只规定访问多少男性、多少城市居民等，具体访问谁由访员自己挑，访员往往挑了更容易接触的人。
\end{itemize}

\section{概率抽样的四种方法}

避免偏差的关键是\textbf{让“谁被抽中”由概率决定，而不是由人决定}。常用的四种方法见图~\ref{fig:es-sampling}（60 名学生，每两行是一个年级）。

\begin{figure}[!htbp]
\centering
\newcommand{\popgrid}{%
  \fill[mjAccent!7] (-.3,-.3) rectangle (9.3,1.3);
  \fill[mjPine!8] (-.3,1.7) rectangle (9.3,3.3);
  \fill[mjOchre!10] (-.3,3.7) rectangle (9.3,5.3);
  \foreach \x in {0,...,9} \foreach \y in {0,...,5} \fill[mjInk!30] (\x,\y) circle (2.6pt);}
\newcommand{\pick}[1]{\foreach \x/\y in {#1} \fill[mjSecond] (\x,5-\y) circle (4.2pt);}
\begin{subfigure}{.48\textwidth}\centering
\begin{tikzpicture}[x=.5cm, y=.42cm]
  \popgrid
  \pick{1/0, 3/0, 6/1, 3/2, 4/2, 7/2, 2/3, 4/3, 0/4, 2/4, 5/5, 8/5}
\end{tikzpicture}
\caption{简单随机：每人被抽中的机会相同}
\end{subfigure}\hfill
\begin{subfigure}{.48\textwidth}\centering
\begin{tikzpicture}[x=.5cm, y=.42cm]
  \popgrid
  \pick{4/0, 5/0, 2/1, 9/1, 0/2, 1/3, 3/3, 6/3, 0/4, 2/4, 6/4, 7/5}
  \foreach \y/\l in {4.5/高一, 2.5/高二, .5/高三} \node[slab, right] at (9.4,\y) {\l};
\end{tikzpicture}
\caption{分层：每个年级各随机抽 4 人}
\end{subfigure}

\vspace{.6\baselineskip}
\begin{subfigure}{.48\textwidth}\centering
\begin{tikzpicture}[x=.5cm, y=.42cm]
  \popgrid
  \foreach \y in {0,...,5} { \draw[mjInk!35, line width=.4pt, rounded corners=2pt] (-.35,\y-.35) rectangle (4.35,\y+.35);
                             \draw[mjInk!35, line width=.4pt, rounded corners=2pt] (4.65,\y-.35) rectangle (9.35,\y+.35); }
  \pick{5/1, 6/1, 7/1, 8/1, 9/1, 0/4, 1/4, 2/4, 3/4, 4/4}
\end{tikzpicture}
\caption{整群：随机抽 2 个班（框），班里全部调查}
\end{subfigure}\hfill
\begin{subfigure}{.48\textwidth}\centering
\begin{tikzpicture}[x=.5cm, y=.42cm]
  \popgrid
  \pick{2/0, 8/0, 4/1, 0/2, 6/2, 2/3, 8/3, 4/4, 0/5, 6/5}
\end{tikzpicture}
\caption{系统：按名单顺序，从第 3 人起每 6 人抽 1 人}
\end{subfigure}
\caption{四种概率抽样（红点为被抽中的学生）}
\label{fig:es-sampling}
\end{figure}

\begin{center}
\small
\begin{tabular}{p{2cm}p{5.3cm}p{5.3cm}}
  \toprule
  方法 & 优点 & 缺点 \\
  \midrule
  简单随机 & 最简单，理论最清楚 & 需要完整名单；可能碰巧某个年级抽得太少 \\
  分层 & 各层都有代表，层间差别大时更精确 & 要事先知道每个人属于哪一层 \\
  整群 & 省时省钱，不需要全体名单，只要班级名单 & 同一个班的人往往相似，精度较低 \\
  系统 & 操作方便 & 名单有周期性时可能出问题（例如每 6 人恰好是一个寝室） \\
  \bottomrule
\end{tabular}
\end{center}

实际调查常把几种方法组合起来，例如先分区、再抽学校、再抽班级。各层抽样比例不同时，最后要\textbf{按比例加权}还原全体的情况。

\section{问卷本身带来的偏差}

即使抽样完美，问的方式也会造成偏差：
\begin{itemize}
  \item \textbf{措辞}：“你是否支持增加对困难家庭的帮助？”和“你是否支持增加福利开支？”说的是同一件事，支持率却可能差很多。
  \item \textbf{顺序}：前面的问题会影响后面问题的回答。
  \item \textbf{社会期望}：人们倾向于少报不光彩的行为（作弊、抽烟），多报体面的行为（读书、锻炼、投票）。
  \item \textbf{回忆误差}：“上周你平均每天睡几个小时？”很难答准，改用每天记录的睡眠日志会好得多。
\end{itemize}

\paragraph{随机化回答：让人放心说真话}
想知道“考试中作过弊的学生占多少”，直接问很难得到真话。一个巧妙的设计是：每人先私下抛一枚硬币，
\begin{itemize}
  \item 正面：如实回答“你作过弊吗？”；
  \item 反面：再抛一次，正面答“是”，反面答“否”。
\end{itemize}
调查者看不到硬币，单个人回答“是”什么也说明不了，大家就敢说真话。设真实比例为 $p$，由图~\ref{fig:es-rr}，
\[
  P(\text{答“是”})=\tfrac12p+\tfrac14 .
\]
若 35\% 的人答“是”，则 $\frac12p+\frac14=0.35$，$p=0.2$，即约 20\% 的人作过弊。代价是随机性带来了额外的误差，需要更大的样本。

\begin{figure}[!htbp]
\centering
\begin{tikzpicture}[
    box/.style={draw=mjInk!50, rounded corners=2pt, font=\small, inner sep=4pt, fill=white, align=center},
    yes/.style={box, draw=mjSecond, fill=mjSecond!10, line width=.9pt},
    e/.style={draw=mjInk!50, line width=.6pt}]
  \node[box] (s) at (0,0) {抛硬币};
  \node[box] (t) at (-3,-1.3) {如实回答};
  \node[box] (f) at (3,-1.3) {再抛一次};
  \node[yes] (ty) at (-4.3,-2.9) {答“是”\\{\footnotesize 概率 $\frac12p$}};
  \node[box] (tn) at (-1.7,-2.9) {答“否”};
  \node[yes] (fy) at (1.7,-2.9) {答“是”\\{\footnotesize 概率 $\frac14$}};
  \node[box] (fn) at (4.3,-2.9) {答“否”};
  \draw[e] (s) -- node[above left, slab] {正面 $\frac12$} (t);
  \draw[e] (s) -- node[above right, slab] {反面 $\frac12$} (f);
  \draw[e] (t) -- node[left, slab] {$p$} (ty);
  \draw[e] (t) -- node[right, slab] {$1-p$} (tn);
  \draw[e] (f) -- node[left, slab] {$\frac12$} (fy);
  \draw[e] (f) -- node[right, slab] {$\frac12$} (fn);
\end{tikzpicture}
\caption{随机化回答：答“是”的总概率为 $\frac12p+\frac14$}
\label{fig:es-rr}
\end{figure}

\section{设计一项调查的清单}

\begin{enumerate}
  \item \textbf{明确总体和指标}：调查谁？要估计什么（平均数、比例）？
  \item \textbf{找抽样框}：名单是否完整？漏了谁、多了谁？
  \item \textbf{用概率抽样}：简单随机、分层、整群或组合；不要“方便抽样”（街头随便拦人、网上自愿填写）。
  \item \textbf{提高回答率}：匿名、简短、追访没回答的人，并比较回答者与不回答者的差别。
  \item \textbf{设计问题}：措辞中立，避免诱导；敏感问题考虑随机化回答；能测量就不靠回忆。
  \item \textbf{加权与报告}：按各层比例加权，报告样本量、回答率和误差范围。
\end{enumerate}

\begin{pitfall}
\begin{itemize}
  \item “样本量大就可靠”——错。样本量只管随机误差，管不了偏差。
  \item “网上投票有十万人参加，很有代表性”——错。自愿参加的人往往是对话题特别有兴趣或特别有意见的人（自愿回应偏差）。
  \item 回答“如何减少偏差”时，要针对具体的偏差来源给出具体做法，而不是笼统地说“扩大样本”。
\end{itemize}
\end{pitfall}

\begin{oneline}
代表性比数量重要：样本量治随机误差，概率抽样和好问卷治系统偏差。
\end{oneline}

\begin{exercises}[练一练]
  \item 某 App 在首页弹窗询问“你对本 App 满意吗？”，95\% 的人选“满意”。指出这个调查中至少两种偏差。
  \item 想把误差范围从 $\pm4\%$ 缩小到 $\pm2\%$，样本量要变为原来的多少倍？
  \item 用骰子做随机化回答：掷出 1--4 点如实回答，5 点答“是”，6 点答“否”。若 40\% 的人答“是”，估计真实比例 $p$。
  \item 学校想了解全校 3000 名学生每天的睡眠时间。有同学建议“晚自习后在宿舍楼门口随机拦 100 个人问”。评价这个方案，并提出改进办法。
\end{exercises}

% =============================================================================
\lecture[number=11, short=相关与因果]{相关与因果：一起变，不等于谁导致谁}
% =============================================================================

\begin{keypoints}
\begin{itemize}
  \item 相关系数 $r$ 衡量两个量之间\textbf{直线关系}的强弱和方向。$r\approx0$ 不代表没关系，可能是曲线关系。
  \item 相关不等于因果，常见的原因有：\textbf{共同原因}（混杂）、\textbf{因果倒置}、\textbf{自我选择}、巧合。
  \item \textbf{随机对照实验}用抽签分组，让两组在各方面（包括看不见的方面）平均相同，是判断因果的“金标准”。
  \item 做不了实验时，\textbf{双重差分}、\textbf{断点回归}、\textbf{自然实验}都在设法找“像抽签一样”的比较。
\end{itemize}
\end{keypoints}

\section{相关：两个量一起变}

身高高的人体重往往也大，气温高的日子冰淇淋卖得多，这叫两个量\textbf{相关}。把每个人（每一天）画成一个点，就得到\textbf{散点图}（图~\ref{fig:es-scatter}）。
衡量相关程度的常用指标是\textbf{相关系数} $r$：
\begin{itemize}
  \item $r$ 在 $-1$ 和 $1$ 之间。$r>0$ 叫正相关（一个大另一个也大），$r<0$ 叫负相关；
  \item $|r|$ 越接近 1，点越紧密地排在一条直线附近；
  \item $r$ 只衡量\textbf{直线}关系。图~\ref{fig:es-scatter}(d) 中两个量的关系非常明显（U 形），$r$ 却几乎是 0。
\end{itemize}

\begin{figure}[!htbp]
\centering
\pgfplotsset{sc/.style={mojing box, width=5.2cm, height=4.4cm, xmin=0, xmax=10, ymin=0, ymax=10,
    xtick=\empty, ytick=\empty, every axis plot/.append style={only marks, mark=*, mark size=1.3pt}}}
\begin{subfigure}{.24\textwidth}\centering
\begin{tikzpicture}
\begin{axis}[sc, width=4.3cm, height=4.1cm]
\addplot[mjAccent] coordinates {(8.01,6.85) (5.87,4.63) (3.10,3.52) (0.89,4.42) (9.26,7.45) (5.87,5.33) (7.61,7.28) (8.69,8.08) (6.69,7.44) (2.21,2.58) (9.33,8.98) (3.06,3.45) (6.16,4.31) (5.73,4.27) (5.90,6.30) (5.32,5.63) (9.46,8.57) (5.02,3.85) (7.44,5.84) (4.95,3.55) (9.48,5.49) (9.31,7.24) (4.04,5.65) (3.40,4.83) (8.26,5.91) (7.69,7.42) (6.72,5.19) (4.18,4.53) (4.01,4.92) (1.68,2.38) (6.13,4.55) (1.24,3.12) (2.97,3.35) (6.41,5.02) (0.63,2.74) (8.01,6.29) (1.15,2.81) (5.22,3.93) (5.42,5.65) (2.55,2.37)};
\end{axis}
\end{tikzpicture}
\caption{$r\approx0.85$}
\end{subfigure}\hfill
\begin{subfigure}{.24\textwidth}\centering
\begin{tikzpicture}
\begin{axis}[sc, width=4.3cm, height=4.1cm]
\addplot[mjAccent] coordinates {(1.16,7.60) (3.70,7.11) (3.80,4.98) (1.95,6.99) (2.07,4.61) (9.18,6.63) (5.31,4.29) (0.79,5.77) (1.42,4.71) (1.13,6.39) (1.87,6.26) (8.28,2.59) (0.50,6.00) (8.20,7.39) (6.43,4.16) (7.35,6.52) (4.49,6.81) (9.09,5.96) (3.54,2.58) (4.66,5.69) (6.65,7.66) (8.33,5.44) (3.92,4.50) (3.77,3.57) (3.90,5.81) (8.86,9.33) (7.19,3.69) (4.16,6.08) (2.64,5.25) (0.76,2.29) (4.59,3.78) (5.33,6.64) (6.68,4.84) (4.44,5.79) (2.72,5.04) (3.11,7.26) (9.01,1.36) (4.09,6.23) (2.23,3.70) (2.50,7.29)};
\end{axis}
\end{tikzpicture}
\caption{$r\approx0$}
\end{subfigure}\hfill
\begin{subfigure}{.24\textwidth}\centering
\begin{tikzpicture}
\begin{axis}[sc, width=4.3cm, height=4.1cm]
\addplot[mjSecond] coordinates {(6.64,5.51) (0.98,4.98) (2.48,8.24) (2.16,6.09) (2.08,5.37) (7.81,3.05) (8.81,3.47) (2.99,5.40) (7.88,7.24) (8.51,2.65) (5.12,5.14) (2.70,6.13) (7.92,3.97) (2.42,6.28) (7.17,3.54) (6.17,3.81) (8.85,2.30) (2.59,3.99) (7.69,4.33) (5.16,6.49) (2.58,5.35) (1.99,5.58) (4.98,4.69) (5.74,4.17) (2.16,4.24) (0.63,8.91) (4.74,5.88) (7.05,3.26) (8.77,4.01) (6.13,4.04) (8.75,3.28) (8.28,4.30) (2.46,7.56) (8.30,5.96) (7.08,6.73) (3.00,5.60) (7.67,3.31) (8.29,3.92) (3.19,6.07) (5.24,3.98)};
\end{axis}
\end{tikzpicture}
\caption{$r\approx-0.6$}
\end{subfigure}\hfill
\begin{subfigure}{.24\textwidth}\centering
\begin{tikzpicture}
\begin{axis}[sc, width=4.3cm, height=4.1cm]
\addplot[mjPine] coordinates {(6.13,1.81) (8.57,5.40) (7.48,3.39) (2.53,2.60) (3.20,2.52) (8.36,5.71) (0.55,6.52) (7.89,4.52) (7.67,3.72) (4.71,1.14) (3.23,3.64) (3.01,3.15) (2.79,2.24) (4.51,1.62) (5.04,1.85) (5.48,1.46) (9.46,7.88) (7.63,3.54) (6.10,2.26) (9.40,8.17) (2.44,3.06) (1.94,4.43) (6.01,1.53) (0.90,6.63) (0.82,6.03) (5.13,1.16) (4.70,1.41) (8.75,6.27) (6.16,2.59) (5.13,0.71) (4.97,1.02) (2.73,3.44) (0.61,6.10) (2.23,3.52) (6.73,2.34) (2.31,4.43) (3.83,2.33) (0.53,7.29) (7.97,3.93) (1.89,4.25)};
\end{axis}
\end{tikzpicture}
\caption{U 形，$r\approx0$}
\end{subfigure}
\caption{四组数据的散点图与相关系数}
\label{fig:es-scatter}
\end{figure}

\begin{remark}
相关系数的公式是
$r=\dfrac{\sum(x_k-\bar x)(y_k-\bar y)}{\sqrt{\sum(x_k-\bar x)^2}\sqrt{\sum(y_k-\bar y)^2}}$。
分子里，点落在“$x$ 和 $y$ 都高于平均”或“都低于平均”的区域时贡献正数，否则贡献负数。由柯西不等式可知 $|r|\le1$。
\end{remark}

\section{相关为什么不等于因果}

“冰淇淋卖得多的月份，溺水事故也多，所以应该限制卖冰淇淋。”这显然荒唐。可是换一个包装，同样的错误就不那么容易看出来了。常见的原因有四种（图~\ref{fig:es-dag}）。

\begin{figure}[!htbp]
\centering
\tikzset{v/.style={node box, minimum width=1.9cm, minimum height=.75cm},
         q/.style={-{Stealth[length=2.4mm,width=1.8mm]}, draw=mjSecond, line width=.8pt, dash pattern=on 3pt off 2pt}}
\begin{subfigure}{.32\textwidth}\centering
\begin{tikzpicture}
  \node[v, fill=mjOchre!15] (t) at (0,1.5) {气温};
  \node[v] (i) at (-1.25,0) {冰淇淋销量};
  \node[v] (d) at (1.25,0) {溺水事故};
  \draw[cause] (t) -- (i); \draw[cause] (t) -- (d);
  \draw[q] (i) -- node[below, slab, mjSecond] {?} (d);
\end{tikzpicture}
\caption{共同原因（混杂）}
\end{subfigure}\hfill
\begin{subfigure}{.32\textwidth}\centering
\begin{tikzpicture}
  \node[v] (p) at (-1.25,.75) {警察人数};
  \node[v] (c) at (1.25,.75) {犯罪数量};
  \draw[cause] (c.north) to[bend right=35] node[above, slab] {真实方向} (p.north);
  \draw[q] (p.south) to[bend right=35] node[below, slab, mjSecond] {误读的方向} (c.south);
\end{tikzpicture}
\caption{因果倒置}
\end{subfigure}\hfill
\begin{subfigure}{.32\textwidth}\centering
\begin{tikzpicture}
  \node[v, fill=mjOchre!15] (m) at (0,1.5) {学习动力};
  \node[v] (b) at (-1.25,0) {参加补习};
  \node[v] (s) at (1.25,0) {成绩};
  \draw[cause] (m) -- node[left, slab] {自己选择} (b); \draw[cause] (m) -- (s);
  \draw[q] (b) -- node[below, slab, mjSecond] {?} (s);
\end{tikzpicture}
\caption{自我选择}
\end{subfigure}
\caption{三种常见的“假因果”。黑箭头是真实的影响，红色虚线是容易被误读的因果。}
\label{fig:es-dag}
\end{figure}

\begin{enumerate}
  \item \textbf{共同原因（混杂因素）}：气温高时，买冰淇淋的人多，去游泳的人也多。冰淇淋和溺水之间没有因果关系，是气温同时影响了两者。
  \item \textbf{因果倒置}：犯罪多的街区警察也多，不是警察导致了犯罪，而是犯罪多才派了更多警察。
  \item \textbf{自我选择}：参加补习的学生成绩更好，可能只是因为学习动力强的学生更愿意去补习，而动力本身就会提高成绩。
  \item \textbf{巧合}：比较的变量足够多时，总能找到一些毫无关系却“高度相关”的数据，例如某国人均奶酪消费量和某种意外死亡人数。
\end{enumerate}

还有一种和自我选择相近的\textbf{幸存者偏差}：二战时，有人统计返航飞机身上的弹孔，发现机翼上多、发动机上少，于是建议加固机翼。
统计学家沃尔德指出，应该加固的恰恰是发动机：发动机中弹的飞机大多没能飞回来，所以在返航的飞机上看不到。我们只看到了“幸存”下来的样本。

\section{随机对照实验}

要知道“参加补习能不能提高成绩”，最可靠的办法是\textbf{随机对照实验}（图~\ref{fig:es-rct}）：
\begin{enumerate}
  \item 招募愿意参加的学生，比如 200 人；
  \item \textbf{抽签}决定谁参加（实验组 100 人）、谁暂时不参加（对照组 100 人）；
  \item 一段时间后，比较两组成绩的变化。
\end{enumerate}

\begin{figure}[!htbp]
\centering
\begin{tikzpicture}[st/.style={node box, minimum height=.85cm, minimum width=2.5cm},
    fl/.style={-{Stealth[length=2.4mm,width=1.8mm]}, draw=mjInk!70, line width=.8pt}]
  \node[st] (a) at (0,0) {招募 200 名\\ 自愿者};
  \node[st, draw=mjOchre, fill=mjOchre!12, minimum width=1.5cm] (b) at (3.1,0) {抽签};
  \node[st, draw=mjSecond, fill=mjSecond!8] (c) at (6.3,.85) {实验组 100 人\\ 参加补习};
  \node[st, draw=mjAccent, fill=mjAccent!8] (d) at (6.3,-.85) {对照组 100 人\\ 不参加};
  \node[st, minimum width=2.2cm] (e) at (9.6,0) {比较两组\\ 成绩变化};
  \draw[fl] (a) -- (b); \draw[fl] (b) -- (c.west); \draw[fl] (b) -- (d.west);
  \draw[fl] (c.east) -- (e); \draw[fl] (d.east) -- (e);
  \node[slab, align=center, mjOchre] at (3.1,-1.55) {两组在动力、基础等\\ 各方面平均相同};
\end{tikzpicture}
\caption{随机对照实验：抽签让两组除了“是否参加”之外，在其他方面都差不多}
\label{fig:es-rct}
\end{figure}

为什么要抽签？因为抽签和学生的动力、基础、家庭背景都没有关系。只要人数够多，两组在这些方面\textbf{平均相同}，包括那些我们想不到、测不出来的因素。
于是两组之间的差别，只能归因于“是否参加补习”。

医学试验还会加上两样东西：
\begin{itemize}
  \item \textbf{安慰剂}：对照组也吃一片外观相同但没有药效的药片，因为“以为自己在治疗”本身就会让人感觉变好；
  \item \textbf{双盲}：病人和医生都不知道谁在哪一组，避免有意无意的影响。
\end{itemize}

\begin{pitfall}
\textbf{随机抽样}和\textbf{随机分组}是两回事：
\begin{itemize}
  \item 随机\textbf{抽样}（第 10 讲）是为了让样本\textbf{代表总体}，结论能推广到全体；
  \item 随机\textbf{分组}是为了让两组\textbf{可以比较}，结论能说明因果。
\end{itemize}
一项研究可以只做到其中一样：从全国随机抽了 1000 名学生，却让他们自己选择是否参加训练营，结论能代表全国，却不能说明训练营的效果。
\end{pitfall}

\section{做不了实验怎么办}

很多问题不能做实验：不能抽签让一些人抽烟，也不能抽签决定哪个城市遭受经济危机。这时研究者会寻找“\textbf{像抽签一样}”的情形。

\paragraph{双重差分}
某市部分学校从今年起推行一项新的教学方法，另一部分学校没有。只比较今年两类学校的成绩不行，它们原本就可能不一样；
只比较试点学校今年和去年的成绩也不行，今年的题目可能更容易。\textbf{双重差分}把两种比较合在一起（图~\ref{fig:es-did}）：
\[
  \text{效果}=(\text{试点校的变化})-(\text{对照校的变化})=(80-70)-(70-65)=5 .
\]
对照校的变化（$+5$）代表“就算不推行新方法也会发生的变化”。这个方法依赖一个假设：如果没有新方法，两类学校的成绩会\textbf{平行地}变化（\textbf{平行趋势}）。

\begin{figure}[!htbp]
\centering
\begin{tikzpicture}
\begin{axis}[mojing box, width=8.6cm, height=6cm, xmin=-.4, xmax=1.4, ymin=60, ymax=85,
    xtick={0,1}, xticklabels={推行前,推行后}, ylabel={平均成绩}, ytick={65,70,75,80}, legend pos=north west]
  \addplot[mjSecond, mark=*, line width=1.1pt] coordinates {(0,70) (1,80)};
  \addplot[mjAccent, mark=*, line width=1.1pt] coordinates {(0,65) (1,70)};
  \addplot[mjSecond, dashed, line width=.8pt] coordinates {(0,70) (1,75)};
  \legend{试点学校,对照学校,试点校若不推行（推测）}
  \draw[decorate, decoration={brace, amplitude=4pt}, mjPine, line width=.9pt] (axis cs:1.04,80) -- (axis cs:1.04,75)
    node[midway, right=5pt, slab, mjPine] {效果 5};
\end{axis}
\end{tikzpicture}
\caption{双重差分：用对照学校的变化推测试点学校“如果不推行”会怎样}
\label{fig:es-did}
\end{figure}

\paragraph{断点回归}
某奖学金规定考试 500 分以上才能获得。考 501 分和考 499 分的学生，能力几乎没有差别，差别只在于是否拿到了奖学金。
比较刚好过线和刚好没过线的学生后来的情况，就像做了一次抽签实验（图~\ref{fig:es-rd}）：如果在 500 分处出现了一个“跳跃”，跳跃的高度就是奖学金的效果。

\begin{figure}[!htbp]
\centering
\begin{tikzpicture}
\begin{axis}[mojing box, width=9.5cm, height=6cm, xmin=468, xmax=532, ymin=22, ymax=62,
    xlabel={考试分数}, ylabel={后来的某项结果}, xtick={470,480,490,500,510,520,530}, ytick=\empty]
  \addplot[only marks, mark=*, mark size=1.3pt, mjAccent!70] coordinates {(475.14,32.38) (484.21,30.87) (475.65,33.14) (495.99,34.12) (498.74,37.55) (479.58,34.42) (476.82,27.14) (493.47,36.62) (495.84,37.20) (487.05,33.16) (487.56,34.83) (470.09,25.75) (487.90,37.42) (488.84,37.75) (498.28,37.99) (471.82,28.38) (492.45,34.15) (475.45,30.10) (482.43,35.12) (487.89,31.96) (483.12,35.76) (495.71,40.30) (476.14,31.41) (493.64,38.27) (498.78,45.18) (478.78,35.87) (487.52,31.14) (486.52,33.94) (493.98,32.88) (481.80,35.98)};
  \addplot[only marks, mark=*, mark size=1.3pt, mjSecond!80] coordinates {(518.08,52.46) (504.93,45.60) (514.07,53.47) (501.00,47.07) (505.21,51.32) (514.27,48.37) (527.38,56.36) (508.91,52.21) (511.77,50.07) (528.41,59.64) (523.50,56.73) (505.11,48.56) (516.40,49.98) (512.42,53.36) (509.63,48.66) (525.89,55.55) (507.81,45.98) (514.51,51.07) (513.33,53.34) (519.79,53.40) (509.46,48.61) (510.97,50.48) (519.20,52.19) (515.52,49.69) (522.71,54.24) (520.99,54.54) (511.91,53.46) (522.27,53.05) (503.71,48.47) (506.77,50.95)};
  \addplot[mjAccent, line width=1.1pt, domain=470:500, samples=2] {40+0.35*(x-500)};
  \addplot[mjSecond, line width=1.1pt, domain=500:530, samples=2] {46+0.35*(x-500)};
  \draw[mjInk!60, dashed] (axis cs:500,22) -- (axis cs:500,62);
  \node[slab, anchor=south west] at (axis cs:500.5,23) {500 分线};
  \draw[decorate, decoration={brace, amplitude=4pt, mirror}, mjPine, line width=.9pt] (axis cs:501.5,40) -- (axis cs:501.5,46)
    node[midway, right=5pt, slab, mjPine] {跳跃 = 效果};
  \node[slab, mjAccent, anchor=west] at (axis cs:470,56) {蓝：没拿到奖学金};
  \node[slab, mjSecond, anchor=west] at (axis cs:470,52) {红：拿到奖学金};
\end{axis}
\end{tikzpicture}
\caption{断点回归：分数线两侧的学生几乎一样，结果在分数线处的跳跃就是奖学金的效果}
\label{fig:es-rd}
\end{figure}

\paragraph{自然实验}
有时现实本身就提供了“抽签”：热门学校用摇号决定录取，摇中和没摇中的学生可以比较；某项政策先在一部分地区试点，而试点地区的选择与要研究的结果无关……
这些都叫\textbf{自然实验}。

\paragraph{控制变量}
最常见的做法是在比较时“控制”可以观测到的因素，例如比较前测成绩、家庭背景相同的学生。它比什么都不做好，但\textbf{控制不了没有观测到的因素}（比如学习动力），
所以结论要说明这一局限。

\section{判断一项因果说法的清单}

\begin{enumerate}
  \item 有没有\textbf{对照}？只看一组人的变化，说明不了什么。
  \item 两组是怎么分的？是\textbf{随机分组}，还是自己选择的？
  \item 有没有可能的\textbf{共同原因}？
  \item 因果的\textbf{方向}会不会反过来？
  \item 样本能不能\textbf{代表}要推广的人群？
  \item 效果有\textbf{多大}？在统计上可靠吗？有没有被其他研究重复验证？
\end{enumerate}

\begin{oneline}
相关不等于因果：先找共同原因、反向因果和自我选择；要说因果，最好抽签分组，做不到就找“像抽签一样”的比较。
\end{oneline}

\begin{exercises}[练一练]
  \item “研究发现，每天吃早餐的学生成绩更好，所以学校应该强制所有学生吃早餐。”指出这个推理可能存在的问题，并设计一个能检验因果关系的方案。
  \item 某市从 2025 年起在部分学校推行“每天户外活动 2 小时”。试点学校学生的近视率从 30\% 变为 32\%，非试点学校从 28\% 变为 33\%。用双重差分估计这项措施的效果，并说明这个估计依赖什么假设。
  \item 一项研究从全国随机抽取了 1000 名高中生，让他们自己决定是否参加一个暑期训练营，结果参加者的成绩提高得更多。这项研究的结论能推广到全国吗？能说明训练营的效果吗？
\end{exercises}

% =============================================================================
\lecture[number=12, short=答题方法与答案]{答题方法与练一练答案}
% =============================================================================

\begin{keypoints}
\begin{itemize}
  \item 计算题：先画图，再列式；写清设定和单位；用“楔子 $x+y=t$”“剩余加起来不变”等关系检查结果。
  \item 分析题四步：\textbf{结论 $\to$ 机制 $\to$ 条件 $\to$ 另一面}。用模型说明“为什么”，再说明模型在什么条件下成立、现实中还要考虑什么。
  \item 社科题先找“关键词”，想到对应的工具：比率想到分母和权重，调查想到偏差，“有关”想到混杂和随机分组。
\end{itemize}
\end{keypoints}

\section{计算题}

\begin{enumerate}
  \item \textbf{先画图}。哪怕是草图，也能帮你看清哪条线动、面积是哪一块。经济题里很多“面积”就是三角形和矩形。
  \item \textbf{写清设定}。例如“设买方价格为 $P_{\text{买}}$，卖方到手价格为 $P_{\text{买}}-t$”；单位（元/件、月利率还是年利率）写出来。
  \item \textbf{检查}。几个随手可用的检查：买方多付的加卖方少得的等于税额；各块面积加起来等于原来的总剩余；贴现后的还款总现值等于借款额；
        两两表决的票数加起来等于总人数。
\end{enumerate}

\section{分析题怎么写}

分析题、论述题最怕“只有结论没有道理”或“只有一堆名词”。一个稳妥的结构是四步（图~\ref{fig:es-steps}）。

\begin{figure}[!htbp]
\centering
\begin{tikzpicture}[st/.style={node box, minimum height=1.35cm, minimum width=2.9cm},
    fl/.style={-{Stealth[length=2.4mm,width=1.8mm]}, draw=mjInk!70, line width=.8pt}]
  \node[st, draw=mjAccent, fill=mjAccent!8] (a) at (0,0) {\textbf{1. 结论}\\ 一句话回答问题};
  \node[st, draw=mjPine, fill=mjPine!8] (b) at (3.55,0) {\textbf{2. 机制}\\ 用哪个模型、\\哪条线怎样动};
  \node[st, draw=mjOchre, fill=mjOchre!10] (c) at (7.1,0) {\textbf{3. 条件}\\ 模型的假设、\\成立的前提};
  \node[st, draw=mjSecond, fill=mjSecond!8] (d) at (10.65,0) {\textbf{4. 另一面}\\ 现实的复杂性、\\局限与建议};
  \draw[fl] (a) -- (b); \draw[fl] (b) -- (c); \draw[fl] (c) -- (d);
\end{tikzpicture}
\caption{分析题的四步结构}
\label{fig:es-steps}
\end{figure}

\begin{example}[一个示范]
\textbf{问}：为什么节假日的热门车票总是“一票难求”？

\textbf{答}：（结论）因为票价被固定在低于高峰期均衡价格的水平，形成了短缺。
（机制）节假日出行需求大幅增加，需求曲线右移，而运力短期内难以增加，供给曲线接近竖直；票价不随之上涨，想买的人数远多于座位数（第 4 讲最高限价）。
价格不再承担分配的功能，只好靠排队、抢票来分配，抢票的时间和精力就是隐性的代价。
（条件）这一分析假设票价不能灵活调整、运力短期固定。
（另一面）完全按供求涨价能减少短缺，但会让收入较低、又必须回家的人承受很高的票价，所以现实中常保持票价稳定，同时增开临时列车、实行实名制和分时段放票，在效率和公平之间做权衡。
\end{example}

\section{看到什么，想到什么}

\begin{center}
\small
\begin{tabular}{p{4.2cm}p{8.8cm}}
  \toprule
  题目里出现 & 想到 \\
  \midrule
  价格、征税、补贴、限价 & 供求图；哪条线动；弹性；剩余与无谓损失（第 1--4 讲） \\
  污染、拥堵、疫苗 & 外部性；社会成本；庇古税、排污交易（第 4 讲） \\
  两家企业、合作与背叛 & 收益表、纳什均衡、囚徒困境；重复博弈与贴现因子（第 5 讲） \\
  利息、分期、贷款、“哪个划算” & 现值；年金公式；同一时刻比较（第 6 讲） \\
  分工、贸易、“谁更强” & 机会成本、比较优势；交换比例范围；分配效应（第 7 讲） \\
  投票、表决、规则 & 孔多塞循环、议程、不同规则不同赢家、策略性投票、阿罗定理、中位选民（第 8 讲） \\
  比率、百分比、“合计” & 分母；加权平均；辛普森悖论；基础比率（第 9 讲） \\
  调查、民调、问卷 & 抽样框、无回答、措辞；随机误差与偏差；概率抽样（第 10 讲） \\
  “A 和 B 有关”“因此应该……” & 混杂、反向因果、自我选择；随机对照实验；双重差分（第 11 讲） \\
  \bottomrule
\end{tabular}
\end{center}

\begin{pitfall}
\begin{itemize}
  \item 名词堆砌不得分：写了“阿罗不可能定理”，却没有说明它和本题的关系，等于没写。
  \item 开放题要有\textbf{具体}的做法：“扩大样本”“加强监管”这类笼统的话，换成“按年级分层、随机抽取班级、匿名填写一周的睡眠日志”这样具体的步骤。
  \item 题目给了定义（如“纳什均衡指……”），就严格按题目的定义来做，不要套用自己记得的其他说法。
\end{itemize}
\end{pitfall}

\section{练一练答案}

\paragraph{第 1 讲}
\begin{enumerate}[label=1.\arabic*, leftmargin=2.6em]
  \item $120-4P=2P$，均衡 $P=20$，$Q=40$。价格定为 15 元时，需求量 60，供给量 30，\textbf{短缺} 30。
  \item (a) 电池是电动车的主要成本，成本下降使\textbf{供给}右移：价格下降，数量增加。
        (b) 汽油涨价使燃油车用起来更贵，电动车作为替代品的\textbf{需求}右移：价格上升，数量增加。
        (c) 两者都使数量增加，所以数量一定增加；一个使价格下降、一个使价格上升，价格变化不确定。
  \item 观察到的“价格上涨、成交量增加”，是不同时期的供求曲线的不同交点，不是同一条需求曲线上的点。很可能是收入增长、人口流入、对房价继续上涨的预期等使需求曲线右移了，
        新交点的价格和数量都更高。需求曲线向下倾斜说的是“其他条件不变时”价格与需求量的关系，这个现象并不与它矛盾。
\end{enumerate}

\paragraph{第 2 讲}
\begin{enumerate}[label=2.\arabic*, leftmargin=2.6em]
  \item 需求曲线与纵轴交于 $P=30$。均衡 $P=20$，$Q=40$。消费者剩余 $\frac12\times40\times(30-20)=200$，生产者剩余 $\frac12\times40\times20=400$。
  \item 价格上涨 $20\%$，游客减少 $10\%$，$|E|\approx\frac{10\%}{20\%}=0.5$，缺乏弹性。门票收入从 $200000$ 元增加到 $216000$ 元。
  \item 临时订票的多是商务出行或有急事的人，没有别的选择，需求缺乏弹性，航空公司定高价也不会少卖多少；提前订票的多是可去可不去的旅游者，需求富有弹性，降价能吸引更多人。
        对弹性不同的人收不同的价格（\textbf{价格歧视}），能增加航空公司的收入。
\end{enumerate}

\paragraph{第 3 讲}
\begin{enumerate}[label=3.\arabic*, leftmargin=2.6em]
  \item $b=4$，$d=2$，$t=3$。买方价格上升 $x=\frac{2}{6}\times3=1$，卖方价格下降 $y=\frac46\times3=2$。原均衡 $P=20$，所以买方共付 21 元（其中标价 18 元，另交税 3 元），卖方到手 18 元，
        成交量 $2\times18=36$（检验：$120-4\times21=36$）。税收 $3\times36=108$ 元。买方每件负担 1 元，卖方负担 2 元。无谓损失 $\frac12\times3\times(40-36)=6$。
  \item 记 $k=\frac{bd}{b+d}$，则无谓损失 $=\frac12kt^2$，$t$ 加倍后变为 $\frac12k(2t)^2$，是原来的 4 倍。税收 $R(t)=t(Q^*-kt)$，
        $R(2t)=2t(Q^*-2kt)=2R(t)-2kt^2<2R(t)$。
  \item 富人购买豪华跑车，有很多替代选择（别的奢侈品、去国外买、推迟购买），需求的弹性可能很大，那么税负会有相当一部分落在生产商、经销商和相关工人身上，未必主要由富人承担。
        要判断，还需要知道需求和供给两边的弹性：例如短期内工厂产能固定、供给缺乏弹性，卖方负担会更多。
\end{enumerate}

\paragraph{第 4 讲}
\begin{enumerate}[label=4.\arabic*, leftmargin=2.6em]
  \item (1) 价格 25 元时，需求量 20，供给量 50，过剩 30。(2) $30\times25=750$ 元。
  \item (1) 工厂只考虑自己的成本，不考虑渔民的损失，每吨的社会成本比工厂的成本高 50 元，所以生产到了对社会来说不值得的数量。
        (2) 每吨征税 50 元后，工厂自己的成本等于社会成本，它会自己把产量减到对社会最好的水平。
        (3) 工厂每生产 1 吨要赔 50 元，效果和征税 50 元一样，产量同样减到有效率的水平；区别是这 50 元给了受损的渔民而不是政府。反过来，即使法律允许工厂排污，只要协商成本低，渔民也会出钱请工厂减产（每少产 1 吨，渔民最多愿出 50 元），产量同样会达到有效率的水平；
        产权给谁只影响钱归谁。这就是科斯的思想：产权明确、协商成本低时，私人之间就能解决外部性。
  \item 偏少。个人决定是否接种时只考虑自己得到的保护，不考虑对周围人的保护（正外部性），所以接种的人数会少于对社会最好的人数。这是政府补贴或免费提供疫苗的理由。
\end{enumerate}

\paragraph{第 5 讲}
\begin{enumerate}[label=5.\arabic*, leftmargin=2.6em]
  \item 收益表中，（猎鹿，猎鹿）和（抓兔，抓兔）都是纳什均衡；另两格中，猎鹿的一方得 0，会改去抓兔。（猎鹿，猎鹿）对双方都更好；
        但抓兔更“保险”：不论对方怎么做，抓兔都能得 3，猎鹿则可能因为对方不合作而一无所获。所以对对方没有把握时，人们可能停在较差的均衡上。
  \item (1) $\delta\ge\frac{6-4}{6-2}=\frac12$。(2) 继续的概率 0.8 起到贴现因子的作用，$\delta=0.8\ge\frac12$，合作可以维持。
  \item 在位者打价格战得 6，容忍得 5，会打价格战；进入者预见到这一点，进入得 $-2$，不进入得 0，选择不进入。结果是（不进入）。
        扩建产能改变了在位者事后的利益，使“打价格战”的威胁变得可信，从而把潜在的竞争者挡在门外。
\end{enumerate}

\paragraph{第 6 讲}
\begin{enumerate}[label=6.\arabic*, leftmargin=2.6em]
  \item 4 年后的 10000 元，现值为 $\frac{10000}{1.04^4}\approx8548$ 元，比现在领的 8500 元略多，4 年后领更值钱（前提是利率确实是 4\%，且奖学金一定能兑现）。
  \item 记 $v=\frac1{1+i}$，现值 $=Av+Av^2+\cdots+Av^n=Av\cdot\frac{1-v^n}{1-v}$。因为 $\frac{v}{1-v}=\frac{1/(1+i)}{i/(1+i)}=\frac1i$，所以现值 $=A\cdot\frac{1-(1+i)^{-n}}{i}$。
        $n\to\infty$ 时 $(1+i)^{-n}\to0$，现值趋于 $\frac Ai$。
  \item 每月还本金 1000 元，第 $k$ 个月利息为 $(12000-1000(k-1))\times1\%$，总利息 $=120\times\frac{12+1}{2}=780$ 元。第 1 个月还款 $1000+120=1120$ 元。
  \item $72\div3=24$ 年。精确值 $\frac{\ln2}{\ln1.03}\approx23.4$ 年。
\end{enumerate}

\paragraph{第 7 讲}
\begin{enumerate}[label=7.\arabic*, leftmargin=2.6em]
  \item 小王写 1 页作文要放弃 2 道数学题，小李只要放弃 1 道。小李写作文的机会成本低，小王做数学题的机会成本低（$\frac12$ 页作文，小李是 1 页）。
        所以小王做数学题、小李写作文，虽然小王两样都比小李快。
  \item 不会。1 单位粮换 2.5 单位布，比甲国自己种粮的机会成本（2 单位布）还高，甲国宁可自己种粮，不会买粮；只有乙国愿意卖。交换比例必须在 1 和 2 之间。
  \item 工资低反映的是劳动生产率低，它不可能在所有产品上都有比较优势：两国的机会成本不同，就存在双方都获益的分工。它在劳动密集型产品上有比较优势，
        我们就在其他产品上有比较优势，通过贸易双方都能消费得更多。但国内与进口竞争的行业会受损，需要转岗培训、社会保障等政策来帮助受损者。
\end{enumerate}

\paragraph{第 8 讲}
\begin{enumerate}[label=8.\arabic*, leftmargin=2.6em]
  \item 相对多数：A 4 票，B 3 票，C 2 票，\textbf{A 胜}。两轮决选：A 与 B 进入第二轮，C 的 2 个支持者更喜欢 B，B 以 5 比 4 胜，\textbf{B 胜}。
        波达计数：A $4\times2=8$ 分，B $4+3\times2+2=12$ 分，C $3+2\times2=7$ 分，\textbf{B 胜}。两两表决：B 对 A 为 5 比 4，B 对 C 为 7 比 2，\textbf{B 是孔多塞赢家}。
  \item 例：3 人排序 A，B，C；2 人排序 B，C，A。两两表决 A 对 B、A 对 C 都是 3 比 2，A 是孔多塞赢家；但波达计数 A 得 6 分，B 得 $3+4=7$ 分，B 胜。
  \item 中位数是 5 公里，选 5 公里。加入两位后，理想位置为 2、3、5、9、10、10、10，中位数变为 9 公里，选 9 公里。
\end{enumerate}

\paragraph{第 9 讲}
\begin{enumerate}[label=9.\arabic*, leftmargin=2.6em]
  \item 小结石：A $\frac{81}{87}\approx93\%$，B $\frac{234}{270}\approx87\%$；大结石：A $\frac{192}{263}\approx73\%$，B $\frac{55}{80}\approx69\%$；
        全部：A $\frac{273}{350}=78\%$，B $\frac{289}{350}\approx83\%$。A 在每组都更好，合计却更差：A 法多用于较难治的大结石（263 人），B 法多用于小结石（270 人）。
        结石大小同时影响“用哪种方法”和“治愈率”，应当分组比较。
  \item (1) 设 100000 人：100 人有病，99 人阳性；99900 人没病，999 人阳性。阳性者有病的概率 $\frac{99}{1098}\approx9.0\%$。
        (2) 第二次仍阳性：有病的 99 人中约 98 人，没病的 999 人中约 10 人，概率 $\approx\frac{98.01}{98.01+9.99}\approx90.7\%$。
  \item 可以追问：中位数是多少？包括管理层吗？是税前还是税后，含不含奖金？我应聘的这个岗位的薪酬范围是多少？统计的是哪一年、多少人？
\end{enumerate}

\paragraph{第 10 讲}
\begin{enumerate}[label=10.\arabic*, leftmargin=2.6em]
  \item 只有打开 App 的人才会看到弹窗，不满意而不再使用的人已经离开了（覆盖偏差／幸存者偏差）；很多人直接关闭弹窗，回答的是自愿回应的一部分人（无回答偏差）；
        在 App 里被 App 自己询问，人们可能顺手点“满意”，选项的设计也可能有诱导。
  \item 误差与 $\frac1{\sqrt n}$ 成正比，误差减半，样本量要变为 4 倍。
  \item $P(\text{答“是”})=\frac46p+\frac16=0.4$，$p=(0.4-\frac16)\times\frac32=0.35$。
  \item 问题：只调查了住校生；晚自习后的这个时间、这个地点出现的人不能代表全体；“随机拦人”实际上由访问者挑选，不是概率抽样；一次口头回忆误差大。
        改进：用全校名册按年级分层，随机抽取班级（或学生）；匿名填写，连续记录一周的睡眠日志；追访没有交回的人；按各年级人数加权估计全校平均。
\end{enumerate}

\paragraph{第 11 讲}
\begin{enumerate}[label=11.\arabic*, leftmargin=2.6em]
  \item 可能的问题：家庭条件好、作息规律、自律的学生更可能每天吃早餐，成绩也更好（共同原因、自我选择）；即使吃早餐有用，强制执行的效果也可能不同。
        方案：在学校随机选择一部分班级提供免费早餐并督促食用，其他班级作对照，比较一段时间前后两组成绩的变化。
  \item 效果 $=(32\%-30\%)-(33\%-28\%)=-3$ 个百分点，即这项措施使近视率的上升少了约 3 个百分点。假设：如果不推行这项措施，两类学校近视率的变化会相同（平行趋势），
        并且同一时期没有其他只影响试点学校的变化。
  \item 样本是随机抽取的，所以“参加者的成绩提高更多”这一现象可以推广到全国高中生；但是否参加是学生自己选择的，参加者可能本来就更有动力，
        所以不能说明成绩提高是训练营造成的。要判断效果，应当在报名者中随机决定谁参加。
\end{enumerate}


\end{document}
